\documentclass[11pt,reqno]{amsart}
\usepackage{amsmath,amssymb,amsthm}
\usepackage[margin=1.1in]{geometry}
\usepackage{graphicx}
\usepackage{hyperref}
\usepackage{xr}
\externaldocument[M-]{erdos-634}
\hypersetup{colorlinks=true,linkcolor=blue,citecolor=blue,urlcolor=blue}

\theoremstyle{plain}
\usepackage{booktabs}
\newtheorem{theorem}{Theorem}
\newtheorem{lemma}[theorem]{Lemma}
\newtheorem{corollary}[theorem]{Corollary}
\newtheorem{proposition}[theorem]{Proposition}
\theoremstyle{definition}
\newtheorem{remark}[theorem]{Remark}
\newtheorem{conjecture}[theorem]{Conjecture}
\newtheorem{hypothesis}[theorem]{Hypothesis}
\newtheorem{definition}[theorem]{Definition}

\newcommand{\R}{\mathbb{R}}
\newcommand{\Z}{\mathbb{Z}}
\newcommand{\Q}{\mathbb{Q}}
\newcommand{\N}{\mathbb{N}}
\DeclareMathOperator{\Area}{Area}

\title[A signed-direction invariant for triangle tilings]{A signed-direction invariant for triangle tilings, and the exclusion of primes congruent to $3$ modulo $4$}
\author{Vico Bonfioli}
\date{\today}

\begin{document}

\title{The base-$\beta$ branch of $3\alpha+2\beta=\pi$: boundary walks, and the structure of a
hypothetical prime tiling}

\author{Vico Bonfioli}
\email{vicobonfioli@gmail.com}

\date{\today}

\begin{abstract}
This is a companion note to \emph{A signed-direction invariant for triangle tilings, and the
exclusion of primes $\equiv3\pmod4$}. That paper excludes every prime $N\equiv3\pmod4$ exceeding
$3$ except the base-$\beta$ candidates $N=3f^2-e^2$, which are exactly the primes
$\equiv11\pmod{12}$. Here we record what is known about the internal structure of a hypothetical
tiling in that one open branch: the boundary walk equations and their solutions, the apex-mismatch
configuration, and the alignment and far-region theorems. None of this is needed for the main
theorem; it is collected here so that the main paper can state the exception cleanly and so that
the structural facts are available to anyone attacking the branch.

Two results are worth isolating. First, at $m=1$ the walk equations are completely solved in a
regime: if $f^2>2e^2$ then each equal side of the target carries no $b$-edge, and if
$f^2>2ef+e^2$ then the base carries exactly $e$ of them, so $R_{\mathrm{base}}\in\{e,2e\}$. Second,
for $e=1$ this pins the base walk to a single family, reducing that entire subfamily --- the
candidates $N=3f^2-1$, among them $11,47,107,191,431,587,971$ below $1000$ --- to the question of
whether one explicit walk is realizable.

We also record, as prominently as the positive results, which attacks on this branch are now known
to be dead: the whole class of linear direction invariants, the metric bounds derived from them,
and the hope of closing $e=1$ by proving that every side carries at least two $c$-edges.
\end{abstract}

\maketitle

\section{Scope and conventions}\label{sec:compscope}

Throughout, the tile is the primitive $3\alpha+2\beta=\pi$ triangle
\[
(a,b,c)=(ef,\ f^2-e^2,\ f^2),\qquad \gcd(e,f)=1,\quad 1\le e<f,
\]
with $a$ opposite $\alpha$, $b$ opposite $\beta$, $c$ opposite $\gamma$, and $2\gamma=\pi+\alpha$.
The base-$\beta$ target is the isosceles triangle with base angles $\beta$,
\[
\text{equal sides } S=f^3m,\qquad \text{base } L=em(3f^2-e^2),\qquad N=m^2(3f^2-e^2),
\]
and $N$ is prime exactly when $m=1$. We write $D=4f^2-e^2$, so $\sin\alpha=e\sqrt D/(2f^2)$ and the
tiling's coordinates lie in $\Q(\sqrt D)$. Numbering of the form ``Proposition~4.3'' refers to the
main paper~\cite{Main}; results are stated here only when they are not there.

We use freely from~\cite{Main}: the dissection basics and vertex figures, the corner figures (a
base corner carries a single tile presenting $\beta$; the apex carries exactly three tiles, each
presenting $\alpha$), the $\gamma$-trap (every side of the target carries at least one $c$-edge),
the unsplittability of $b$, and the corner parallelogram.

\section{Boundary walks}\label{sec:compwalks}

A side of the target of length $L$ is exactly partitioned by the tile edges lying on it, so it
determines a \emph{walk equation}
\[
n_a\,ef+n_b\,(f^2-e^2)+n_c\,f^2=L,\qquad n_a,n_b,n_c\ge0 ,
\]
together with $n_c\ge1$ from the $\gamma$-trap. The results of this section solve those equations
and then constrain the order in which the edges occur.

\begin{theorem}[The base-$\beta$ family at $e=1$, $f=2$]\label{thm:basebeta-e1}
The isosceles triangle with base angles $\beta$ and sides $(8,8,11)$ admits no tiling by $N=11$
copies of the tile $(2,3,4)$. Consequently the base-$\beta$ branch is closed for $N=11$,
unconditionally.
\end{theorem}

\begin{proof}
Write $\omega=\alpha/2$. The law of cosines gives $\cos\alpha=(2f^2-e^2)/(2f^2)$, hence
$\sin\omega=e/(2f)$; as $1\le e<f$ places this strictly in $(0,\tfrac12)$, \emph{Niven's theorem}
forces $\alpha/\pi$ irrational --- for every $(e,f)$, with no appeal to \cite{Laczkovich1995}
(\texttt{tile\_alpha\_irrational}). Irrationality pins the vertex figures: a vertex
$x\alpha+y\beta+z\gamma=k\pi$ forces $y+z=2k$ and $2x+z=3y$, so a $\pi$-vertex is $\{3\alpha,2\beta\}$
or $\{\alpha,\beta,\gamma\}$, each base corner carries \emph{exactly one} tile (its $\beta$-vertex) and
the apex \emph{exactly three} (\texttt{vertex\_pi}, \texttt{vertex\_beta\_corner},
\texttt{vertex\_apex}); in particular a $\pi$-vertex carries at most one $\gamma$, since
$2\gamma=\pi+\alpha>\pi$. Each side of $ABC$ is partitioned into whole tile edges (a supporting line
meets a convex tile in a face).

Two citation-free constraints now cut the base. (a) \emph{Corner parallelogram}: the unique
$\beta$-corner tile $T_B$ has its $b$-edge a chord $[P,Q]$, matched (by the mod-$f$ edge arithmetic)
by a single tile $T'$ whose $\alpha,\gamma$ are swapped, so $T_B\cup T'$ is a parallelogram with sides
$a,c$; hence $P,Q$ are $(1,1,1)$-vertices whose third tile carries $\beta$, and the second base edge
from each corner lies in $\{a,c\}$. So a $b$-edge can occupy only positions $3,\dots,k-2$ of a base of
$k$ edges. (b) \emph{$\gamma$-trap}: no boundary junction carries $\gamma$ on both sides, and only a
$c$-edge leaves the $\gamma$-state, so every side carries at least one $c$-edge.

At $e=1$, $m=1$ the base $Y=3f^2-1$ has $q_{\mathrm{base}}\equiv1\pmod f$, and the only compositions
are $B_0=\{b,c,c\}$ ($k=3$), $B_1=(a^{f},b,c)$ ($k=f+2$) and $B_2=(a^{2f},b)$ ($k=2f+1$). $B_2$ dies
by (b). $B_0$ dies by (a) for \emph{every} $f$: its single $b$ would need a position in
$3,\dots,1$. At $f=2$, $B_1$ also dies by (a): $k=4$ leaves positions $3,\dots,2$. So no base
survives, and $N=11$ admits no tiling. \qedhere
\end{proof}

\begin{theorem}[Boundary walks of the base-$\beta$ target at $m=1$]\label{thm:walks}
Let $(a,b,c)=(ef,\,f^2-e^2,\,f^2)$ be the primitive $3\alpha+2\beta$ tile ($1\le e<f$,
$\gcd(e,f)=1$) and let $m=1$, so that the target is $(f^3,f^3,e(3f^2-e^2))$ and $N=3f^2-e^2$. Write a
walk along a side as $P\cdot a+Q\cdot b+R\cdot c$ equal to that side's length. Then both walk sets are
cut out by the \emph{same} linear form $\langle e,f,1\rangle$, the base at level $2e$ and the equal
side at level $f$:
\[
\begin{array}{llll}
\text{base:} & (P,Q,R)=(je+fp,\;e+fj,\;R), &\quad pe+jf+R=2e, &\quad j\in\mathbb{Z}_{\ge0},\\[2pt]
\text{side:} & (P,Q,R)=(qe+fp,\;fq,\;R), &\quad pe+qf+R=f,  &\quad q\in\mathbb{Z}_{\ge0},
\end{array}
\]
with $p\in\mathbb{Z}$ and $P,R\ge0$. In the thin regime $f>2e$ the levels force short lists:
\begin{enumerate}
\item[\textup{(i)}] \textup{(trichotomy)} the base is exactly one of
$\{b^{e},c^{2e}\}$, $\{a^{f},b^{e},c^{e}\}$, $\{a^{2f},b^{e}\}$; in particular it carries
\emph{exactly $e$} $b$-edges;
\item[\textup{(ii)}] \textup{(dichotomy)} each equal side is $\{a^{e},b^{f}\}$ or
$\{a^{fp},c^{\,f-pe}\}$ for some $p\ge0$.
\end{enumerate}
Feeding in the $\gamma$-trap ($R\ge1$ on every side, Remark~\ref{rem:e2}) removes the thinness
hypothesis altogether. For \emph{every} $(e,f)$ at $m=1$:
\begin{enumerate}
\item[\textup{(iii)}] \emph{no equal side carries a $b$-edge} --- each is $\{a^{fp},c^{\,f-pe}\}$,
$p\ge0$; so \emph{all} boundary $b$-edges lie on the base;
\item[\textup{(iv)}] the base's $b$-count $Q=e+fj$ obeys $j(f-e)\le e-1$.
\end{enumerate}
In particular $e=1$ forces $j=0$ for every $f$ --- \emph{the base carries exactly one $b$-edge} --- and
$f>2e$ forces $j=0$, recovering \textup{(i)} minus its $R=0$ member.
\end{theorem}

\begin{proof}
Base. Reducing $P\cdot ef+Q(f^2-e^2)+R f^2=e(3f^2-e^2)$ mod $f$ gives $Qe^2\equiv e^3$, so
$f\mid Q-e$ by $\gcd(e,f)=1$; as $0\le Q$ and $0<e<f$, the least nonnegative residue is $e$, whence
$Q=e+fj$ with $j\ge0$. Substituting and cancelling one factor $f$ leaves the \emph{halved equation}
\begin{equation}\label{eq:halfbase}
Pe+j(f^2-e^2)+Rf=2ef .
\end{equation}
Reducing \eqref{eq:halfbase} mod $f$ gives $e(P-je)\equiv0$, so $f\mid P-je$, say $P=je+fp$,
$p\in\mathbb{Z}$; substituting into \eqref{eq:halfbase} and cancelling $f$ gives $pe+jf+R=2e$.
Conversely every such triple solves the original equation. For $f>2e$: from \eqref{eq:halfbase},
$j(f^2-e^2)\le2ef$; if $j\ge2$ then $f^2-e^2\le ef$, i.e. $f^2-ef-e^2\le0$, whereas
$(f-2e)(f+e)=f^2-ef-2e^2>0$ gives $f^2-ef-e^2>e^2>0$ --- contradiction.
If $j=1$ then $f\mid P-e$ forces $P\ge e$, and \eqref{eq:halfbase} gives
$e^2+(f^2-e^2)\le2ef$, i.e. $f^2\le2ef$, i.e. $f\le2e$ --- contradiction. So $j=0$, whence $f\mid P$,
$P=fp$ with $p\ge0$, and $pe+R=2e$ forces $p\in\{0,1,2\}$ and $R=e(2-p)$, giving the three walks.

Side. The same reduction on $P\cdot ef+Q(f^2-e^2)+Rf^2=f^3$ gives $Qe^2\equiv0$, so $f\mid Q$,
$Q=fq$; cancelling $f$ leaves $Pe+q(f^2-e^2)+Rf=f^2$, whose reduction mod $f$ gives $f\mid P-qe$,
say $P=qe+fp$, and cancelling $f$ again gives $pe+qf+R=f$. For $f>2e$: $q\ge2$ would give
$2(f^2-e^2)\le f^2$, i.e. $f^2\le 2e^2<f^2$, absurd. If $q=1$ then $Pe+Rf=e^2$ and $f\mid P-e$ force
$P=e$, $R=0$; if $q=0$ then $f\mid P$, $P=fp$ and $pe+R=f$.

(iii) and (iv) are one computation, and need no thinness. Suppose a side walk has $q\ge1$. From
$P=qe+fp\ge0$ we get $fp\ge-qe>-qf$ (using $e<f$ and $q\ge1$), so $p>-q$, i.e. $p\ge1-q$. The level
equation then gives
\[
R \;=\; f-pe-qf \;\le\; f-(1-q)e-qf \;=\; (1-q)(f-e)\;\le\;0,
\]
contradicting $R\ge1$. Hence $q=0$, which is (iii). Identically for the base with $P=je+fp\ge0$: if
$j\ge1$ then $p\ge1-j$ and
\[
R \;=\; 2e-pe-jf \;\le\; 2e-(1-j)e-jf \;=\; e-j(f-e),
\]
so $R\ge1$ forces $j(f-e)\le e-1$, which is (iv) (and it holds trivially at $j=0$ since $e\ge1$). For
$e=1$ this reads $j(f-1)\le0$, so $j=0$; for $f>2e$ we get $je<j(f-e)\le e-1$, so $j=0$.
\end{proof}

\begin{theorem}[Apex mismatch: the pierced corner]\label{thm:pierce}
Let a base-$\beta$ target be tiled at any scale $m$ such that both equal sides end (at the apex) with
a $c$-edge --- automatic at $m=1$, where the sides are $b$-free. Then, writing $T_2$ for the middle
of the three apex tiles:
\begin{enumerate}
\item[\textup{(i)}] the outer apex tiles carry $c$ on the boundary and $b$ inward, and $T_2$ carries
$b$ and $c$ inward, so exactly one inner apex ray carries $T_2$'s $c$-edge against a neighbour's
$b$-edge: a T-junction at distance $b$ from the apex, whose vertex figure on the neighbour's side is
$(1,1,1)$ (one $\alpha$, one $\beta$, plus the neighbour's $\gamma$);
\item[\textup{(ii)}] the remaining segment of $T_2$'s $c$-edge has length $c-b=e^2$, and
$n_a\,ef+n_b(f^2-e^2)+n_c f^2=e^2$ has \emph{no} solution in nonnegative integers, for any
$1\le e<f$ with $\gcd(e,f)=1$;
\item[\textup{(iii)}] hence some far-side edge crosses the far endpoint of $T_2$'s $c$-edge
\emph{straight through}: every such tiling contains a \textbf{pierced $\beta$-corner} at distance
$f^2$ along an inner apex ray, and the angles on $T_2$'s side of the piercing line beyond the corner
form $\{\alpha,\gamma\}$ or $\{3\alpha,\beta\}$.
\end{enumerate}
\end{theorem}

\begin{proof}
(i) The apex is $3\alpha$ (three tiles, each $\alpha$ there); the two edges of an $\alpha$-corner are
$b$ and $c$. The boundary edges at the apex are the sides' last edges, $c$ by hypothesis, so the outer
tiles point $b$ inward and $T_2$ has $\{b,c\}$ inward. Whichever inner ray pairs $T_2$'s $c$ with an
outer $b$, the $b$-edge ends at distance $b$: its far endpoint is the $\gamma$-corner of its tile
($b$ joins $\alpha$ to $\gamma$), and the vertex figure on that side of the straight $c$-edge sums to
$\pi$ with a $\gamma$ present, hence is $(1,1,1)$ by the vertex classification. (ii) Any solution has
$n_a=n_c=0$ since $ef,f^2>e^2$; then $n_b(f^2-e^2)=e^2$ gives $n_b f^2=(n_b+1)e^2$, and
$\gcd(e,f)=1$ forces $f^2\mid n_b+1$, so $n_b\ge f^2-1$ and
$e^2=n_b(f^2-e^2)\ge f^2-1\ge e^2$, with equality forcing $f^2-e^2=1$ --- impossible, as
$f^2-e^2\ge 2e+1\ge 3$. (iii) The far side of the leftover segment is covered by consecutive whole
tile edges starting at the T-junction; by (ii) no partial sum lands exactly on the segment's far
endpoint, so some edge strictly crosses it. That endpoint is the $\beta$-corner of $T_2$ ($c$ joins
$\alpha$ to $\beta$). The angles on $T_2$'s side of the crossing line beyond the corner fill
$\pi-\beta$; writing them as $x\alpha+y\beta+z\gamma$ and clearing $\beta,\gamma$, irrationality of
$\alpha/\pi$ forces $y+z=1$, $2x+z=3y+3$, whose only solutions are $(3,1,0)$ and $(1,0,1)$.
\end{proof}

\begin{remark}[Validation on the $44$-tiling]\label{rem:pierce}
The theorem can be checked against the genuine $44$-tiling ($m=2$; its equal sides do end with $c$),
where it holds exactly: the apex rays pair one $b\,|\,b$ and one mismatch $c\,|\,b$; the middle
tile's $c$-endpoint $V=(10,2\sqrt{15})$ is pierced by a straight $b$-edge, and the sector on the
middle tile's side is $\beta+\{\alpha,\gamma\}$ --- the first continuation type. Parts (ii) and the
vertex figures in (i), (iii) are machine-checked (\texttt{apex\_leftover\_nonrepresentable},
\texttt{pierced\_corner\_types}, \texttt{BaseBetaWalks.lean}). The configuration is an interior
constraint, uniform in $(e,f)$ and unconditional at $m=1$: every hypothetical thick-regime tiling,
$N=59$ included, must contain it.
\end{remark}

\begin{proposition}[Rung two: the pre-piercer chain]\label{prop:rung2}
In the configuration of Theorem~\ref{thm:pierce}, every full far-side edge strictly between the
T-junction and the piercer lies inside $[b,f^2]$, hence has length $\le e^2$; among the tile edges
only $b=f^2-e^2$ can satisfy this, and only in the \emph{super-thick} regime $f^2\le 2e^2$. So the
piercer starts at $(t+1)b$ with $tb<e^2$, where $t=0$ outside super-thick; and if the piercer is
itself a $b$-edge, its corner at its starting junction is exactly $\alpha$ (the $(1,1,1)$ figure
offers $\{\alpha,\beta\}$, while a $b$-edge's endpoints are $\{\alpha,\gamma\}$). In particular for
$N=59$ ($(e,f)=(4,5)$, super-thick: $25\le32$) the piercer starts at distance $9$ or $18$, and for
every non-super-thick member it starts at the T-junction itself.
\end{proposition}

\begin{theorem}[Alignment of the mismatch ray]\label{thm:align}
Let $m=1$, $e\ge2$, and let $r$ be the mismatch inner apex ray of Theorem~\ref{thm:pierce}, of inner
length $f\cdot b$ (Proposition~\ref{prop:rung2}, $b=f^2-e^2$). Along $r$, the near side is $T2$'s
single unsplittable $c$-edge covering $[0,f^2]$, and the far side begins with $T1$'s $b$-edge. Then:
\begin{enumerate}
\item[\textup{(i)}] \emph{no common junction} lies in $(0,f^2)$: for $1\le k\le f-1$, $k\cdot b$ is
never $n_a\,ef+n_b\,b+n_c f^2$ with $n_c\ge1$ (so every far-side junction is a genuine T-junction,
the near $c$-edge passing straight through it);
\item[\textup{(ii)}] the far side is exactly $b^{f}$: the only representation of $f\cdot b$ as
$n_a\,ef+n_b\,b+n_c f^2$ with $n_b\ge1$ is $(0,f,0)$;
\item[\textup{(iii)}] $V=f^2$ is pierced: $b\nmid f^2$, so no far-side junction $k\cdot b$ equals $f^2$
and a far $b$-edge straddles $T2$'s $\beta$-corner.
\end{enumerate}
All three are machine-checked (\texttt{far\_near\_disjoint}, \texttt{far\_is\_bpow},
\texttt{b\_not\_dvd\_fsq}); (ii) alone uses $\gamma$-vertex induction (\emph{via} the walk arithmetic)
to fix the first far edge, the rest is pure number theory.
\end{theorem}

\begin{proof}
(i) If $k b = n_a ef + n_b b + n_c f^2$ then $(k-n_b)b = f(n_a e + n_c f)$, so $f\mid(k-n_b)e^2$ and
$\gcd(e,f)=1$ give $f\mid k-n_b$; writing $k-n_b=fs$, the equation reads $sb=n_a e+n_c f\ge f>0$
(as $n_c\ge1$), so $s\ge1$ and $k\ge n_b+f\ge f$, against $k\le f-1$. (ii) $f\mid n_b$ by the same
reduction, and $n_b\le f$ by size, so $n_b\in\{0,f\}$; $n_b\ge1$ forces $n_b=f$, and then
$n_a ef+n_c f^2=0$ gives $n_a=n_c=0$. (iii) $b\mid f^2$ would give $b\mid e^2$, hence
$b\mid\gcd(f^2,e^2)=1$, so $b=1$, impossible as $b\ge2e+1\ge5$. The far junctions are the multiples
$kb$, none equal to $f^2$, so the $b$-edge containing $f^2$ straddles it.
\end{proof}

\begin{theorem}[The mismatch ray is completely determined]\label{thm:ray}
Let $m=1$, $e\ge2$, and let the family be thick, $f\le2e$. Then along the mismatch inner apex ray,
of inner length $f\,b$:
\[
\text{far side}=b^{\,f}\ (f\text{ edges}),\qquad
\text{near side}=\{a^{\,f-e},\,c^{\,f-e}\}\ (2(f-e)\text{ edges}),
\]
and the near side is the \emph{unique} possibility. In particular the near side carries equally many
$a$- and $c$-edges, and no $b$-edge.
\end{theorem}

\begin{proof}
The far side is Theorem~\ref{thm:align}(ii). For the near side, $T2$'s $c$-edge covers $[0,f^2]$ and
the remainder $f\,b-f^2=f(b-f)$ is covered by whole edges, say $n_a$ of length $a$, $n_b$ of length
$b$ and $n_c$ of length $c$. Reducing modulo $f$ kills $n_b$ as in Theorem~\ref{thm:align}(i), and
cancelling one factor $f$ leaves
\[
n_a e+n_c f+f=b .
\]
Now $(n_a,n_c)=(f-e,\,f-e-1)$ solves it, since $(f-e)e+(f-e-1)f=f^2-e^2-f=b-f$. For uniqueness,
$\gcd(e,f)=1$ and the equation give $f\mid n_a+e$, so $n_a+e=ft$; $t=0$ is impossible and $t\ge2$
would force $n_ae\ge(2f-e)e$, whence $2ef\le f^2-f$, i.e. $f\ge2e+1$, contradicting thickness. Hence
$t=1$, $n_a=f-e$, and then $n_cf=f(f-1-e)$ gives $n_c=f-e-1$. Adding $T2$'s own $c$-edge produces the
multiset $\{a^{f-e},c^{f-e}\}$, whose total length is $(f-e)(ef+f^2)=f(f-e)(e+f)=f\,b$ as required.
\end{proof}

\begin{theorem}[Unique representation of $k\,b$]\label{thm:kbrep}
For $1\le k\le f-1$ the only way to write $k\,b$ as $n_a\,ef+n_b\,b+n_c f^2$ in nonnegative integers
is $b^{k}$. \textup{(}Machine-checked, \texttt{kb\_unique\_rep}; it generalizes
Theorem~\ref{thm:align}\textup{(ii)}, which needed $n_b\ge1$.\textup{)}
\end{theorem}

\begin{theorem}[The angle $(\pi-\alpha)/2$]\label{thm:halfangle}
A vertex figure filling $(\pi-\alpha)/2$ is $\alpha+\beta$, and nothing else.
\textup{(}Machine-checked, \texttt{halfpi\_minus\_alpha\_unique}.\textup{)}
\end{theorem}

\begin{theorem}[The $b$-run orientation lemma]\label{thm:chain}
Let a run of $b$-edges be laid end to end along a ray, seeded at an apex tile. Then every $b$-edge of
the run carries $\alpha$ at its apex end and $\gamma$ at its far end, and every junction interior to
the run is an $(\alpha,\beta,\gamma)$ $\pi$-vertex. In particular the $3\alpha+2\beta$ figure never
occurs along such a run.
\end{theorem}

\begin{proof}
A $b$-edge joins its tile's $\alpha$- and $\gamma$-corners, so at a junction the incoming tile
presents $\alpha$ or $\gamma$, and likewise the outgoing one. The $3\alpha+2\beta$ figure contains no
$\gamma$, and two $\gamma$'s cannot share a side of a straight line since $2\gamma=\pi+\alpha>\pi$.
Hence if the incoming corner is $\gamma$ the figure is $(1,1,1)$, whose single $\alpha$ must be the
outgoing corner, so the next edge again ends in $\gamma$. An apex tile has $\alpha$ at the apex, hence
$\gamma$ at the far end of its $b$-edge, which seeds the induction; it proceeds with no branching.
\end{proof}

\begin{theorem}[The far region is a scaled tile]\label{thm:farregion}
Suppose the mismatch apex ray is a complete wall, of inner length $f\,b$. Then it cuts off a triangle
with angle $\alpha$ at the apex, $\beta$ at the base corner and hence $\gamma$ at its third vertex,
and with sides $f\,c,\ f\,a,\ f\,b$ --- that is, \emph{the tile scaled by $f$}. It therefore has area
$f^2$ times the tile's and is tiled by the standard $f$-subdivision into $f^2$ copies, whose ray side
is exactly $b^{f}$ with the junction chain of Theorem~\ref{thm:chain}.
\end{theorem}

\begin{corollary}[The far side can never give a contradiction]\label{cor:farvacuous}
\emph{Under the hypothesis of Theorem~\textup{\ref{thm:farregion}}} --- that the mismatch apex ray is
a complete wall --- no obstruction to a base-$\beta$ tiling can be derived from the far side of that
ray, at any $(e,f)$: that side admits a genuine tiling. The hypothesis may not be dropped; if the ray
is not a complete wall the far region is not a scaled copy of the tile and the conclusion does not
follow.
\end{corollary}

\begin{remark}[Where the obstruction must instead live]\label{rem:live}
Corollary~\ref{cor:farvacuous} is a negative result and an important one: it explains why the
pierced-corner analysis resisted three independent attempts. The configuration it pins down is not
contradictory but \emph{realizable}, so no amount of further chasing on that side can close the
branch. The check is exact for every family --- for $(e,f)=(4,5)$ the far region has sides
$(45,100,125)=5\cdot(9,20,25)$ and contains $25=f^2$ of the $59$ tiles --- and it is confirmed on the
genuine $44$-tiling, whose mismatch ray has \emph{no} straddling tile and whose far region contains
$16=(mf)^2$ tiles with sides $4\cdot(c,a,b)$.

The live direction is the other ray. For every thick member with $m=1$, $e\ge2$, the \emph{matched}
inner apex ray is \emph{not} a wall. Indeed if it were, the target would split into $f^2+f^2+(f^2-e^2)$
tiles, the last a middle triangle $M$ with apex $\alpha$, legs $f\,b$ and base $e\,b$. Since $k\,b$
has the unique representation $b^{k}$ for $1\le k\le f-1$ (Theorem~\ref{thm:kbrep}), both edges of $M$
at its far corner are $b$-edges; but the angle there is $(\pi-\alpha)/2$, whose only decomposition is
$\alpha+\beta$ (Theorem~\ref{thm:halfangle}), and a $b$-edge endpoint is an $\alpha$- or
$\gamma$-corner, so no tile present can supply the $\beta$. This is the precise mechanism behind the
failure of the clean-ray corollary retracted in Remark~\ref{rem:cleanfail}, and it is confirmed on the
$44$-tiling, whose matched ray carries exactly three straddling tiles against the mismatch ray's zero.

Consequently the branch now reduces to a single question: \emph{can the residual region --- the
target minus the far wall, containing $N-f^2$ tiles and crossed by a necessarily straddled matched
ray --- be tiled?} Every tool developed here (walk arithmetic, vertex figures, unsplittability) is
wall-based and goes silent once a tile straddles. What is needed is a bound on where a tile may
straddle a ray whose two ends are pinned by $b$-edges.
\end{remark}

\begin{remark}[Tiles do straddle: the no-straddle hypothesis is false]\label{rem:straddle}
Remark~\ref{rem:live} asks for a bound on where a tile may straddle a ray whose two ends are pinned
by $b$-edges. The strongest form one might hope for is that no such straddle occurs at all:
\[
(\mathrm{DL})\quad\text{if a segment carries a whole } b\text{-edge at each end, no tile straddles it.}
\]
$(\mathrm{DL})$ would make the walk arithmetic of \S\ref{sub:funnels} applicable to interior
segments, and it is the form the funnel arguments implicitly assume. It is false.

Testing it against the three tilings of base-$\beta$ targets exhibited in this note and the main
paper --- all with tile $(2,3,4)$, all verified exactly --- for every $b$-edge we take the line
through it and count tiles with interior on both sides whose centroid projects near the segment:

\begin{center}
\begin{tabular}{lccc}
\toprule
tiling & $b$-edges & lines with a straddling tile & worst \\
\midrule
$44$-tiling \texttt{N44B} of $(16,16,22)$, $m=2$ & $44$ & $19$ & $3$\\
$44$-tiling \texttt{N44C} of $(16,16,22)$, $m=2$ & $44$ & $30$ & $4$\\
$99$-tiling of $(24,24,33)$, $m=3$              & $99$ & $64$ & $4$\\
\bottomrule
\end{tabular}
\end{center}

Straddling near $b$-edges is common rather than exceptional. Consequently the missing ingredient is
not a prohibition but a \emph{quantitative} bound: a description of the set in which a straddling
tile's crossing points must lie. Since every vertex of a tiling lies in a finitely generated
$\mathbb Z$-module over $\Q(\sqrt D)$, and $b$-pinning fixes the segment's endpoints inside that
module, such a bound is the natural object to seek; we have not found it. The computation is
\texttt{rust/walls} (binary \texttt{straddle}), and it is recorded here because a negative result of
this shape saves the reader the attempt, and because it explains why the interior argument of
\S\ref{sub:funnels} has never been written out.
\end{remark}


\begin{remark}[Two corrections this forces]\label{rem:cleanfail}
Theorem~\ref{thm:chain} is uniform in $(e,f)$ and in $m$, and is confirmed on the genuine
$44$-tiling, where all four far $b$-edges of the mismatch ray and both sides of the clean ray read
$\alpha\to\gamma$. It strictly generalizes Proposition~\ref{prop:rung2}, which established only the
first step. But it also refutes two statements of an earlier version, which we withdraw.

First, the assertion that the \emph{clean} apex ray is $b^{f}$ against $b^{f}$. Its justification ---
``by the symmetry of the isosceles target both inner rays have the same inner length $f\,b$'' ---
conflates two objects. The symmetry does apply to the geometric rays; it does not apply to the
tiling's maximal edge \emph{segment} along a ray, which may terminate early at a vertex where the
edges turn away. In fact the clean ray's segment can never reach the base: if it did, both final
$b$-edges would contribute $\gamma$ at an interior point of the base, where the tiles above sum to
$\pi$, contradicting $2\gamma>\pi$. On the $44$-tiling the two segments have lengths $12$ and $9$
respectively, and the clean ray terminates at an interior $\beta+3\gamma$ vertex --- which is
precisely the $M_4$ vertex the census forces.

Second, and more seriously, the same objection applies to the hypothesis under which
Theorems~\ref{thm:align} and~\ref{thm:ray} were stated: that the \emph{mismatch} ray's maximal
segment attains the full length $f\,b$. That too was justified by symmetry, and it too now requires a
genuine proof. The arithmetic of those theorems is unaffected --- given a chain of total length
$f\,b$ containing a $b$-edge, the conclusions stand and remain machine-checked --- but their
geometric hypothesis must be stated explicitly. Establishing that the mismatch segment reaches the
base is the single highest-value open step in this branch: it would restore Theorems~\ref{thm:align}
and~\ref{thm:ray} unconditionally for every thick member.
\end{remark}

\begin{remark}[Scope and consequences]\label{rem:ray}
Thickness is used exactly once, and sharply: the next solution would be $n_a=2f-e$, which needs
$f\ge2e+1$ --- the thin regime. So in the regime that is open, the entire mismatch ray is now pinned
on \emph{both} sides, with no freedom beyond the order of the near-side edges. Explicitly:
$N=23$ has far $5^3$, near $\{6,9\}$; $N=59$ far $9^5$, near $\{20,25\}$; $N=66$ far $16^5$, near
$\{15^2,25^2\}$; $N=83$ far $11^6$, near $\{30,36\}$; $N=167$ far $39^8$, near $\{40^3,64^3\}$.
The count of interior junctions is therefore also determined --- $f-1$ on the far side and
$2(f-e)-1$ on the near side, none shared (Theorem~\ref{thm:align}(i)) --- so a hypothetical thick
tiling carries exactly $3f-2e-2$ T-junctions along this ray, each a $\pi$-vertex. The arithmetic is
machine-checked (\texttt{near\_side\_unique}, \texttt{BaseBetaWalks.lean}).
\end{remark}

\begin{remark}[Scope, and why $m=1$ is the rigid case]\label{rem:walkscope}
Theorem~\ref{thm:walks}(i) generalizes the $e=1$ base analysis used for
Theorem~\ref{thm:basebeta-e1} to the infinite family $f>2e$ --- $e=1$, $f\ge3$; $e=2$, $f\ge5$;
$e=4$, $f\ge9$; \dots --- covering the primes $N=47,71,107,191,227,239,359,\dots$, and its parts (iii)--(iv),
which need no thinness, deliver the $e\ge2$ analogue of the $e=1$ structure theorem of
Remark~\ref{M-rem:lem6} for \emph{every} $(e,f)$: the equal sides are $b$-free and the base's $b$-count
is bounded by $j(f-e)\le e-1$. The step that fails for $m\ge2$ is the very first one: there
$Q\equiv em\pmod f$, so $j$ may be \emph{negative} and the level rises to $2em$. This is not an
artifact. The genuine $44$-tiling ($e{=}1,f{=}2,m{=}2$) has base walk $aaaaccca$, i.e.
$(P,Q,R)=(5,0,3)$, which is the parameter $j=-1$, $p=3$; and the $99$-tiling
($e{=}1,f{=}2,m{=}3$) has base walk $aabbbbbbbcc$, i.e. $(2,7,2)$, the parameter $j=2$, $p=0$. Both
satisfy $R\ge1$, as the $\gamma$-trap requires.
\end{remark}

\begin{theorem}[Walk structure at $m=1$]\label{thm:walkstruct}
Let $m=1$, $\gcd(e,f)=1$, $1\le e<f$, and consider a tiling of the base-$\beta$ target.
\begin{enumerate}
\item[\textup{(i)}] If $f^2>2e^2$, then each equal side carries \emph{no} $b$-edge. Its edges are
$a$'s and $c$'s with $n_a=fk$ and $n_c=f-ke\ge1$, and it ends with a $c$-edge at the apex.
\item[\textup{(ii)}] If $f^2>2ef+e^2$ --- equivalently $f/e>1+\sqrt2$ --- then the base carries
\emph{exactly} $e$ $b$-edges, and $n_a=f\ell$, $n_c=e(2-\ell)$ with $\ell\in\{0,1\}$. In particular
\[
R_{\mathrm{base}}\in\{e,\,2e\}.
\]
\end{enumerate}
\end{theorem}

\begin{proof}
A side of length $L$ is exactly partitioned by the tile edges on it, giving
$n_a\,ef+n_b\,b+n_c\,f^2=L$ with $n_a,n_b,n_c\ge0$, and $n_c\ge1$ by
Proposition~\ref{M-prop:gammatrap}.

\emph{(i)} Here $L=f^3$. Rearranging with $b=f^2-e^2$ gives
$n_b e^2=f\,(n_a e+n_b f+n_c f-f^2)$, so $f\mid n_b e^2$ and hence $f\mid n_b$ by coprimality;
write $n_b=ft$. If $t\ge2$ the $b$-edges alone cost at least $2fb=2f^3-2fe^2$, which together with
$n_cf^2\ge f^2$ exceeds $f^3$ unless $f^2\le2e^2$. If $t=1$, cancelling $f$ leaves
$n_ae+n_cf=e^2$, so $n_a\le e$ and $f\mid e(e-n_a)$; coprimality with $0\le e-n_a\le e<f$ forces
$n_a=e$ and then $n_c=0$, contradicting $n_c\ge1$. So $t=0$. With $n_b=0$ the equation is
$n_ae+n_cf=f^2$, whence $f\mid n_a$; writing $n_a=fk$ gives $n_c=f-ke$. Finally, the apex tile on
that side presents $\alpha$ (Proposition~\ref{M-prop:cornerfig}), whose incident edges are $b$ and
$c$; as the side has no $b$-edge, it ends with a $c$-edge.

\emph{(ii)} Here $L=e(3f^2-e^2)$, and the same rearrangement gives
$(n_b-e)e^2=f\,(n_ae+n_bf+n_cf-3ef)$, so $f\mid n_b-e$; write $n_b=e+ft$. Negative $t$ makes
$n_b<0$. For $t\ge1$ the $b$-edges alone cost at least
\[
(e+f)b=(f+e)^2(f-e)=e(3f^2-e^2)+f\,(f^2-2ef-e^2),
\]
which already exceeds $L$ under the hypothesis, before the mandatory $c$-edge is placed. So
$n_b=e$; subtracting and cancelling $f$ leaves $n_ae+n_cf=2ef$, whence $f\mid n_a$, and writing
$n_a=f\ell$ gives $n_c=e(2-\ell)$. Then $n_c\ge1$ forces $\ell\le1$.

Machine-checked (\texttt{equal\_side\_no\_b}, \texttt{equal\_side\_shape}, \texttt{base\_b\_count}
in \texttt{lean/BaseBetaWalkArith.lean}), and cross-checked by enumerating all walk solutions for
the $198$ coprime pairs with $f\le25$: no discrepancies.
\end{proof}

\begin{remark}\label{rem:walkscope2}
Both hypotheses are sharp: $n_b=2f$ on an equal side is admissible exactly when $f^2\le2e^2$, and
$n_b=e+f$ on the base exactly when $f^2\le2ef+e^2$, by the displayed identity. Of the $42$
base-$\beta$ prime candidates below $1000$, $28$ satisfy \textup{(i)} --- including the thick values
$23$, $131$ and $167$, since \textup{(i)} needs only $f/e>\sqrt2$ --- and $17$ satisfy
\textup{(ii)}. The theorem structures those cases; it does not by itself exclude any of them.
\end{remark}

\begin{theorem}[The $e=1$ subfamily reduces to one walk]\label{thm:e1reduce}
Let $e=1$, $m=1$ and $f\ge3$, so the target is $(f^3,f^3,3f^2-1)$, the tile is $(f,f^2-1,f^2)$ and
$N=3f^2-1$. Both hypotheses of Theorem~\ref{thm:walkstruct} hold ($f^2>2$ and $f^2>2f+1$). In any
tiling:
\begin{enumerate}
\item[\textup{(i)}] neither equal side carries a $b$-edge; each equal side begins and ends with a
$c$-edge, and therefore carries at least two;
\item[\textup{(ii)}] the base carries exactly one $b$-edge and exactly one $c$-edge --- so
$R_{\mathrm{base}}=1$ --- and its first and last edges are $a$-edges.
\end{enumerate}
Thus the base walk is a permutation of $(a^f,b,c)$ beginning and ending with $a$, with the $b$ in
position $3,\dots,f$.
\end{theorem}

\begin{proof}
\emph{(i)} An equal side has length $f^3$ and its edges lie in $\{f,\ f^2-1,\ f^2\}$. Reducing
$n_af+n_b(f^2-1)+n_cf^2=f^3$ modulo $f$ gives $f\mid n_b$. If $n_b\ge2f$ then
$n_b(f^2-1)\ge2f^3-2f>f^3$ for $f\ge2$. If $n_b=f$ then $f(f^2-1)=f^3-f$ leaves $f$, forcing
$n_c=0$ and $n_a=1$, hence no $c$-edge on that side --- contradicting
Proposition~\ref{M-prop:gammatrap}. So $n_b=0$, and then $n_a+n_cf=f^2$, so $f\mid n_a$; write
$n_a=fk$, $n_c=f-k$.

By Proposition~\ref{M-prop:cornerfig} the apex tile on that side presents $\alpha$, whose incident
edges are $b$ and $c$; since the side carries no $b$-edge, it ends with a $c$-edge. Its other end is
handled in \emph{(ii)}: there we show the base begins and ends with $a$-edges, so the base-corner
tile's $c$-edge lies on the equal side, which therefore begins with a $c$-edge. Were $n_c=1$ the
single $c$-edge would be both the first and the last, forcing the side to consist of one edge and
$f^3=f^2$; so $n_c\ge2$.

\emph{(ii)} The base has length $3f^2-1$. Reducing modulo $f$ gives $n_b\equiv1$. If $n_b\ge f+1$
then $n_b(f^2-1)\ge f^3+f^2-f-1>3f^2-1$ for $f\ge3$, so $n_b=1$. Subtracting,
$2f^2=n_af+n_cf^2$, so $n_a=f(2-n_c)$ with $n_c\in\{0,1,2\}$, and $n_c\ge1$ by
Proposition~\ref{M-prop:gammatrap}.

If $n_c=2$ then $n_a=0$ and the base is three edges, a permutation of $(b,c,c)$. By
Proposition~\ref{M-prop:cornerpara} the first two edges lie in $\{a,c\}$ and so do the last two; with
no $a$-edge available this forces all three to be $c$-edges, contradicting $n_b=1$. Hence $n_c=1$
and $n_a=f$, so the base is a permutation of $(a^f,b,c)$ with $n=f+2$ edges, and
Proposition~\ref{M-prop:cornerpara} places the $b$ away from the first two and last two positions.

It remains to exclude $E_1=c$. Suppose it. Then the unique $c$-edge is $E_1$, so every one of
$E_2,\dots,E_n$ is an $a$- or $b$-edge and places a $\gamma$ at one of its endpoints; there are
$f+1$ of them and $f+1$ interior junctions $J_1,\dots,J_{n-1}$. The corner tile $T_A$ has its
$c$-edge on the base, so its $\alpha$ vertex is $J_1$ and its $\gamma$ vertex lies on the equal
side; by Proposition~\ref{M-prop:cornerpara} its partner $T_3$ presents $\gamma$ at $J_1$. So $J_1$
already carries a $\gamma$, and since a $\pi$-vertex carries at most one
(Proposition~\ref{M-prop:vertexfigures}), $E_2$ must place its $\gamma$ at $J_2$; inductively $E_i$
places at $J_i$, and $E_n$ places at $J_n$, a corner --- contradicting
Proposition~\ref{M-prop:cornerfig}. Hence $E_1$ is an $a$-edge, and the same argument at $B$ gives
$E_n$ an $a$-edge. This is what \emph{(i)} used.
\end{proof}

\begin{remark}\label{rem:e1status}
Theorem~\ref{thm:e1reduce} reduces the whole $e=1$, $m=1$ subfamily --- which contains the base-$\beta$
prime candidates $N=3f^2-1$, among them $11$, $47$, $107$, $191$, $431$, $587$ and $971$ below
$1000$ --- to the single question of whether that one base walk is realizable. It also corrects a
plausible-looking plan: one might hope to close the subfamily by proving $R\ge2$ on every side. That
is false for the base. The walk with $R_{\mathrm{base}}=2$ is exactly the one the corner
parallelogram kills, and the survivor has $R_{\mathrm{base}}=1$; on the base, ``$R\ge2$'' is not a
lemma reducing the problem but a restatement of it. On the two equal sides $R\ge2$ does hold, and is
proved above.
\end{remark}

\begin{remark}[The $e\ge2$ residue]\label{rem:e2}
The $e\ge2$ residue ($N=23,59,71,83,131,167,\dots$; the members $23$ and $71$ are now settled by search, Table~\ref{M-tab:basebeta}) resists all three natural invariant methods, and in two
cases \emph{provably}:
\begin{itemize}
\item \textbf{Linear/colouring invariants} --- dead by Remark~\ref{M-rem:nogo}.
\item \textbf{$b$-edge census.} On a maximal segment of length $L$ the two sides carry $q_1\equiv q_2
\pmod f$ (both $\equiv -e^{-2}L$), and the entire census --- in \emph{every} modulus, and in its
per-direction refinement --- is a formal consequence of the single parity identity
$N(a+b+c)\equiv2X+Y \pmod 2$, valid for all $(e,f,m)$. So no congruence of this type can bite.
\item \textbf{Euler / vertex counting.} For any such tiling $V-E+F=2$ is \emph{identically} the sum
of the three corner identities $\#\alpha=\#\beta=\#\gamma=N$; the vertex-type system collapses to
$Q_{013}=1+P_{320}+2Q_{640}+Q_{431}+W_{320}$ and
$3P_{320}+P_{111}+6Q_{640}+4Q_{431}+2Q_{222}+3W_{320}+W_{111}=N-3$, which is
feasible for every $e\ge2$, $m=1$ target. So Euler carries no information here.
Here $W_{320}$ and $W_{111}$ count the interior vertices at which some tile is met in the
\emph{relative interior} of one of its edges: such a tile contributes a straight angle, so the
remaining tile angles sum to $\pi$ rather than $2\pi$ and the figure is $(3,2,0)$ or $(1,1,1)$
(Proposition~\ref{M-prop:vertexfigures}). These types are not optional --- the dissections here are
not edge-to-edge (Remark~\ref{M-rem:parityfail}) --- and an earlier version of this census omitted
them. Their presence only enlarges the feasible set, so the conclusion is unaffected.
\end{itemize}
What \emph{is} gained is citation-free structure, valid for all $(e,f,m)$: the $\gamma$-trap
(Proposition~\ref{M-prop:gammatrap}) and the corner parallelogram
(Proposition~\ref{M-prop:cornerpara}), both proved above from the corner figures and the
unsplittability of $b$, and both used by the exact search as prunes. They were checked against the
two genuine certificates before being proved: the first-two-edges rule of
Proposition~\ref{M-prop:cornerpara} holds on all six sides of the $44$- and $99$-tilings, and at each
of the four base corners exactly one of the two sides begins with a $c$-edge.

This isolates the residue to one clean statement. Since $q_{\mathrm{base}}\equiv em \pmod f$ while
$q_{\mathrm{equal}}\equiv0$, we have:
\[
\textbf{if the base carries no } b\textbf{-edge, then } f\mid em,\ \text{hence } f\mid m,
\ \text{hence no tiling with } m=1 \ \text{--- for every } e\ \text{at once.}
\]
The $44$-tiling ($m=2$) indeed carries \emph{zero} $b$-edges on its entire boundary, consistent with
this. But Remark~\ref{M-rem:thm14false} now shows the converse direction is a dead end: the
$99$-tiling has $f\nmid m$ and its base carries seven $b$-edges, so ``the base never carries a
$b$-edge'' is \emph{false} in general, and the displayed implication --- whose hypothesis is
equivalent to $f\mid m$ --- cannot be used to exclude $m=1$ without already knowing $f\mid m$, which
is false. A fourth method also fails, and provably: \emph{boundary counting}. Each side $S$ of $ABC$, cut into
$k_S$ whole edges, forces at least $2k_S-1$ tiles to meet it (the $k_S$ inner tiles, plus a distinct
``middle'' tile at each of the $k_S-1$ junctions, since a triangle has no two collinear edges); this
is exactly tight on the $44$-tiling (base $k=8$, $15=2\cdot8-1$ tiles). But at $m=1$ the boundary is
\emph{minimal} --- for $e=1$ the least feasible total is $\sum k_S=3f+2$ --- so counting sees only
$O(f)$ of the $N=\Theta(f^2)$ tiles, and $\sum(2k_S-1)=6f+1$ exceeds $N=3f^2-1$ only at $f=2$, the
case already closed. The ratio $(6f+1)/(3f^2-1)\to0$: $m=1$ is the case \emph{least} constrained by
counting, not the most.

So no \emph{invariant} route survives: linear invariants, the census, Euler and counting are all
provably vacuous, and the one arithmetic statement that would settle it ($f\mid m$) is now known to
be false. What does survive is \emph{structure}. Theorem~\ref{thm:walks} classifies the boundary walks
for every $(e,f)$ at $m=1$, and in the thin regime $f>2e$ reduces the whole branch to two possible
bases and $b$-free equal sides --- the $e\ge2$ analogue of Remark~\ref{M-rem:lem6}, and the first
statement in this branch that is uniform in $e$. Parts (iii)--(iv) need no thinness at all: at $m=1$ \emph{no equal side carries a $b$-edge}, for every
$(e,f)$, so all boundary $b$-edges lie on the base; and the base's $b$-count $Q=e+fj$ obeys
$j(f-e)\le e-1$. It does not close the branch, but it does locate the difficulty exactly. The bound
$j\le(e-1)/(f-e)$ is uniform, and it degrades precisely as $f$ falls towards $e$: at $e=1$ it pins
$j=0$ for every $f$ (one $b$-edge on the base, the $e=1$ rigidity of Remark~\ref{M-rem:lem6},
now a corollary); at $f>2e$ it still pins $j=0$; but in the \emph{thick} regime $f\le2e$ it leaves
$\lfloor(e-1)/(f-e)\rfloor+1$ values of $j$ and the trichotomy dissolves. That is precisely where the
open members sit --- $N=59$ is $(e,f)=(4,5)$, so $j\le3$; $N=66$ is $(3,5)$, so $j\le1$; likewise
$N=23,83,131,167$. We record the branch as open and settle its members one at a time by exhaustive
search. Table~\ref{M-tab:basebeta} lists every member below $110$; of the eleven, eight are settled and
$59$ and $66$ are now excluded by exhaustive search (apex-down with the walk prune; $1\,838\,175$ and
$7\,232\,464$ nodes), leaving $83$ under search. ($N=83$, $(e,f)=(5,6)$, is the smallest base-$\beta$
prime not yet settled.)
The walk classification also feeds back into the search: the $\gamma$-trap and the corner
parallelogram have the multiset corollaries $R\ge1$ on every side, $Q_{\mathrm{base}}\le k-4$ and
$Q_{\mathrm{side}}\le k-2$ ($k$ the number of edges of the walk), and the engine now prunes any
partial boundary walk that fits no surviving multiset (prune P5). The rule was validated positionally
on all six sides of the two genuine tilings before use --- the $44$-tiling's base $aaaaccca$ and the
$99$-tiling's base $aabbbbbbbcc$, which is exactly tight for $Q\le k-4$ --- the strengthened engine
re-finds the $44$-tiling from scratch ($852{,}093$ nodes; the found tiling passes the independent
verifier: congruence, exact area, containment, pairwise disjointness) --- and the corner rule is
applied only after checking that the tile's $b$-edge is unsplittable ($n_a a + n_c c = b$ has no
solution), which is what forces the parallelogram. For $N=59$ this cuts the admissible bases from
seven to four, for $N=66$ from four to three.
\end{remark}

\section{The $e=1$ family}\label{sec:compe1}

\begin{theorem}[The $e=1$ base-$\beta$ family]\label{thm:e1family}
Let $f\ge 2$, and consider the base-$\beta$ instance $(e,f)=(1,f)$ at $m=1$: the tile
$(f,\,f^2-1,\,f^2)$ on the isosceles target with base angles $\beta$, base $Y=3f^2-1$ and equal
sides $X=f^3$, with tile count $N=3f^2-1$. No such tiling exists. In particular no prime
$p=3f^2-1$ whose only base-$\beta$ representation is $(1,f)$ is a number of congruent triangles
into which a triangle can be cut; this covers $p=11,47,107,191,431$.
\end{theorem}

The proof is a finite chain of forcing lemmas; we state the architecture and prove the three
arithmetic pivots, deferring the full strip-and-column exposition to the extended version of this
note. The chain: the walk trichotomy (base walks $(2f,1,0)$, $(f,1,1)$, $(0,1,2)$, the last dead
by Theorem~\ref{thm:e1reduce}); the corner and partner forcing; for each corner, either the
quadratic growth completes ($k=f$) or a first break occurs, on the base (then the walk letter is
forced into $\{a,c\}$, and two base-breaks are incompatible with the single $c$ of $(f,1,1)$) or
on a side. A side-break at level $j$ forces a horizontal column of $f$ tiles whose junction
fillers are exactly the triangles on consecutive apexes (Lemma~\ref{lem:filler}); the structure
recurses one level per step and reaches the base, where the apexes must be walk junctions (no
vertex of a tile can lie interior to a boundary edge), pinning the orientations
(Lemma~\ref{lem:offsets}) and forcing the letters $a^f$ on $[0,f^2]$
(Lemma~\ref{lem:interior}). In $(f,1,1)$ the walk then ends in $b$ or $c$, contradicting the
corner forcing; in $(2f,1,0)$ the letter after $f^2$ is the $b$, while the terminal vertex at
$(f^2,0)$ carries exactly $\gamma+\alpha$ from the filler and column, leaving a wedge of exactly
$\beta$, whose tile presents $a$ or $c$ --- never $b$. Every case terminates in one of these
contradictions.

\begin{lemma}[Filler identity]\label{lem:filler}
With $J_i$ the $i$-th column junction and $\mathrm{ap}_i$ the $i$-th apex, the triangle
$(J_i,\mathrm{ap}_i,\mathrm{ap}_{i+1})$ has sides exactly $c,b,a$:
$\bigl(\tfrac{3f^2-1}{2f}\bigr)^2+c^2\sin^2\beta=f^4$ and
$\bigl(\tfrac{f^2-1}{2f}\bigr)^2+c^2\sin^2\beta=(f^2-1)^2$, both identities in $f$
(clear denominators and expand; machine-checked in \texttt{lean/W2Core.lean}, axiom-free).
Hence the column and its fillers tile the full inter-level strip and lay a run of $a$-edges one
level down, whose far sides decompose only as single $a$-edges.
\end{lemma}

\begin{lemma}[Offset congruence]\label{lem:offsets}
The terminal column's base apex lies at $\bigl[j(3f^2-1)+\sum_{i}\varepsilon_i\bigr]/(2f)+kf$ with
$\varepsilon_i\in\{-(f^2-1),\,3f^2-1\}$ per level. Modulo $f$ --- for every $f$, of either parity,
since $f^2\equiv0$ --- these are $+1$ and $-1$. Let $p$ be the number of levels taking $+1$ and
$q=j-p$ the rest, so $\sum\varepsilon_i\equiv p-q$ and
\[
\sum\varepsilon_i-j=-2q .
\]
Integrality of a base vertex gives $f\mid 2q$ when $f$ is odd, and $2f\mid 2q$ when $f$ is even;
either way $f\mid q$, and $q\le j<f$ forces $q=0$. Hence $\sum\varepsilon_i=j$, every level is
$(\gamma,\beta)$-oriented, and the apexes lie at consecutive multiples of $f$.
\end{lemma}

\begin{remark}[The evenness hypothesis is removable]\label{rem:parity}
Earlier versions of Lemma~\ref{lem:offsets} worked modulo $2f$, where the two offsets are $+1$ and
$-1$ only when $f$ is even, since that congruence needs $f^2\equiv0\pmod{2f}$; for odd $f$ the
residues are $f+1$ and $f-1$ and the computation does not transfer. That is why
Theorem~\ref{thm:e1family} was previously stated for even $f$ only, leaving the members
$N=26,74,146,\dots$ claimed by Theorem~\ref{thm:basebeta-full} but uncovered.

Both defects disappear at modulus $f$. There the residues are $\pm1$ with no parity hypothesis, and
parametrising by the numbers $p,q$ of $+1$ and $-1$ levels makes $\sum\varepsilon_i-j=-2q$
manifestly even, so the parity check that the odd case appeared to require never has to be made: it
is absorbed into the parametrisation. The forcing is
\textup{\texttt{lean/OffsetForcing.lean}} (\texttt{forcing\_odd}, \texttt{forcing\_even},
\texttt{forcing}), general in $f$ and axiom-clean, superseding the per-member even-$f$ checks of
\texttt{lean/W2Core.lean}.
\end{remark}

\begin{lemma}[Interior multiples]\label{lem:interior}
A base edge of length $f^2-1$ or $f^2$ placed with its left end below $f^2$ contains a multiple
of $f$ in its open interior, which would be a column apex interior to a boundary edge ---
impossible. Hence the letters covering $[0,f^2)$ are exactly $f$ copies of $a$.
\end{lemma}

The corner cases $(1,2)$, $(1,3)$, $(1,4)$ are additionally settled by exhaustive certificates
(the seeded refutations of the side-break configurations; $267$, $1809$ and $1728$ nodes, all
branches dead), and the members $11,26,47$ by the full engine searches recorded in the main
paper. The arithmetic layers of the chain are machine-checked in
\texttt{BAdjacency.lean}, \texttt{Rigidity.lean}, \texttt{W2Core.lean}, \texttt{MidTriangle.lean}
and \texttt{SurplusLattice.lean} (all kernel-only, axiom-clean).

\section{The general base-$\beta$ architecture}\label{sec:compgeneral}

The $e=1$ machinery of Section~\ref{sec:compe1} extends to every coprime pair. Four pillars are
identities in $(e,f)$: the column reach $b\sin\gamma=c\sin\beta$ (law of sines); the filler-mesh
identity (the column's junction fillers are exactly the triangles on consecutive apexes, sides
$c,b,a$; \texttt{lean/GeneralPillars.lean}); the sector inequalities $\gamma-(\alpha+\beta)=\alpha$
and $\cos\gamma=-e/2f<0$; and the interior-multiple bounds. Two master lemmas complete the frame
(\texttt{lean/MasterLemmas.lean}): a base segment of length $a$ between forced junctions
decomposes only as a single $a$ (two applications of $\gcd(e,f)=1$), and the two corner-block
types --- $f$ $a$-feet, or dually $e$ $c$-feet --- have identical footprints $ef^2$ (the relation
$ec=fa$), so in every walk of the trichotomy the two corner structures consume exactly the
non-$b$ letters and the middle is exactly the $e$ $b$-letters. Walks containing an $a$-block die
at the middle's $\alpha+\beta$ corner with $(b,b)$ rays; the walk $(0,e,2e)$ is forced to
$c^e b^e c^e$, its two dual blocks complete, and the resulting pentagon is refuted (certificate at
$(2,3)$: $85$ nodes, all branches dead; the general argument runs through the dual-junction
overhangs). The orientation congruences force $(\gamma,\beta)$ at every level for both parities of
$f$. The assembled exposition of the general recursion is in preparation; every arithmetic
component cited above is kernel-checked, and the certified refutations cover
$(1,3)$, $(1,4)$, $(2,3)$.

\section{The base-$\beta$ branch: no instance at $m=1$}\label{sec:compfull}

This section assembles the complete argument. Throughout, $(e,f)$ is coprime with $1\le e<f$, the
tile is $(a,b,c)=(ef,\,f^2-e^2,\,f^2)$ with angles $(\alpha,\beta,\gamma)$, $\gamma=2\alpha+\beta$,
$3\alpha+2\beta=\pi$, $\alpha/\pi$ irrational; the target $\Delta$ is isosceles with base angles
$\beta$, base $Y=e(3f^2-e^2)$, equal sides $X=f^3$, and $N=3f^2-e^2$.

The argument below is \emph{conditional}, on one geometric statement about the base corners that we
isolate here and do not prove. We state it first, so that nothing downstream is read as
unconditional.

\begin{hypothesis}[Complete corner walls]\label{hyp:walls}
Let $\mathcal T$ be a tiling of the base-$\beta$ target at $m=1$. Then neither base-corner structure
is starved or broken: each is the complete block --- $f$ $a$-feet at the west corner, $e$ $c$-feet at
the east --- whose hypotenuse is a straight segment, of length $f\,b$ resp.\ $e\,b$, partitioned into
whole $b$-edges.
\end{hypothesis}

\begin{theorem}[The base-$\beta$ branch, conditional on Hypothesis~\ref{hyp:walls}]\label{thm:basebeta-full}
Assume Hypothesis~\ref{hyp:walls}. Then $\Delta$ admits no tiling by $N$ congruent copies of the
tile. Consequently no prime $N\equiv 11 \pmod{12}$ is a number of congruent triangles into which a
triangle can be cut, and with the results of the main paper the folklore conjecture would follow:
no prime $N\equiv 3\pmod 4$, $N>3$, is such a number.
\end{theorem}

\begin{remark}[Why the hypothesis is needed, and where it bites]\label{rem:whyhyp}
Hypothesis~\ref{hyp:walls} is exactly the step Remark~\ref{rem:cleanfail} already identifies as
``the single highest-value open step in this branch''; isolating it here makes the dependence
explicit rather than leaving it distributed over the funnels. Two places consume it.

First, the funnels of \S\ref{sub:funnels} dispose of a walk by completing both corner blocks; if a
corner is \emph{starved or broken} the text routes the case to ``the walk arithmetic \dots\ exactly
as in the $e=1$ case'', and the $e=1$ case is Theorem~\ref{thm:e1family}, whose proof defers the
strip-and-column exposition. At $e=1$ this is not a corner case but the whole problem:
Theorem~\ref{thm:e1reduce} forces the base walk to begin and end with $a$-edges, which is
incompatible with the complete-block reading of the base, so \emph{every} hypothetical $e=1$ tiling
falls into the starved-or-broken clause.

Second, Lemma~\ref{lem:pentagon} is stated ``after both dual blocks complete'', and its stubs are
region-boundary segments only because the opposite side is a wall of whole $b$-edges. In a
dissection that is not edge-to-edge --- which the main paper explicitly permits --- a far-side edge
may begin inside a stub and overhang it, so no such segment need be created. The arithmetic of the
lemma is unaffected and is machine-checked general in $(e,f)$
(Lemma~\ref{lem:stubgap}, \texttt{lean/PentagonLemma.lean}); what is conditional is that a stub
arises at all.

The unconditional consequences of this section are therefore the arithmetic ones, together with the
members settled by certified exhaustive search independently of any of it: $N=11,23,26,47,59,66,71,107$.
The member $(1,2)$, i.e.\ $N=11$, is moreover closed unconditionally by
Theorem~\ref{thm:basebeta-e1}, whose argument is boundary-only and needs no corner-block hypothesis.
\end{remark}

\begin{remark}[A corner block can break: the hypothesis needs $m=1$ essentially]\label{rem:blockbreaks}
It is worth recording what the corner blocks are, and that Hypothesis~\ref{hyp:walls} is genuinely a
statement about $m=1$ and not a local fact.

First the identification. At a base corner the complete block with $f$ $a$-feet has footprint $fa=ef^2$
and, by the law of cosines against the base angle $\beta$, its far side has length exactly $f^3=fc$;
so \emph{the block is the tile scaled by $f$}, and it contains $f^2$ tiles. Dually the $c$-feet block
is the tile scaled by $e$, with footprint $ec=ef^2$, far side $e^2f=ea$, and $e^2$ tiles. At $m=1$ the
$f$-block's far side is the whole equal side, so its hypotenuse is the cevian from the apex. The two
counts and the middle add up,
\[
N \;=\; f^2 \;+\; e^2 \;+\; 2\,(f^2-e^2),
\]
which is $3f^2-e^2$, as it must be.

Now the point. The local ingredients of the strip argument --- the corner decomposition
$\{\beta\}$ with flanks $\{a,c\}$, the single-$b$ partner, the parallelogram mate, the kite kills,
the surplus lattice --- are all statements about the tile and about angles. None of them mentions $m$.
If they sufficed to force a complete block they would force one at every $m$. They do not: among the
three genuine tilings of base-$\beta$ targets that we possess, all with tile $(2,3,4)$ and all
verified exactly, the $44$-tiling \texttt{N44B} of $(16,16,22)$ has base walk $a^4c^3a$, and at its
right-hand corner \emph{neither} candidate block is complete --- the $f$-block ($4$ straddling tiles)
nor the $e$-block ($1$ straddler) --- because the last base letter there is a single $a$ where a
block would need the footprint $ef^2=4$ to be cleanly composed. The other two, \texttt{N44C} and the
$99$-tiling of $(24,24,33)$, do have both corners complete.

So a broken corner is not merely permitted in principle; it occurs in a tiling we have exhibited.
Any proof of Hypothesis~\ref{hyp:walls} must therefore use $m=1$ in an essential way --- the natural
candidate being the base $b$-count, which is $e$ at $m=1$ but is $0$ for both $44$-tilings and $7$
for the $99$-tiling --- and no argument phrased purely in the local toolkit above can succeed. We
record this because it is the sharpest constraint we have on how the gap can be closed, and because
it disqualifies the shortest-looking route.
\end{remark}

\subsection{The angle calculus}\label{sub:anglecalc}
Every corner argument below is an instance of one finite computation, which we isolate once.

Since $\alpha/\pi$ is irrational (Sec.~\ref{sec:compscope}), so is $\alpha/\beta$: from
$\alpha/\beta=p/q$ and $3\alpha+2\beta=\pi$ one gets $\beta=q\pi/(3p+2q)$, making $\beta$ and hence
$\alpha$ rational multiples of $\pi$. Consequently an angle has \emph{at most one} representation
$x\alpha+y\beta$ with $x,y\ge0$: two would give $(x-x')\alpha=(y'-y)\beta$. Encoding
\[
\alpha\mapsto(1,0),\qquad \beta\mapsto(0,1),\qquad \gamma=2\alpha+\beta\mapsto(2,1),\qquad
\pi=3\alpha+2\beta\mapsto(3,2),
\]
a vertex figure at an angle $(X,Y)$ with $n_\alpha$, $n_\beta$, $n_\gamma$ tile corners is exactly a
solution of $n_\alpha+2n_\gamma=X$, $n_\beta+n_\gamma=Y$ in nonnegative integers. All of the
following are that system, solved; each is machine-checked in
\textup{\texttt{lean/AngleArithmetic.lean}} (core toolchain, no \texttt{sorry}).

\begin{lemma}[Angle calculus]\label{lem:anglecalc}
For every member $(e,f)$:
\begin{enumerate}
\item \textup{(No right angle.)} $\pi/2$ has no representation: it would force $2x=3$. Hence no
piece of any dissection of a base-$\beta$ region has a right angle. In particular the target cannot
be split along its altitude, and no perpendicular cut is available anywhere in the branch.
\item \textup{(Apex.)} At the apex, $(X,Y)=(3,0)$, so $n_\beta=n_\gamma=0$ and $n_\alpha=3$: the
apex carries exactly three $\alpha$-corners.
\item \textup{(Base corner.)} At $(0,1)$ the only figure is a single $\beta$-corner, with flanks in
$\{a,c\}$.
\item \textup{($\gamma$-trap.)} At a straight angle, $(3,2)$, the only figures are
$\{\alpha,\alpha,\alpha,\beta,\beta\}$ and $\{\gamma,\alpha,\beta\}$; in particular at most one
$\gamma$ meets any interior boundary point.
\item \textup{(Wedge $\alpha+\beta$.)} At $(1,1)$ the only figure is one $\alpha$ and one $\beta$.
\end{enumerate}
\end{lemma}

\begin{remark}[Where $\beta$ sits]\label{rem:betapi3}
$3\alpha>\beta$ iff $\beta<\pi/3$ iff $\cos\beta>\tfrac12$ iff $e(3f^2-e^2)>f^3$. For $e=1$ this is
$f^3-3f^2+1<0$, which holds only at $f=2$; in general it is $f/e<2.879\ldots$, the real root of
$t^3-3t^2+1$. So $(1,2)$ --- the one member whose instance spectrum is completely determined
(Theorem~\ref{thm:realize12}) --- is the unique member of its family with $\beta<\pi/3$, and should
not be treated as typical of $e=1$.
\end{remark}

\subsection{The \texorpdfstring{$\gamma$}{gamma}-grading of directions}\label{sub:gammagrading}
Before the forcing arguments, one structural fact about the whole branch, which does not seem to have
been noticed and which costs two lines.

\begin{proposition}[$\gamma$-grading]\label{prop:gammagrading}
Let $\mathcal T$ be any tiling of any target by a base-$\beta$ tile. Then every edge direction
occurring in $\mathcal T$ is an integer multiple of $\gamma$ modulo $\pi$. Writing a direction as
$x\alpha+y\beta$, the integer is $L=2x-3y$; it is well defined because $3\alpha+2\beta=\pi$.
\end{proposition}
\begin{proof}
From $\gamma=2\alpha+\beta$ and $3\alpha+2\beta=\pi$:
\[
2\gamma=4\alpha+2\beta=4\alpha+(\pi-3\alpha)=\alpha+\pi\equiv\alpha,
\qquad
-3\gamma-\beta=-6\alpha-4\beta=-2(3\alpha+2\beta)=-2\pi\equiv0 ,
\]
so $\alpha\equiv2\gamma$ and $\beta\equiv-3\gamma$ modulo $\pi$. Hence
$\langle\alpha,\beta\rangle=\langle2\gamma,3\gamma\rangle=\langle\gamma\rangle$, as $\gcd(2,3)=1$.

It remains to see that all directions lie in $\theta_0+\langle\alpha,\beta\rangle$ for one reference
$\theta_0$. Read around a tile, the direction turns by the exterior angle, $-(\text{interior})$
modulo $\pi$, so a tile's three directions are $\theta,\ \theta-\alpha,\ \theta-\alpha-\beta$: a
translate of a fixed three-element set $S$. If two tiles share a boundary segment of positive length
they share a direction, so their offsets differ by an element of $S-S\subseteq\langle\alpha,\beta\rangle$.
The tile adjacency graph is connected, the target being connected and the tiles finitely many closed
sets with disjoint interiors covering it; propagating along a chain from a reference tile gives the
claim.
\end{proof}

Two consequences are used below. Round a tile the label shifts are $-1,-2,+3$ at the $\gamma$-,
$\alpha$- and $\beta$-vertices, so a tile's three edge labels are $\{L-2,L,L+1\}$, or
$\{L-1,L,L+2\}$ in the mirrored chirality: \emph{span exactly three}, with the $b$-edge carrying $L$
in either case. And the target's own sides are graded $0$ (base) and $\mp3$ (the equal sides).

\begin{remark}
The point of Proposition~\ref{prop:gammagrading} is that $\alpha/\pi$ is irrational, so the group
generated by $\alpha$ and $\beta$ is dense modulo $2\pi$; one expects no finite structure on
directions at all, and indeed $\langle\gamma\rangle$ \emph{is} dense in $\R/\pi\Z$. What the grading
says is not that the direction set is discrete but that it is \emph{cyclic}: the directions are the
multiples of one angle, so they carry a canonical integer label, and a given tiling uses only a
narrow window of those labels. This is worth isolating: the invariant no-go of the main paper is
proved for functionals on a dense direction set, and does not automatically apply to functionals
graded by $L$. The arithmetic is machine-checked in
\textup{\texttt{lean/LabelCalculus.lean}}, and the grading was tested against every tiling in this
work --- the $44$-tilings, the $99$-tiling, the $28$-tiling on a cevian target and the scalene
four-component $77$-tiling --- with no direction failing.
\end{remark}

\begin{proposition}[The direction group of a branch]\label{prop:dirgroup}
Let $T$ be a triangle whose angles satisfy a single primitive relation, i.e.\ the lattice
$\Lambda=\{(x,y)\in\Z^2:\ x\alpha+y\beta\equiv0 \pmod\pi\}$ is generated by one vector $(p,q)$. Then
in any tiling by congruent copies of $T$ all edge directions lie in one coset of
$\langle\alpha,\beta\rangle\subset\R/\pi\Z$, and
\[
\langle\alpha,\beta\rangle\ \cong\ \Z^2/\langle(p,q)\rangle\ \cong\ \Z\oplus\Z/d,
\qquad d=\gcd(p,q).
\]
Writing $(p,q)=d\,(p',q')$, the free coordinate is the label $L(x,y)=q'x-p'y$, which is surjective
onto $\Z$ with kernel exactly $\langle(p',q')\rangle$.
\end{proposition}
\begin{proof}
The coset statement is the propagation argument of Proposition~\ref{prop:gammagrading}, which used
no property of the angles. For the quotient, $\gcd(p',q')=1$ gives $u,v$ with $up'+vq'=1$; then
$L(vn,-un)=n$, so $L$ is onto, and $L(x,y)=0$ means $q'x=p'y$, whence $p'\mid x$ and $(x,y)$ is a
multiple of $(p',q')$. The remaining generator contributes $\Z/d$.
\end{proof}

\begin{remark}[A torsion dichotomy across the classification]\label{rem:torsion}
Solving for $\Lambda$ in each branch of the Laczkovich classification:
\begin{center}
\begin{tabular}{lccl}
\toprule
branch & relation & $(p,q)$ & direction group\\
\midrule
base-$\beta$      & $3\alpha+2\beta=\pi$   & $(3,2)$ & $\Z$\\
$\gamma=2\alpha$  & $3\alpha+\beta=\pi$    & $(3,1)$ & $\Z$\\
$\gamma=2\pi/3$   & $\alpha+\beta=\pi/3$   & $(3,3)$ & $\Z\oplus\Z/3$\\
$\gamma=\pi/3$    & $\alpha+\beta=2\pi/3$  & $(3,3)$ & $\Z\oplus\Z/3$\\
right angle       & $\alpha+\beta=\pi/2$   & $(2,2)$ & $\Z\oplus\Z/2$\\
\bottomrule
\end{tabular}
\end{center}
The two branches in which $\alpha/\pi$ is irrational --- base-$\beta$ and $\gamma=2\alpha$ --- are
exactly those with $d=1$, that is, with torsion-free direction group. Every special-angle branch
carries torsion. We record the coincidence without asserting a mechanism.

The phenomenon that tiles related by rotation through a single irrational angle generate a cyclic
orientation group is not itself new; it is familiar from substitution tilings, where the pinwheel
tiling's orientations are the multiples of $2\arctan(1/2)$ and are dense in the circle. What is
recorded here is the computation of $\Lambda$ for each branch, the resulting torsion dichotomy, and,
in the base-$\beta$ case, the identification of the free generator as $\gamma$ itself. The arithmetic
is machine-checked in \textup{\texttt{lean/DirectionGroup.lean}}.
\end{remark}

\subsection{The forcing toolkit}
\emph{Corners.} A vertex of angle $\theta$ between two boundary rays admits only the tile-angle
decompositions of $\theta$, and the extreme tiles must present flanks compatible with the rays'
first edges; for $\theta=\beta$ the unique decomposition is $\{\beta\}$ and the flanks are
$\{a,c\}$, one per ray (the two orientations). \emph{Partners.} The corner tile's $b$-edge is a
boundary-anchored chord; both sides partition into whole edges of total $b$, and
$x\,a+y\,b+z\,c=b$ forces $(x,y,z)=(0,1,0)$ (reduce mod $f$; sizes), so the far side is a single
$b$-edge; of the two congruent mates the reflected one repeats the corner angle and overflows
($2\gamma=\pi+\alpha$), leaving the parallelogram partner. \emph{Strips.} Inductively, the $k$-th
chord ($k<f$ resp.\ $k<e$ for the dual) carries $k$ aligned $b$-edges; the surplus lattice
(Sec.~\ref{sec:compgeneral}) forbids $0<|dQ|<f$, so the far side matches $b$-for-$b$; kites die
bottom-up ($\gamma$-repeat at the base, then $2\gamma>\pi$ at each junction), partners are forced,
and each junction wedge is exactly $\beta$ resp.\ the dual angles. All statements are symmetric
under the $a$-feet/$c$-feet duality; footprints agree by $ec=fa$.

\subsection{Runs, columns, fillers, recursion}
Along any covered run of an anchored line both sides partition into whole edges of equal total,
and the surplus lies in $\Lambda=\mathbb Z(f,0,-e)\oplus\mathbb Z(f{-}e,-f,f{-}e)$. A break tile's
$c$-edge admits no whole-edge completion of $c-a=f(f-e)$ (mod $f$ and $\gcd$), so the run closes
minimally as the generalized column: $e$ $c$-edges against $f$ $a$-edges ($ec=fa$). The column's
junction fillers are exactly the triangles on consecutive apexes with sides $c,b,a$ (the filler
identity, an identity in $(e,f)$), so column and fillers tile the inter-level strip completely and
lay an $a$-run one level down; an interior $a$-edge's far side decomposes only as a single $a$
(Master Lemma 1: $\gcd(f^2-e^2,ef)=1$), so the run forces the next level's column. The recursion
reaches the base; a tile vertex on the base is a walk junction (no one-sided $T$-junction), and
the accumulated orientation offsets satisfy, modulo $2f$, the congruence that forces every level
to the $(\gamma,\beta)$ orientation --- for both parities of $f$ --- placing the terminal apexes
at consecutive multiples of $a$. Between consecutive apexes the walk letters decompose only as
single $a$'s (Master Lemma 1 again), so the letters spanning the apex range are forced.

\subsection{The funnels}\label{sub:funnels}
By Master Lemma 2 the two corner structures --- $f$ $a$-feet or $e$ $c$-feet --- have equal
footprints $ef^2$, so in each walk of the trichotomy they consume exactly the non-$b$ letters and
the middle is the $e$ $b$-letters:
$(2f,e,0)\mapsto(a,a)$, $(f,e,e)\mapsto(a,c)$, $(0,e,2e)\mapsto(c,c)$.
(The step just used --- that a stub is a genuine region-boundary segment, so that whole-edge
partition applies --- assumes no tile straddles the ray. That assumption is false in general;
see Remark~\ref{rem:straddle}, and Remark~\ref{rem:whyhyp} for how it is absorbed into
Hypothesis~\ref{hyp:walls}.)
If a corner structure is starved or broken, the recursion of the previous subsection forces the
$a$-prefix regardless (the apex lattice), and the $\beta$-wedge at a column terminal (filler
$\gamma$ + column $\alpha$ leave exactly $\beta$, whose tile presents $a$ or $c$) forbids a
$b$-letter there; the walk arithmetic then kills exactly as in the $e=1$ case.

\begin{lemma}[Terminal wedge]\label{lem:termwedge}
Let a base-reaching column terminate at base vertex $V=(x_0+fa,0)$, $x_0$ its anchor offset. Then
the angle at $V$ decomposes as $\gamma$ (the last filler) $+\ \alpha$ (the last column tile)
$+\ \beta$, and the base edge east of $V$ is an $a$- or a $c$-letter.
\end{lemma}
\begin{proof}
The filler $(J_{f-1},\mathrm{ap}_{f-1},\mathrm{ap}_f)$ has its $a$-edge on the base ending at $V$
with the $a$--$b$ corner there, i.e.\ $\gamma$; the column tile ends at $V$ with its apex angle
$\alpha$. The remaining budget is $\pi-\gamma-\alpha=\beta$, whose unique decomposition is a single
$\beta$-tile; its flanks are $a$ and $c$, one of which lies along the base east of $V$. A
$b$-letter there would have to be that flank, and $b\notin\{a,c\}$.
\end{proof}

\begin{lemma}[Pentagon]\label{lem:pentagon}
In the walk $(0,e,2e)$, after both dual blocks complete, the middle region --- base $b^e$, dual
hypotenuses $b^e$, side segments of length $fb$, apex $3\alpha$, containing $3(f^2-e^2)$ tiles ---
admits no tiling.
\end{lemma}
\begin{proof}
Consider the west base corner $W_0=(ef^2,0)$, of interior angle $\pi-\alpha$ (the dual block's
$\alpha$ sits west of it). Its decompositions are $\{\gamma,\beta\}$ and $\{2\alpha,2\beta\}$; in
each, the tile adjacent to the dual hypotenuse presents along it a flank from $\{a,c\}$ ---
$\beta$-tiles carry $(a,c)$, $\gamma$-tiles $(a,b)$ with the $b$ unavailable (it would repeat the
hypotenuse edge inside the wedge), $\alpha$-tiles $(b,c)$ likewise contribute $c$. The hypotenuse
is partitioned into whole $b$-edges of the dual strips, so a flank of length $a$ or $c$ never
terminates at a hypotenuse junction: each placement leaves a stub of length
$s_1\in\{c-b,\,|a-b|\}=\{e^2,\,|ef-f^2+e^2|\}$ against the next $b$-edge. Such a stub is a region
boundary segment and must be partitioned into whole tile edges. Neither can be: the first lies in
$(0,\min(a,b))$, which contains no element of the numerical semigroup $\langle a,b,c\rangle$, and
the second is a gap of that semigroup by Lemma~\ref{lem:stubgap}. So no partition exists and the
region dies. The certified refutation at $(2,3)$ ($85$ nodes, every branch dead, all kills of run or
wedge type) instantiates the argument; the stub arithmetic above is member-independent.
\end{proof} For the middles:
any walk containing an $a$-block dies at the middle's $a$-side corner (Lemma~\ref{lem:termwedge} and the $(b,b)$ row of the vertex table), of angle $\alpha+\beta$
with both rays presenting $b$ --- the $T_{\mathrm{mid}}$ kill of the main paper. The walk
$(0,e,2e)$ is forced to $c^e b^e c^e$ (the dual break flanks are $\{a,c\}$ and no $a$ exists);
both duals complete and the middle is the pentagon: base $b^e$, two dual hypotenuses $b^e$, two
side segments of length $fb$, apex $3\alpha$. Its west wedge is $\pi-\alpha$; every decomposition
places against the dual hypotenuse a flank from $\{a,c\}$, never matching $b$: each placement
leaves a stub, and the stubs walk in steps of $e^2$ (through-vertex overhangs) or $|a-b|$;
each such stub is a gap of the numerical semigroup $\langle a,b,c\rangle$ --- the first,
$e^2 \bmod b$, because it lies in $(0,\min(a,b))$, and the second, $|a-b|$, by
Lemma~\ref{lem:stubgap} --- so no whole-edge partition exists and the region dies. No iteration
occurs: both stubs are gaps at the first step, for every member. (Certificate at
$(2,3)$: $85$ nodes, all dead, kills exclusively of run/wedge type.)

\begin{lemma}[The second stub is a gap]\label{lem:stubgap}
Let $\gcd(e,f)=1$ with $1\le e<f$, and let $(a,b,c)=(ef,\,f^2-e^2,\,f^2)$. Then $|a-b|$ is a gap of
the numerical semigroup $\langle a,b,c\rangle$: it is nonzero and admits no representation
$xa+yb+zc$ with $x,y,z\ge0$.
\end{lemma}
\begin{proof}
First $\gcd(a,b)=1$: a common divisor of $e$ and $b$ divides $b+e^2=f^2$, hence divides
$\gcd(e,f^2)=1$; a common divisor of $f$ and $b$ divides $f^2-b=e^2$, hence divides
$\gcd(f,e^2)=1$; so $\gcd(ef,b)=1$. Next both generators exceed $1$: $a=ef\ge2$, and
$b=(f-e)(f+e)\ge3$ since $f-e\ge1$ and $f+e\ge3$. Hence $a\ne b$ --- equality with $\gcd(a,b)=1$
would force $a=b=1$ --- so $t:=|a-b|>0$.

Now $t<c$: if $a>b$ then $t=ef+e^2-f^2<f^2$ because $ef<f^2$ and $e^2<f^2$; if $b>a$ then
$t=f^2-e^2-ef<f^2$. So $z=0$ in any representation. Also $t<\max(a,b)$ trivially, so the larger of
$a,b$ has coefficient $0$ as well. What remains is $t=k\cdot\min(a,b)$ for some $k\ge1$, i.e.\
$\min(a,b)$ divides $|a-b|$ and therefore divides $\max(a,b)$ too; with $\gcd(a,b)=1$ this forces
$\min(a,b)=1$, contradicting $a\ge2$, $b\ge3$.
\end{proof}

\begin{remark}
The lemma is stronger than the pentagon argument needs and removes its only iterative step. The
earlier treatment sent the stub sequence walking by $e^2$ modulo $b$ until it entered
$(0,\min(a,b))$, a gap whose reachability had to be argued geometrically; the statement above shows
the second stub is \emph{already} a gap, for every member, so no walk occurs. This is not merely a
shorter route to the same place: an exhaustive check over every member with $f\le200$
($12\,231$ pairs, \texttt{rust/stubgap}) finds that \emph{both} initial stubs are gaps at the first
step in every single case, so the walk the earlier argument invoked never executes even one step.
The mechanism can be deleted from the lemma rather than justified. It is formalized in
\textup{\texttt{lean/PentagonLemma.lean}} (\texttt{pentagon\_lemma}, general in $(e,f)$, standard axioms only),
together with $\gcd(a,b)=1$, and cross-checked by exhaustive semigroup search over all coprime pairs
with $f\le400$ (\texttt{rust/stubgap}). Note that $(0,\min(a,b))$ and $\{|a-b|\}$ are genuinely
different gaps: at $(1,3)$ the tile is $(3,8,9)$, $\min(a,b)=3$, and $|a-b|=5$ lies outside
$(0,3)$ while still being a gap of $\langle3,8,9\rangle$.
\end{remark}

\begin{lemma}[Two $c$-edges are forced; $p\le f-2$]\label{lem:ple}
At $e=1$, $m=1$, the first and last edges of each equal side are $c$-edges: the first is the
base-corner tile's side flank (its base flank is an $a$-edge by Theorem~\ref{thm:e1reduce}), and the
last belongs to the apex figure $\{3\alpha\}$, whose side edge is $b$ or $c$, with $b$ excluded by
Lemma~\ref{lem:sidenob}. Hence $f-p\ge2$, i.e.\ $p\le f-2$. In particular $p=0$ for $f=2$: the
member $(1,2)$ satisfies Hypothesis~\ref{hyp:walls} outright, by boundary analysis alone, and the
closure of $N=11$ in Theorem~\ref{thm:basebeta-e1} is recovered through
Remark~\ref{rem:sidenoa}. For $f=3$ the only surviving value is $p=1$, whose side word is forced to
$c\,a\,a\,a\,c$; by the run rigidity its chirality word is $(-4)^j(-2)^{3-j}$ for some
$0\le j\le3$, so the entire boundary configuration space at $(1,3)$ consists of four explicit
placements.
\end{lemma}

\begin{lemma}[The equal sides carry no $b$-edge]\label{lem:sidenob}
Let $m=1$ and let $(e,f)$ be any member. Then every side walk $P'a+Q'b+R'c=f^3$ obeying the
$\gamma$-trap $R'\ge1$ has $Q'=0$.
\end{lemma}
\begin{proof}
Reducing the walk modulo $f$ kills $P'a$ and $R'c$ and leaves $f\mid Q'e^2$; as $\gcd(e,f)=1$,
$f\mid Q'$. Write $Q'=qf$ and suppose $q\ge1$. Dividing the walk by $f$,
\[
P'e+q\,b+R'f=f^2,\qquad b=f^2-e^2 ,
\]
so in particular $q\,b+R'f\le f^2$. Reducing \emph{this} modulo $e$ gives
$f(qf+R')\equiv f^2$, and $f$ is invertible mod $e$, so $R'\equiv f(1-q)\pmod e$.

\emph{Case $q=1$.} Then $R'\equiv0\pmod e$, and $R'\ge1$ forces $R'\ge e$, so $R'f\ge ef$. But
$q\,b+R'f\le f^2$ gives $R'f\le f^2-b=e^2$, whence $ef\le e^2$ and $f\le e$, contradicting $e<f$.

\emph{Case $q\ge2$.} From $q\,b+R'f\le f^2$ with $q\ge2$, $R'\ge1$ we get $2(f^2-e^2)+f\le f^2$,
i.e.\ $f^2+f\le 2e^2$; hence $f<e\sqrt2<2e$, so $u:=f-e$ satisfies $1\le u\le e-1$. Put $t:=q-1\ge1$.
Since $f\equiv u\pmod e$, the congruence above reads $R'\equiv-ut\pmod e$. Write $ut=v+we$ with
$0\le v\le e-1$ and $w\ge0$. Substituting $f=e+u$ and $b=u(2e+u)$, the constraint is
\[
(t+1)\,u\,(2e+u)+R'(e+u)\;\le\;(e+u)^2 .
\]
If $v=0$ then $e\mid R'$ and $R'\ge e$, so the display gives $(t+1)u(2e+u)\le(e+u)^2-e(e+u)=u(e+u)$,
i.e.\ $(t+1)(2e+u)\le e+u$, impossible as $t+1\ge2$. If $v\ge1$ then $R'\ge e-v$, so
$(t+1)u(2e+u)\le(e+u)(u+v)$. But $(t+1)u=ut+u=v+we+u$, so
\[
(t+1)u(2e+u)-(e+u)(u+v)=(v+u)\bigl[(2e+u)-(e+u)\bigr]+we(2e+u)=e\bigl[(v+u)+w(2e+u)\bigr]>0,
\]
a contradiction. Hence $q=0$, i.e.\ $Q'=0$.
\end{proof}

\begin{remark}
The special case $e^2<f$ --- which already covers every $e=1$ member --- has the one-line proof
``$Q'\ne0\Rightarrow Q'\ge f\Rightarrow Q'b+R'c\ge f^3-fe^2+f^2>f^3$'', and it is that case which is
machine-checked in \texttt{lean/SideNoB.lean} as \texttt{side\_no\_b}, with
\texttt{side\_no\_b\_e\_one} the $e=1$ specialisation. The unconditional statement is
\texttt{side\_no\_b\_uncond} in the same file, formalised by the identity $(\star)$
$t\,b+R'f+P'e=e^2$ used in the proof above; both depend only on the standard axioms. Cross-checked
by exhaustive search over every member with $f\le200$.
\end{remark}

\begin{remark}[At $e=1$, Hypothesis~\ref{hyp:walls} is exactly ``the sides carry no $a$-edge'']\label{rem:sidenoa}
Lemma~\ref{lem:sidenob} is one of the few statements in this development genuinely special to $m=1$
--- at $m=2$ it is false, by the identity $e(ef)+f(f^2-e^2)+f\cdot f^2=2f^3$ --- and by
Remark~\ref{rem:blockbreaks} any proof of Hypothesis~\ref{hyp:walls} must use $m=1$ essentially.
It is therefore worth following where it leads at $e=1$, where it reduces the hypothesis to a single
integer.

With $Q'=0$ the side equation $P'f+R'f^2=f^3$ gives $P'+R'f=f^2$, so $f\mid P'$; writing $P'=fp$,
\[
(P',R') = (fp,\; f-p), \qquad 0\le p\le f-1,
\]
the $\gamma$-trap being $R'\ge1$. The complete west block is the tile scaled by $f$, so its far side
--- which at $m=1$ is the whole equal side --- is subdivided into exactly $f$ $c$-edges, i.e.\ $p=0$.
Hence, at $e=1$,
\[
\textup{Hypothesis~\ref{hyp:walls} at the west corner} \iff p=0 \iff
\textup{the equal side carries no } a\textup{-edge}.
\]
Two corner facts pin the ends of that walk. At the base corner the unique figure is a single
$\beta$-tile with flanks $\{a,c\}$, and Theorem~\ref{thm:e1reduce} puts an $a$-edge at the base end,
so the side's first edge is a $c$. At the apex the figure is exactly three $\alpha$-corners
(Lemma~\ref{lem:anglecalc}(2)), and the two edges at a tile's $\alpha$-vertex are $b$ and $c$, so
with $Q'=0$ the side's last edge is a $c$ as well. Every equal side therefore reads
$c\,(\cdots)\,c$ with $fp$ interior $a$-edges.

So the open step at $e=1$ is the single statement ``$p=0$'', a direct analogue of
Lemma~\ref{lem:sidenob} with $a$ in place of $b$, and it is not arithmetic: the walk equation admits
every $p\le f-1$. Finally, note that $p=0$ closes the subfamily outright, without any further
funnel analysis: complete blocks put $f$ $a$-feet at one base corner and $e=1$ $c$-foot at the
other, so the base reads $a^f\,b\,c$ and ends in a $c$-edge, whereas Theorem~\ref{thm:e1reduce}
forces $a$-edges at both ends. The contradiction is immediate. Hypothesis~\ref{hyp:walls} is thus not
a step towards the $e=1$ case; at $e=1$ it \emph{is} the case.
\end{remark}

\begin{lemma}[First-run orientation]\label{lem:firstrun}
In any tiling of the $(1,f)$ target at $m=1$, every tile of the first $a$-run of an equal side is
mirrored: its $\gamma$ sits at the upper junction and its label is $-2$. In particular at $(1,3)$,
where Lemma~\ref{lem:ple} forces the side word $c\,a\,a\,a\,c$, the boundary configuration is
\emph{unique}: both $c$-tiles direct with label $-1$, all three $a$-tiles mirrored with label $-2$.
\end{lemma}
\begin{proof}
The corner tile's $b$-edge is a chord with both endpoints on the target boundary: its upper end is
the first side junction $J$, its lower end the base point $(f,0)$, and the ray beyond either end
leaves the target (past $J$ it crosses the side, the two directions being non-parallel; past
$(f,0)$ it descends below the base). So no covering edge overshoots, the far side partitions into
whole edges of total $b$, and by the partner relation $xa+yb+zc=b\Rightarrow(x,y,z)=(0,1,0)$
(machine-checked, general in $(e,f)$: \texttt{PentagonLemma.partner\_unique}) the far side is a
single $b$-edge: a partner tile $P$ shares the whole edge. The ends of a $b$-edge host $\alpha$ and
$\gamma$, so $P$ contributes $\alpha$ or $\gamma$ at $J$; the corner tile contributes $\alpha$
there. If the first $a$-tile were direct it would contribute $\gamma$ at $J$, and the straight-angle
figure at $J$ would contain one $\alpha$, one $\gamma$, and a further corner from
$\{\alpha,\gamma\}$; neither admissible figure does
(machine-checked: \texttt{OrderForcing.first\_run\_kill}). So the first $a$-tile is mirrored, and
the run rigidity (\texttt{gamma\_far\_absorbing}: mirrored is absorbing along the run) makes the
whole run mirrored.
\end{proof}

\begin{lemma}[Top junction; $p\le f-3$ for $f\ge4$]\label{lem:topjunction}
Let an $a$-run end at the last side junction $J$ (final $c$-block of length $1$) with the run's last
tile mirrored. The junction figure is $\{\gamma,\alpha,\beta\}$, its interior $\alpha$-filler has
its edges along the two junction rays, and the edge along the top tile's $a$-ray is $c=f^2$
(mirrored filler) or $b=f^2-1$ (direct filler), with endpoint $J+(f^2,0)$ resp.\ $J+(f^2-1,0)$.
Testing these against the right side of the target reduces, exactly, to the signs of
\[
f^3-3f^2+1 \quad\text{(mirrored filler outside} \iff >0 \iff f\ge3\text{)},\qquad
f^3-3f^2-f+1 \quad\text{(direct filler outside} \iff >0 \iff f\ge4\text{)}.
\]
Hence for $f\ge4$ no such junction exists; for $f=3$ the filler is forced direct with label $0$.
Consequently: $p=f-2$ gives $c$-count $2$, forcing the single-run word $c\,a^{f(f-2)}\,c$, which is
all-mirrored by Lemma~\ref{lem:firstrun} and ends at such a junction — dead for $f\ge4$. So
\[
p\ \le\ f-3 \qquad (f\ge4),
\]
and at $(1,4)$ in particular $p\le1$. The first polynomial is the thin/thick discriminant of
Remark~\ref{rem:betapi3}: the member $(1,2)$ is on its other side, which is again why $f=2$ is the
degenerate case. The sign reductions are verified exactly (rational cross products over
$\Q(\sqrt D)$) and the two positivity facts are machine-checked
(\texttt{OrderForcing.mirrored\_filler\_outside}, \texttt{direct\_filler\_outside}), as are the
single-run combinatorics and the count bound (\texttt{two\_c\_single\_run},
\texttt{p\_le\_f\_sub\_three}).
\end{lemma}

\begin{lemma}[c-chord dichotomy]\label{lem:cchord}
Let a chord of length $c$ be blocked at both ends, so that its far side partitions into whole tile
edges of total $c$. Then $xa+yb+zc=c$ forces $(x,y,z)=(0,0,1)$, or, exactly when $e=1$,
$(x,y,z)=(f,0,0)$: the far side is a single $c$-edge or (at $e=1$ only) a run of $f$ $a$-edges.
Machine-checked general in $(e,f)$ (\texttt{CChord.c\_chord\_dichotomy}); the modular step is the
identity $(y-1)f^2=e(ye-xf)$.
\end{lemma}

\begin{remark}[Tile-interior blocking]\label{rem:tib}
Two ends of a chord can be blocked by the target boundary; they can also be blocked by the
\emph{interior of an already-forced tile}, since a covering edge extended past the chord's end would
enter it. Forced tiles thereby create new fully-blocked chords, and the forcing propagates.
\end{remark}

\begin{theorem}[Hypothesis~\ref{hyp:walls} holds at $(1,3)$]\label{thm:walls13}
No tiling of the $(1,3)$ target at $m=1$ carries an $a$-edge on an equal side. Hence, by
Remark~\ref{rem:sidenoa}, no tiling of the $(1,3)$ target at $m=1$ exists at all.
\end{theorem}

\begin{proof}
Suppose an equal side carries an $a$-edge. By Lemma~\ref{lem:ple}, $p=1$ and the side word is
$c\,a\,a\,a\,c$; by Lemma~\ref{lem:firstrun} the boundary configuration is unique, and the corner
tile $T_1$, its $b$-partner $P$, and the three mirrored $a$-tiles are all pinned. Write the base
corner at the origin, $J$ for the first side junction, and $Q=J+(f,0)$ for the end of $P$'s
horizontal $a$-edge, so that $P$'s $c$-edge runs from $(f,0)$ to $Q$ parallel to the side. That
edge is blocked at $(f,0)$ by the base and at $Q$ by the interior of the second side tile
(Remark~\ref{rem:tib}; the directions differ, so the crossing is transversal), and
Lemma~\ref{lem:cchord} applies: its far side is a single $c$-edge or a run of three $a$-edges.

The vertex figure at $(f,0)$ already contains $T_1$'s $\gamma$ and $P$'s $\alpha$, leaving exactly
one $\beta$ (the representation of $\pi-\gamma-\alpha$ is unique). The tile carrying it presents,
from its $\beta$-corner with flanks $\{a,c\}$, the base's second edge.

\emph{Single-$c$ branch.} The $c$-tile $C$ has $\beta$ at $(f,0)$ and its $a$-edge on the base:
the base word is $a\,a\,b\,c\,a$, so the third base edge is the $b$-edge $[2f,\,2f+b]$. Label
consistency forces $C$ direct, and $C=T_1+(f,0)$: the corner block grows by one cell. Now $C$'s
$b$-edge is blocked at $(2f,0)$ by the base and at $Q$ by the same tile interior, so its far side
is a single $b$-edge (Lemma~\ref{lem:sidenob}'s companion fact, \texttt{partner\_unique}), carried
by a tile $Z$ whose corners at the two ends are $\alpha$ and $\gamma$ in one of two orientations.
The reflected orientation puts $\gamma$ at $(2f,0)$, where $C$'s $\gamma$ already sits: the
straight-angle figure admits at most one $\gamma$. So $Z=P+(f,0)$, with $\alpha$ at $(2f,0)$ and
$\gamma$ at $Q$, and the figure at $Q$ closes as $\{\beta,\alpha,\gamma\}$. At $(2f,0)$ the figure
now contains $C$'s $\gamma$ and $Z$'s $\alpha$, leaving exactly one $\beta$, whose tile presents
the base edge from $(2f,0)$ out of flanks $\{a,c\}$. But that edge is the $b$-edge of the forced
base word. Contradiction.

\emph{Three-$a$ branch.} The first $a$-tile $A_1$ has its corner at $(f,0)$ among the $a$-edge's
end angles $\{\beta,\gamma\}$; the single open slot is $\beta$, so $A_1$ has $\beta$ at $(f,0)$
and its other flank, the $c$-edge, on the base: the base word is $a\,c\,b\,a\,a$, with the $b$-edge
at $[f+c,\,f+c+b]=[12,20]$. Label consistency forces $A_1$ mirrored. $A_1$'s $b$-edge runs from the
first trisection point of $P$'s $c$-edge down to $(f+c,0)=(12,0)$; it is blocked at $(12,0)$ by the
base and at the trisection point by $P$'s interior, so its far side is again a single $b$-edge,
carried by $Z'$. The reflected orientation repeats $\gamma$ at the trisection point, where $A_1$'s
$\gamma$ already sits — impossible. So $Z'$ is the parallelogram mate with $\gamma$ at $(12,0)$.
The figure at $(12,0)$ then contains $A_1$'s $\alpha$ and $Z'$'s $\gamma$, leaving exactly one
$\beta$ — but the base edge from $(12,0)$ is the $b$-edge of the forced word, whose ends host
$\alpha$ or $\gamma$. Contradiction.

Both branches of Lemma~\ref{lem:cchord} are impossible, so no such tiling exists; $p=0$.
\end{proof}

\begin{remark}\label{rem:walls13}
Three remarks. First, the terminal mechanism is the same in both branches: the cascade advances one
cell from the corner and the base word, fixed at the fork, collides with it one base point later.
Second, the proof is search-free; it re-derives the exclusion of the $(1,3)$ instance ($N=26$)
established by the certified exhaustion ($2\,025$ nodes), and the two agree. Third, the arithmetic
components are machine-checked (\texttt{partner\_unique}, \texttt{c\_chord\_dichotomy},
\texttt{first\_run\_kill}, \texttt{straight\_junction\_cases}, the slot classifications
\texttt{beta\_slot\_unique}, \texttt{alpha\_slot\_unique}, \texttt{wedge\_gamma\_cases}, and the
polynomial signs of Lemma~\ref{lem:topjunction}); the geometric steps are exact coordinate
computations over $\Q(\sqrt{35})$. Hypothesis~\ref{hyp:walls} is thereby proved for the members
$(1,2)$ (Lemma~\ref{lem:ple}) and $(1,3)$; it remains open for $f\ge4$, where $p\le f-3$, and for
$e\ge2$.
\end{remark}

\begin{remark}[The cascade at $(1,4)$, and the general mechanism]\label{rem:walls14}
Running the same cascade at $(1,4)$ ($p\le1$ by Lemma~\ref{lem:topjunction}) kills five of the
\emph{six} admissible base words (an earlier count of five missed $a\,c\,b\,a\,a\,a$, which dies
at once: its $b$-edge sits at $[f+c,\,f+c+b]$, exactly the four-$a$ branch's first $\beta$-slot). For every side word with initial $c$-block $1$, the corner partner and the
$c$-chord dichotomy force, at each base point $(kf,0)$ reached, a $\beta$-slot whose tile presents
the next base edge from flanks $\{a,c\}$: the words with $b$ in position $3$ die at $(2f,0)$; the
word $a\,a\,a\,b\,c\,a$ dies at $(3f,0)$; the word $a\,a\,c\,b\,a\,a$ dies at $(f+c,0)=(24,0)$
through the $a$-run partner, exactly as at $(1,3)$. Two configurations survive the present tools,
and the obstruction is located precisely in each: for the base word $a\,c\,a\,b\,a\,a$ the deviating
$a$-run's new $c$-chord $[(20,0),\,(20,0)+cu]$ has its upper endpoint beyond the second side tile's
$c$-edge ($25.875>21.875$ exactly), so tile-interior blocking fails and the dichotomy does not
apply; and for the side word $c\,c\,a^4\,c$ the first run does not abut the corner tile, so
Lemma~\ref{lem:firstrun} gives no chirality forcing. Hypothesis~\ref{hyp:walls} at $(1,4)$ is
thereby reduced to these two configurations; the first is now eliminated, and the second is
narrowed, as follows. \emph{The double-$c$ side word dies.} With the second $c$-tile forced direct
and the run all-direct, the $\beta$-slot tile $B$ at the run's first junction has exactly two
label-consistent forms, and both are impossible. The junction ray toward the second $c$-tile's
$b$-edge extends, exactly, to the base point $(2f,0)$: the points $2c\,u$, $Q$ and $(2f,0)$ are
collinear with $|2c\,u\to(2f,0)|=2b$. If $B$ is direct, its $a$-edge covers the first $4$ units of
that line's far side, and the remaining stretch to $Q$ has length $11\notin\langle4,15,16\rangle$;
some covering edge must therefore cross $Q$, a through-edge forces the west $\alpha$-slot tile
$X=T_1+(f,0)$ (the mirrored $X$ dies by base overshoot), so the whole segment $[Q,(2f,0)]$ is an
edge and the far side must partition $26\notin\langle4,15,16\rangle$ with no crossing at the
$\gamma$-end. If $B$ is mirrored, its $c$-edge itself crosses $Q$ by $e^2=1$, forcing $X$ the same
way, and the remaining stretch is $14\notin\langle4,15,16\rangle$ with no crossing at the
$\alpha$-end. Both die; the side word $c\,c\,a^4\,c$ admits no completion for any base word. The
crossing exclusions are instances of the general principle that a through-edge occupies the whole
half-turn on its side (\texttt{through\_edge\_exclusive}); the three gaps are machine-checked
(\texttt{gaps\_4\_15\_16}). Each remaining configuration is narrowed as follows. In the
double-$c$ configuration $c\,c\,a^4\,c$, the second $c$-tile is forced direct: mirrored would place
its $\alpha$ at the first junction beside the corner tile's $\alpha$ and the partner's $\gamma$,
and $\pi-2\alpha-\gamma$ has no representation (\texttt{first\_run\_kill}); moreover the final
$c$-block has length $1$, so by Lemma~\ref{lem:topjunction} the run must be entirely direct, and the
configuration now hinges on the single unforced $\beta$ at the run's first junction. In the base
word $a\,c\,a\,b\,a\,a$, the $\beta$-slot at $(20,0)$ forces the tile $T_1+(20,0)$, whose new
$c$-chord is covered on its far side, in the mirrored-filler sub-case, by the mate's $a$-edge, two
further $a$-edges, and the filler's own $a$-edge, whose endpoint is exactly the escape point
$E=Q+(c,0)=(20,0)+cu$; the direct-filler sub-case leaves $E$ unblocked and remains open.
\end{remark}

\begin{proposition}[The double-$c$ kill at any initial block length]\label{prop:doublec}
Let $e=1$, $m=1$, $f\ge3$, and let an equal side carry an $a$-edge and have initial $c$-block of
length $j\ge2$. Write $u$ for the unit vector along that side and place the base corner at the
origin. Then:
\begin{enumerate}
\item[\textup{(i)}] $j\le f-1$, and the chord from $j\,c\,u$ to the base point $(jf,0)$ has length
exactly $jb$, divided into $j$ segments of length $b$ by $R_i=i\,c\,u+(j-i)(f,0)$ (so $R_1$ is the
point called $Q$ in the $(1,3)$ argument; the $Q_j$ of Lemma~\ref{lem:wallclimb} are a different
family, spaced $c$ apart along the interior wall);
\item[\textup{(ii)}] that chord admits exactly one partition into tile edges, namely $j$ $b$-edges;
\item[\textup{(iii)}] the $\beta$-slot tile at $j\,c\,u$ presents, from its $\beta$-corner, a flank
of $\{a,c\}$, and each choice leaves a run outside $\langle a,b,c\rangle$;
\item[\textup{(iv)}] the $\alpha$-slot tile at $R_{j-1}$ cannot lay a $c$-edge along the chord.
\end{enumerate}
\end{proposition}

\begin{proof}
\emph{(i)} The side begins and ends with a $c$-edge (Theorem~\ref{thm:e1reduce}(i)), so an initial
block of length $\ge2$ leaves a further $c$ at the far end and $n_c\ge3$; with $n_c=f-k$ and $k\ge1$
(the side carries an $a$-edge) this gives $j\le n_c=f-k\le f-1$. For the chord, $\cos\beta$ is
$(3f^2-1)/(2f^3)$ by the law of cosines on $(f,f^2-1,f^2)$, which is the isosceles ratio of the
target, so with $c\,u=(f^2\cos\beta,f^2\sin\beta)$,
\[
|j\,c\,u-(jf,0)|^2=j^2\bigl(f^2-1\bigr)^2=(jb)^2 ,
\]
and $R_i$ is the affine combination dividing the segment in ratio $i:(j-i)$, so consecutive $R_i$
are $b$ apart (\texttt{chord\_jb}, \texttt{chord\_jb\_segment}).

\emph{(ii)} Reducing $x a+y b+z c=jb$ modulo $f$ forces $y\equiv j$, while $yb\le jb$ forces $y\le j$;
as $0\le j\le f-1$ this gives $y=j$, hence $xa+zc=0$ and $x=z=0$ (\texttt{partition\_jb\_gen}, which
is general in $(e,f)$; \texttt{partition\_30} is its case $f=4$, $j=2$).

\emph{(iii)} The $\beta$-slot exists: at the side junction $j\,c\,u$ where the block ends,
Lemma~\ref{lem:wallclimb} gives $C_j$'s $\alpha$ and $M_j$'s $\gamma$, which force the next side
tile's corner to $\beta$, and that lemma is unconditional --- it uses no base-letter hypothesis and
no line closure, so it applies at every block length. Now $\beta$ is opposite $b$, so its incident
edges are $a$ and $c$: the tile in the $\beta$-slot lays an $a$-edge (direct) or a $c$-edge
(mirrored), never a $b$-edge. Both runs left on
the chord are gaps, by the uniform statement that $jb-t\notin\langle a,b,c\rangle$ whenever
$1\le j\le f-1$ and $f\mid t$, $t>0$ (\texttt{gap\_jb\_minus\_multiple}): indeed $a=f$ and $c=f^2$
are both multiples of $f$. Concretely, if the tile is direct its $a$-edge does not reach $R_{j-1}$
and the remainder $b-a$ is a gap, so no edge ends at $R_{j-1}$ and some edge crosses it; if it is
mirrored its $c$-edge overshoots $R_{j-1}$ by $c-b=e^2=1$ and crosses it outright. Either way a
through-edge occupies the whole half-turn on its side
(\texttt{through\_edge\_exclusive}), the segment to $(jf,0)$ is a single edge, and the far side must
partition $jb-a$ resp.\ $jb-c$ --- gaps.

\emph{(iv)} The last segment $[R_{j-1},(jf,0)]$ of the chord has length exactly $b$ by \emph{(i)},
and $c-b=e^2>0$, so a $c$-edge laid from $R_{j-1}$ along the chord passes $(jf,0)$ by $e^2$. The
chord descends: $j\,c\,u$ has ordinate $jf^2\sin\beta>0$ and $(jf,0)$ lies on the base, so the
continuation beyond $(jf,0)$ has negative ordinate and leaves the target. This is the directional
overshoot test of Remark~\ref{rem:whyf3}: a $c$ laid downward onto a base point exceeds the
boundary and dies, whereas the same edge laid upward from the base would merely pass an interior
point. Hence the $\alpha$-slot flank along the chord is a whole $b$-edge, and the reflected placement
is excluded. Nothing here depends on $f$, nor on $e$ beyond $e^2>0$.
\end{proof}

\begin{remark}\label{rem:doublecscope}
Proposition~\ref{prop:doublec} is the $f=4$ argument of Remark~\ref{rem:walls14} with every constant
replaced by its general form: the three runs $11,26,14$ against $\langle4,15,16\rangle$ are
$b-a$, $2b-a$, $2b-c$, i.e.\ the cases $(j,t)=(1,a),(2,a),(2,c)$, and the $2b$ chord is $j=2$. The
bound $j\le f-1$ of \emph{(i)} is sharp for \emph{(iii)}: at $j=f$ both $jb-a=f(f^2-2)$ and
$jb-c=f(f^2-f-1)$ are multiples of $f$ and lie in the semigroup.

Part \emph{(iv)} supplies what the $f=4$ argument records as ``the mirrored $X$ dies by base
overshoot'': there $X=T_1+(f,0)$ is the triangle $(f,0),(2f,0),R_1$, and the exclusion is the case
$j=2$. The general reason is that the last segment of the chord has length exactly $b$ while the
competing edge has length $c=b+e^2$, and the chord's far end is a base point, so the excess leaves
the target. Consequently escape route 2 of Conjecture~\ref{conj:advance} is closed at every initial
block length, for every $f\ge3$. Route 1 --- a deviating $a$-run whose new $c$-chord ends beyond the
forced region --- is what now stands between this and Hypothesis~\ref{hyp:walls} at $e=1$, and the
next remark says exactly how it behaves in $f$.
\end{remark}

\begin{remark}[Route 1 escapes by exactly $a$, uniformly in $f$]\label{rem:route1uniform}
The escape recorded at $(1,4)$ as ``$25.875>21.875$ exactly'' is not a numerical accident of $f=4$,
and it does not improve with $f$. Write $A=c\,u$ for the upper end of the side's first $c$-edge and
$B=(c+a)u$ for the upper end of the run's first $a$-edge. The second side tile carries its $a$-edge
$AB$ on the side and, being mirrored (Lemma~\ref{lem:firstrun}), presents $\beta$ at $A$, so its
$c$-edge runs from $A$; its third vertex is
\[
V=c\,u+(c,0),\qquad |V-A|=c,\quad |V-B|=b ,
\]
so that $c$-edge is \emph{horizontal}, at height $c\sin\beta$. The deviating $a$-run's new $c$-chord
starts at the base junction $(f+c,0)$ and has upper endpoint $E=(f+c,0)+c\,u$, at the same height.
Hence
\[
x(E)-x(V)=f
\]
identically: the chord's endpoint overshoots the blocking edge by exactly one $a$-length, for every
$f$. At $f=4$ this is $25.875-21.875=4$.

Two consequences. The escape condition $x(E)>x(V)$ reads $f>0$, so it holds in every member: there
is no analogue here of the sign conditions $f^3-3f^2+1$ and $f^3-3f^2-f+1$ by which
Lemma~\ref{lem:topjunction} disposes of its case, and no threshold in $f$ beyond which route 1
closes itself. Any proof must therefore recover the blocking by a different mechanism rather than by
showing the overshoot absent.

Second, the configuration on that horizontal line is completely rigid, and says what the mechanism
has to be about. Writing $R_1=c\,u+(f,0)$ for the point called $Q$ in the $(1,3)$ argument (the
midpoint of the $2b$ chord of Proposition~\ref{prop:doublec}), all four relevant points lie at the
common height $c\sin\beta$, in the order
\[
A \ \overset{a}{\text{---}}\ R_1 \ \overset{c-a}{\text{---}}\ V \ \overset{a}{\text{---}}\ E ,
\qquad |AV|=|R_1E|=c ,
\]
so the line carries \emph{two overlapping $c$-segments offset by exactly $a$}: the second side
tile's $c$-edge $[A,V]$, which is a genuine tile edge, and $[R_1,E]$, whose right end is the chord's
unblocked endpoint. In particular the gap $[V,E]$ between the blocking edge's end and the chord's
endpoint has length exactly $a$ --- the length of an $a$-edge. What route 1 needs is therefore a
statement about what covers $[V,E]$: if that segment is forced to be a whole $a$-edge then $E$ is a
junction and the vertex figure there is constrained, which is the substitute for the tile-interior
blocking that fails. We do not prove that here, and record only that the geometry leaves no other
target: every length in the display is exact and independent of $f$.
\end{remark}

\begin{remark}[The blocking hypothesis cannot be weakened to ``both ends are vertices'']\label{rem:noedgetoedge}
It is worth recording what the step left open in Remark~\ref{rem:route1uniform} may \emph{not} appeal
to. In Lemma~\ref{lem:cchord} the hypothesis that the chord is blocked at both ends is what converts
it into a union of whole tile edges; one cannot replace it by the weaker hypothesis that both ends
are vertices of the tiling, because the tilings in this family are not edge-to-edge. The certified
$99$-tiling of member $(1,2)$ settles this concretely. Its $99$ tiles span $216$ distinct edges, and
$52$ of those carry a vertex of the tiling in their relative interior --- each exactly one. Moreover
its $89$ vertices contain $19$ pairs at distance exactly $a$, $23$ at distance exactly $b$ and $45$
at distance exactly $c$ that span no tile edge at all. (These counts are read off the coordinates of
the machine-checked certificate; in that embedding every $x$ is rational and every $y$ a rational
multiple of $\sqrt{15}$, so the chart $(x,y/\sqrt{15})$ is an affine rescaling of the plane and
decides collinearity and betweenness faithfully.)

Hence knowing that a segment has both endpoints at tiling vertices and has the length of a tile side
carries \emph{no} information about how it meets the tiling. Since $V$ and $E$ are vertices at
distance exactly $a$, this is precisely the shape of the hypothesis one is tempted to use, and it is
not available. Whatever forces $[V,E]$ to be a whole $a$-edge must consume the cascade's tile data
--- which tiles have been forced and where --- and not merely its vertex set. The same remark
explains why the tile-interior blocking of Remark~\ref{rem:tib}, and not a vertex-incidence
argument, is the mechanism the cascade actually runs on.
\end{remark}

\begin{conjecture}[Advance-and-collide]\label{conj:advance}
Let $e=1$, $m=1$, $f\ge3$, and suppose an equal side carries an $a$-edge. Then the corner cascade
--- the $b$-partner of the corner tile, the $c$-chord dichotomy on successive both-ends-blocked
chords, and the $\beta$-slot at each base point reached --- terminates in every branch with a
collision: a $\beta$-slot, whose tile presents the next base edge from flanks $\{a,c\}$, at a base
position where the word of Theorem~\ref{thm:e1reduce} carries its $b$-edge, or a repeated $\gamma$.
Consequently $p=0$: Hypothesis~\ref{hyp:walls} holds for every $(1,f)$.

The conjecture is proved here for $f=2$ (Lemma~\ref{lem:ple}), $f=3$ (Theorem~\ref{thm:walls13})
and $f=4$ (Theorem~\ref{thm:walls14}, which exhausts all six admissible base words). The two ways a
branch escapes the collision in the general argument are exactly: a deviating $a$-run whose new
$c$-chord ends beyond the forced region (so that tile-interior blocking fails), and an initial side
$c$-block of length at least $2$ (so that the first-run orientation is not forced). Any proof of the
general case must close these two gaps and no others.

The second gap is now general except for its enumeration. Its ingredients carry no restriction on
$f$: \texttt{first\_run\_kill} and \texttt{through\_edge\_exclusive} are statements in
$(n_\alpha,n_\beta,n_\gamma)$ alone; Lemma~\ref{lem:topjunction} reduces to the signs of
$f^3-3f^2+1$ and $f^3-3f^2-f+1$; the chord carrying the three runs has length exactly $2b$ with $Q$
its midpoint for every $f$, since with the base corner at the origin and $u$ the unit vector along
the equal side,
\[
|2c\,u-(2f,0)|^2=4(f^2-1)^2=(2b)^2,\qquad |2c\,u-Q|^2=|Q-(2f,0)|^2=(f^2-1)^2=b^2
\]
(\texttt{chord\_two\_b}, \texttt{chord\_two\_b\_half}); and the three runs $b-a=f^2-f-1$,
$2b-a=2f^2-f-2$, $2b-c=f^2-2$ are gaps of $\langle f,\,f^2-1,\,f^2\rangle$ for every $f\ge3$
(\texttt{double\_c\_kill}), the values $11,26,14$ against $\langle4,15,16\rangle$ being the case
$f=4$. At $f=2$ the last two runs lie in the semigroup, which is where the hypothesis $f\ge3$ of
Theorem~\ref{thm:e1reduce} is used. What is not yet general is the enumeration: a side with initial
$c$-block at least $2$ begins and ends with a $c$-edge, so $n_c\ge3$ and the double-$c$ walks are
$(fk,0,f-k)$ with $1\le k\le f-3$; at $f=4$ that is the single word $c\,c\,a^4\,c$ treated in
Remark~\ref{rem:walls14}, while for $f\ge5$ there are $f-3$ walks and the argument as written
covers an initial block of exactly $2$ followed by an $a$-run.
\end{conjecture}

\begin{theorem}[Hypothesis~\ref{hyp:walls} holds at $(1,4)$]\label{thm:walls14}
No tiling of the $(1,4)$ target at $m=1$ carries an $a$-edge on an equal side; hence, by
Remark~\ref{rem:sidenoa}, no such tiling exists, re-deriving the exclusion of the $(1,4)$ instance
($N=47$) established by certified exhaustion ($12\,440$ nodes).
\end{theorem}
\begin{proof}
Suppose a side (named the left) carries an $a$-edge; then $p=1$ there
(Lemma~\ref{lem:topjunction}). Both corner partners are forced with $\gamma$ at their first
junctions, unconditionally: the reflected orientation repeats the corner tile's $\gamma$ at the
base end and overflows.

\emph{Left kills.} For left side words with initial $c$-block $1$, Lemma~\ref{lem:cchord} runs on
$P$'s $c$-chord as at $(1,3)$: base words with second letter $a$ die at $(2f,0)$ or $(3f,0)$
($\beta$-slot against the forced $b$-edge); $a\,a\,c\,b\,a\,a$ dies at $(f+c,0)$ through the
$a$-run partner; $a\,c\,b\,a\,a\,a$ dies at $(f+c,0)=(20,0)$, where its $b$-edge occupies the
four-$a$ branch's first $\beta$-slot. Five of the six words die, leaving $a\,c\,a\,b\,a\,a$. The
double-$c$ left word dies for every base word by the line-total kill of Remark~\ref{rem:walls14}.

\emph{Right kill of $a\,c\,a\,b\,a\,a$.} The right cascade is the left cascade of the reversed
word $a\,a\,b\,a\,c\,a$. The right corner tile, its partner $P'$, and the $\beta$-slot at
$(Y-f,0)=(43,0)$ force $C'=T_1'-(f,0)$, whose $b$-edge is collinear with the second right tile's
$b$-line and meets the base at $(Y-2f,0)=(39,0)$ --- an endpoint of the base word's $b$-edge
$[24,39]$. If the right side is all $c$'s, the second right tile is forced direct
(\texttt{first\_run\_kill} with $P'$'s $\gamma$), the doubled $b$-line is edged on its east side,
and its west side partitions $2b=30$, whose unique decomposition is $b+b$
(\texttt{partition\_30}): the split falls at the mid-junction, producing the two $b$-partners. If
the right side carries an $a$-run with initial block $1$, the second right tile's horizontal
$c$-edge blocks the mid-junction and the partner of $C'$ is forced with $\gamma$ there by the
through-edge figure. Either way, $\gamma$ at $(39,0)$ would double $C'$'s $\gamma$, so the
partner's $\alpha$ sits at $(39,0)$; the figure there closes with exactly one $\beta$ adjacent to
the base-left ray, and that tile must present the $b$-edge $[24,39]$ from flanks $\{a,c\}$ ---
impossible. The double-$c$ right word dies by the mirrored line-total kill. Every combination of
side and base words is exhausted; $p=0$.
\end{proof}

\begin{remark}[The reversal principle]\label{rem:reversal}
The mechanism is symmetric: the corner cascade runs from both base corners, and the right cascade
of a base word is the left cascade of its reversal. A base word survives only if both it and its
reversal survive the left cascade; at $(1,4)$ the reversal involution pairs the six words so that
none does. This is the intended proof shape for Conjecture~\ref{conj:advance}: show the left
cascade kills every word with second letter $a$ and every word whose $b$-edge occupies the
four-$a$ branch's first slot, and that the survivors' reversals always lie in the killed class.
\end{remark}

\begin{theorem}[The cascade closes every initial-block-1 configuration, for all $f$]\label{thm:e1cascade}
Let $f\ge3$, $m=1$, and suppose an equal side of the $(1,f)$ target carries an $a$-edge whose first
run is preceded by exactly one $c$-edge, and that the opposite side either has no $a$-edge or also
has initial $c$-block $1$. Then no tiling exists. Consequently Hypothesis~\ref{hyp:walls} fails to
hold at $(1,f)$ only if some hypothetical tiling has a side whose first $a$-run sits behind at least
two $c$-edges; for $f\le4$ no such configuration survives either
(Theorems~\ref{thm:walls13},~\ref{thm:walls14}), so Hypothesis~\ref{hyp:walls} holds outright there.
\end{theorem}
\begin{proof}
Three lemmas. \emph{(L1: the single-$c$ chain.)} With initial block $1$ the second side tile is the
first $a$-tile, mirrored (Lemma~\ref{lem:firstrun}), and its horizontal $c$-edge spans offsets
$[0,c]$ above the corner. The cascade produces cells $C_k=T_1+(kf,0)$ and partners
$Z_k=P+(kf,0)$ by induction: the $k$-th chord's upper end is $cu+(kf,0)$, interior to that
horizontal edge exactly when $kf<f^2$; a base word whose $b$ sits at position $q\le f$ is reached at
step $q-2\le f-2$, so every needed chord is blocked (\texttt{cascade\_reaches}). At each base point
the figure closes with exactly one $\beta$ whose flanks are $\{a,c\}$; reflected partners die by
repeating $\gamma$ at the base end. The letter $b$ therefore collides: every word with its $b$
before its $c$ dies. \emph{(L2: the $f$-$a$ branch.)} If the $c$ immediately precedes the $b$, the
$f$-$a$ branch of Lemma~\ref{lem:cchord} runs at the $c$'s position, its first $a$-tile's
$b$-partner is forced (the reflected orientation repeats $\gamma$ at the chord's division point),
its $\gamma$ lands at the next base point, and the $\beta$-slot there collides with the $b$-edge.
\emph{(L3: reversal.)} Every admissible word has its $b$ before its $c$, or its $c$ immediately
before its $b$, or its $c$ before its $b$ with a gap; in the third case the reversed word's
$b$-position is $f+3-q\in[3,f]$ and precedes its $c$ (\texttt{reversal\_covers}), so the mirrored
cascade from the opposite corner kills it — the opposite corner's chain is available because the
opposite side has no $a$-run (then the doubled $b$-line's far side partitions $2b$ uniquely as
$b+b$, \texttt{partition\_2b}, forcing the partners) or has initial block $1$ (then its horizontal
$c$-edge blocks the mid-junction and the through-edge figure forces the partner). In every case the
collision closes the configuration.
\end{proof}

\begin{lemma}[The $jb$-line partition]\label{lem:jbline}
For $1\le j<f$ the only decomposition of $j\,b$ as $xa+yb+zc$ is $x=z=0$, $y=j$. Hence a line
of length $jb$ that is covered by whole edges on a side splits there into exactly $j$ whole
$b$-edges. (Whether a given line is so covered is a separate question: see
Remark~\ref{rem:straddle}.)
\end{lemma}
\begin{proof}
Reducing $xf+y(f^2-1)+zf^2=j(f^2-1)$ modulo $f$ gives $y\equiv j$; the size bound gives $y\le j$,
and $j<f$ then forces $y=j$, hence $x=z=0$. Machine-checked (\texttt{partition\_jb}).
\end{proof}

\begin{remark}[The block cascade]\label{rem:blockcascade}
In a residual configuration (initial $c$-block of length $\ell\ge2$), the line through $j\,cu$ in
the $b$-direction reaches the base at $(jf,0)$ with total length $jb$: each step is the vector
$cu-(f,0)$, of length $b$, by the translation structure alone. Three consequences, all side-word-independent, and all conditional on the relevant
segment being uncrossed (Remark~\ref{rem:straddle}). (i)~The $(f,0)$-slot closes with exactly one $\beta$ whose flanks are
$\{a,c\}$, so the second base letter is $a$ or $c$, with its tile forced: $X_1=T_1+(f,0)$ or the
run-start $A_1$. (ii)~When the second letter is $a$, the $2b$-line $[2cu,Q_1,(2f,0)]$ is fully
west-edged by $C_2$'s and $X_1$'s $b$-edges, and at $Q_1$ the west side closes exactly
($C_2$'s $\gamma+P$'s $\beta+X_1$'s $\alpha=\pi$); Lemma~\ref{lem:jbline} forces the far side to
split at $Q_1$ into both $b$-partners, the reflected lower partner repeats $X_1$'s $\gamma$ at the
base point $(2f,0)$ and dies, and the $(2f,0)$-slot closes with exactly one $\beta$: \emph{the
third letter is never $b$}. Every residual configuration whose base word has second letter $a$ and
$b$ in position~$3$ is dead. (iii)~At the $2cu$ junction the figure $C_2$'s $\alpha+{}$upper mate's
$\gamma+{}$corner${}=\pi$ forces $C_3$ direct when $\ell\ge3$: the block's chirality forcing
propagates upward with the line, one junction per depth, and Lemma~\ref{lem:jbline} supplies the
partition at every depth $j<f$. The remaining freedom is the row-$2$ chord $[Q_1,Q_1+cu]$, whose
east side forks single-$c$/$f$-$a$ exactly as the base did: the corner block grows as a
two-dimensional array of words, one chord fork per row. Closing the residual class for all $f$
means closing this array; that is the precise remaining content of Conjecture~\ref{conj:advance}
at $e=1$.
\end{remark}

\begin{lemma}[The east fan at the fork is forced]\label{lem:eastfan}
Let $f\ge3$, $\ell\ge2$, and let the second base letter be $a$, so that the brick $X_1$ and the
mate $M$ are in place. At the fork point $Q_1=cu+(f,0)$ the tiles east of the chord form the
figure $\{\gamma,\alpha,\beta\}$ in that angular order from the base side: the mate $M_1$ of
$X_1$, then a $\beta$-tile, then the direct partner $D_2$ of $C_2$. The $\beta$-tile has flanks
$\{a,c\}$, one along the floor and one against $D_2$'s $c$-edge, so exactly two configurations
occur: $a$ on the floor (matching $M_1$'s $a$-edge), or $c$ on the floor --- the tile $A_2$.
\end{lemma}
\begin{proof}
The chord through $Q_1$ joins the side vertex $2cu$ to the base vertex $(2f,0)$, so both ends lie
on the boundary and no edge extends past them; its east side is a union of whole edges of total
length $2b$, which is $b+b$ for $f\ge3$ (Lemma~\ref{lem:jbline}), split exactly at $Q_1$. Hence
$Q_1$ is a vertex of that cover and each east tile meets it at a vertex. The lower half is
$X_1$'s $b$-edge; the single whole $b$ covering it from the east belongs to a partner of $X_1$,
and the reflected partner would repeat $X_1$'s $\gamma$ at the base junction $(2f,0)$, giving
$2\gamma>\pi$ there since $\cos\gamma=-1/(2f)<0$
(\texttt{straight\_junction\_gamma\_bound}). So it is the direct partner $M_1$, which presents
$\gamma$ at $Q_1$ --- its edges there being $b$ and $a$ --- and lays an $a$-edge east along the
floor. A straight-junction figure containing a $\gamma$ is $\{\gamma,\alpha,\beta\}$
(\texttt{straight\_junction\_cases}), so exactly two further tiles meet $Q_1$. The one adjacent to
the chord shares $C_2$'s whole upper $b$, so its angle at $Q_1$ is an endpoint angle of a
$b$-edge, namely $\alpha$ or $\gamma$; $\gamma$ is taken, so it is $\alpha$, and the reflected
alternative would again double $C_2$'s $\gamma$ at $Q_1$. That tile is therefore $D_2$. The
remaining tile carries $\beta$, whose flanks are $\{a,c\}$.
\end{proof}

\begin{theorem}[The row fork kill]\label{thm:forkkill}
Let $f\ge3$. In a residual configuration (initial $c$-block $\ell\ge2$) whose base word has second
letter $a$, the row-$2$ fork resolves single-$c$: the $\beta$-tile lays $a$ along the floor.
\end{theorem}
\begin{proof}
Lemma~\ref{lem:eastfan} leaves exactly two configurations, and
Proposition~\ref{prop:a2branch} excludes the second.
\end{proof}

\begin{proposition}[The $A_2$ branch dies]\label{prop:a2branch}
In a residual configuration (initial $c$-block $\ell\ge2$) whose base word has second letter $a$,
suppose the tile east of the row-$2$ fork is the mirrored $L=-2$ tile $A_2$, laying a horizontal
$c$-edge along the floor line east of the fork point. Then no tiling exists.
\end{proposition}
\begin{proof}
$A_2$'s $c$-edge is an edge of the tiling, so no tile crosses the floor east of the fork point,
and the south side of that edge is a union of whole edges with vertices on it. At a junction $J$
of that cover the fan below the floor is the previous element's $\beta$, the row-$1$ brick's
$\alpha$, and a remaining wedge of exactly $\gamma$. Since $\gamma<\pi$, no tile can meet $J$
through the interior of one of its edges (that contributes exactly $\pi$), so the wedge is filled
by tiles having a vertex at $J$: the fill is $\{\gamma\}$ or $\{2\alpha,\beta\}$
(\texttt{straight\_junction\_cases}). The wedge is bounded by the floor and by the brick's
$b$-edge running \emph{down} to the base. A filler laying $c$ along that ray leaves the target
(it exceeds the ray by $c-b=1$ past the base); a filler laying $a$ along it makes that $a$ the
first element of the $b$-edge's east cover, leaving $f^2-1-f$, which admits no completion
(\texttt{east\_cover\_gap}); a filler laying a whole $b$ along it is the reflected partner of the
brick and repeats $\gamma$ at the base junction
(\texttt{straight\_junction\_gamma\_bound}). Hence every fill is the brick's mate, the cover is
$f$ whole $a$-edges, their feet force base junctions at $(kf,0)$ for $k=2,\dots,f+1$, and the
first $f+1$ letters are all $a$ --- contradicting the $b$ at position at most $f$. (The non-brick
$a$-up placements are excluded by \texttt{alpha\_vertex\_gap} and \texttt{anti\_brick\_side} as
before.)
\end{proof}

\begin{remark}[Why $f\ge3$, and what the $(1,2)$ tilings show]\label{rem:whyf3}
The hypothesis $f\ge3$ in Lemma~\ref{lem:eastfan} is doing real work, and the $(1,2)$ member shows
exactly what it buys. At $f=2$ one has $2b=6=b+b=a+c=3a$: the chord's partition is not unique, the
split need not fall at $Q_1$, and a tile may straddle there. The kernel-checked $99$-tiling does
precisely this --- its tile with vertices $(4,0)$,
$\bigl(\tfrac{23}4,\tfrac{\sqrt{15}}4\bigr)$, $(5,\sqrt{15})$ lays a $c$-edge from the base along
the chord which covers the row-$1$ brick's entire $b$-edge and overshoots $Q_1$ by exactly
$c-b=1$. For $f\ge3$ that escape costs $2b-c=f^2-2$, which lies below $b$ and is not divisible by
$f$: a gap (\texttt{gap\_2b\_minus\_c}). So the phenomenon is real, is realised, and is confined
to $f=2$.

Two methodological points are worth recording. First, the overshoot test is \emph{directional}: a
$c$ laid downward from a floor junction exceeds the boundary and dies, while a $c$ laid upward
from the base overshoots an interior point and is legal. Second, an argument of the form ``a
segment pinned at both ends is covered by whole edges'' is sound only where the ends are genuinely
pinned --- on the boundary, or against a tile edge that blocks continuation --- and should say
which. Both uses in Theorem~\ref{thm:e1cascade} are of that kind: its chords are blocked above by
the first $a$-tile's horizontal $c$-edge, and its doubled $b$-line runs boundary to boundary.
\end{remark}

\begin{lemma}[The wall climb]\label{lem:wallclimb}
For every $\ell\ge2$ and $2\le j\le\ell$, the block tile $C_j$ is direct and its mate $M_j$ is
forced, with $M_j$'s $c$-edge $[Q_{j-1},Q_j]$ continuing the interior wall; the first run tile
above the block is mirrored. This is unconditional: no base-letter hypothesis and no line
closure is used.
\end{lemma}
\begin{proof}
Induction on $j$. $C_2$ is direct by the $cu$-junction figure ($T_1$'s $\alpha+P$'s
$\gamma+\beta$). Given $C_j$ direct, the west fan at $Q_{j-1}$ consists of $\beta$ ($P$'s for
$j=2$, else $M_{j-1}$'s), $C_j$'s $\gamma$, and a wedge between $C_j$'s $b$-ray and the $u$-line
of exactly $\pi-\gamma-\beta=\alpha$: it is filled by a single $\alpha$-tile ($\alpha$ is the
minimal angle). Its flank along $C_j$'s $b$-ray cannot be a $c$-edge — it would overshoot the
boundary point $jcu$ — so it is a whole $b$-edge shared with $C_j$, and the reflected partner is
excluded by the wedge's $\alpha$: the filler is the mate $M_j$, its $c$-edge running up the
$u$-line. At the side junction $jcu$, $C_j$'s $\alpha+M_j$'s $\gamma$ force the next side
tile's corner to $\beta$: $C_{j+1}$ is direct while the block lasts ($\alpha+\gamma+\alpha$ and
$\alpha+\gamma+\gamma$ are impossible straight figures), and the first run tile is mirrored when
it ends.
\end{proof}

\begin{remark}[The pincer]\label{rem:pincer}
The assembly frame for the residual class: each corner runs its own chain, and a chain needs
only its \emph{own} second letter to be $a$. If $cp>bp$ the left chain meets the $b$ first
(second letter $a$ since $cp>bp\ge3$); if $cp<bp$ the reversed word has $b$ before $c$ and the
right corner — whose side is all-$c$, block-$1$, or again residual — runs the mirrored chain
(reversed second letter $a$ since $bp\le f$). Both chains advance through slots, lines and row
forks; with the three leaks excluded the forks resolve to bricks, the lines close at every
depth (Lemma~\ref{lem:jbline}), and the colliding $b$ kills the configuration. Hypothesis
hyp:walls at $e=1$ is therefore reduced to the leak exclusion.
\end{remark}


\begin{theorem}[L2 at every reached slot]\label{thm:l2slot}
Let $f\ge3$ and suppose the chain from a corner has closed slots $1,\dots,k$. If the base letter
presented at slot $k$ is $c$ and the next letter is $b$, no tiling exists. This holds for every
block length of that corner's side.
\end{theorem}
\begin{proof}
The slot-$k$ $\beta$-tile lays $c$ on the base and its $a$ up the $u$-ray from $(kf,0)$, which is
the $Z_{k-1}$-mate's $c$-edge (for $k=1$: $P$'s $c$-edge) --- an edge of the chain. The branch
tile's $b$ runs from the $fu$-division point of that edge to the base point $(kf+c,0)$; its upper
end abuts the edge's interior, so no covering edge crosses it, and the lower end is on the base:
the far side is a single whole $b$ (\texttt{partner\_unique}). The reflected partner repeats
$\gamma$ at the division point, an interior straight junction, and $2\gamma>\pi$; the mate's
$\gamma$ lands at $(kf+c,0)$, the figure there closes with exactly one $\beta$ whose flanks are
$\{a,c\}$, and it must present the $b$-edge: impossible.
\end{proof}

\begin{theorem}[The block-two chain runs to arbitrary depth]\label{thm:elltwo}
Let $f\ge3$ and let a side carry its first $a$-run behind exactly two $c$-edges. Then the chain
from that corner closes every slot it needs: the base letters are presented one by one, a $b$
collides at its slot, and a $c$ followed by $b$ collides by Theorem~\ref{thm:l2slot}. In
particular a base word whose $b$ precedes its $c$ dies from this corner.
\end{theorem}
\begin{proof}
By Lemma~\ref{lem:wallclimb}, $T_{a1}$ is mirrored and lays its horizontal $c$-edge at height
$2cu_y$: the ceiling. The column-$k$ line through $((k+1)f,0)$ meets the ceiling line after two
pieces, at a point interior to $T_{a1}$'s $c$ (offset $kf<c$), so its west-edged portion is
$X_k$'s $b$ plus the row-$2$ brick's $b$, pinned above by the ceiling edge and below by the base:
\texttt{partition\_2b} splits the east side at the joint, the reflected partners die by repeating
$\gamma$ (at the base, resp.\ at the joint), and both mates are forced: the next slot closes. The
row-$2$ bricks advance by the fork at each column: the fork line is pinned the same way, the fan
at the fork point is $\gamma\,|\,\alpha\,|\,\beta$ exactly as in Lemma~\ref{lem:eastfan} with the
ceiling in place of the wall, and the $A$-branch dies by Proposition~\ref{prop:a2branch}, whose
south feet lie on the base and whose leak inventory is closed by \texttt{east\_cover\_gap} and
the directional overshoot (both valid one strip above the base). Induction on the column.
\end{proof}

\begin{theorem}[Reach three behind a thick block, and the pincer window]\label{thm:depthwindow}
Let $f\ge3$. From a corner whose side has initial block $\ell\ge3$ (including the all-$c$ side),
slots $1,2,3$ close: the lines for columns $k\le2$ terminate on the wall
(Lemma~\ref{lem:wallclimb}) with $k+1$ pieces and split by Lemma~\ref{lem:jbline}; the single
row-$2$ brick needed is forced by Theorem~\ref{thm:forkkill}. Consequently every residual
configuration dies unless its base word escapes all four kills: direct-$b$ from the $b$-first
corner at depth $\le4$, and L2 (Theorem~\ref{thm:l2slot}) from the $c$-first corner when the $c$
sits at depth $\le4$ with the $b$ adjacent. For $f\le5$ no admissible $(bp,cp)$ escapes
(\texttt{pincer\_window}).
\end{theorem}

\begin{corollary}\label{cor:walls15}
Hypothesis~\ref{hyp:walls} holds at $(1,5)$: together with blocks $\le2$
(Theorems~\ref{thm:e1cascade},~\ref{thm:elltwo}) the window covers every configuration. This is
the first member closed with no engine certificate.
\end{corollary}

\begin{lemma}[The apex]\label{lem:apex}
The target's apex angle is $\pi-2\beta=3\alpha$, and the only vertex figure filling it is three
$\alpha$-corners (\texttt{apex\_figure}). Since the sides carry no $b$-edge
(\texttt{side\_no\_b\_uncond}), the two outer apex tiles lay $c$-edges on the sides --- the top
letters of both side words are $c$ --- and the three wedges' rays sit at labels $-3,-1,+1,+3$.
The apex therefore admits exactly two configurations, distinguished by the middle tile's
chirality: its $b$ shares the $-1$-ray with the left tile's $b$ (both $\gamma$-ends meeting
below), or its $c$ covers that $b$ with a T-junction at depth $b$. The side conclusion is
specific to $m=1$: the kernel-checked tilings at $m\ge2$ carry $b$-edges on their sides, and
their apex regions, while all realizing the forced three-$\alpha$ figure, diverge from one
another immediately below it. Any forcing chain from the apex must therefore anchor on the
$m=1$ side structure; no $m$-independent reach exists.
\end{lemma}

\begin{remark}[The straddler: the unique leak at rows $\ge3$]\label{rem:straddler}
At every strip junction of a row-$j$ $A$-branch the wedge is exactly $\alpha$, and its filler
admits exactly two label assignments: the mirrored $L=-2$ mate (whole $b$ shared with the run
tile; the strip structure stays rigid) or the direct $L=0$ \emph{straddler} --- $b$ horizontal
under the next $c$-line, $c$ laid along the run tile's $b$-ray, covering that whole $b$ and
overshooting the run tile's $\alpha$ by exactly $c-b=1$, poking below the floor line just past
the horizontal $c$-edge's east end. At row $2$ the analogous fill dies against the boundary
($2b-c=f^2-2$ is a gap, \texttt{gap\_2b\_minus\_c}); at rows $\ge3$ the poke is interior and no
local figure excludes it. The straddler does not cascade: it couples strips only upward, so no
descent-to-the-base argument applies. Killing this one shape at rows $\ge3$ is equivalent to the
full chain reach, hence to Conjecture~\ref{conj:advance} at $e=1$.
\end{remark}

\begin{lemma}[Descent identities and the ladder]\label{lem:ladder}
Let $f\ge3$. (i) A line in the $b$-direction (label $-1$) crosses the horizontal strip lines at
intervals of exactly $b$ along itself and advances exactly $c$ horizontally per crossing:
$\sin\beta=b\sqrt D/(2f^3)$ with $D=4f^2-1$, and $fu_x+b\,|\cos2\gamma|=c$
(\texttt{descent\_ident}, \texttt{sinb\_ident}). (ii) A straddler's $c$-edge therefore covers the
run tile's $b$ and pokes one unit past its $\alpha$-vertex T-junction; at every T-junction on the
line the fan-adjacent tile lays a whole further edge along it, so the straddler forces a doubly
edged ladder: east breaks in $b+\langle a,b,c\rangle$, west breaks in $c+\langle a,b,c\rangle$,
staggered by $c-b=1$. (iii) The ladder cannot terminate on the base: the termination abscissa is
$\bigl(f^3(m+1)-K(f^4-3f^2+1)\bigr)/f^2$ with $K=(j-1)f+i$, and $f^4-3f^2+1\equiv1 \pmod{f^2}$
forces $f^2\mid K$, impossible since $K\equiv i\not\equiv0\pmod f$
(\texttt{ladder\_no\_base}). (iv) The ladder can terminate only at an interior mutual vertex,
where the cumulative east and west sums agree; the stagger makes this $s_1=s_2+1$ with
$s_1,s_2\in\langle a,b,c\rangle$, realized earliest by $(s_2,s_1)=(b,c)$: every straddler forces
at least $b+c=2f^2-1$ of doubly edged ladder, descending two full strips and gliding exactly $c$
east of the strip ends.
\end{lemma}

\begin{lemma}[Termination is one condition, boundary or interior]\label{lem:termination}
A ladder terminates only where \emph{both} of its staggered covers end: at an interior mutual
vertex by definition, and at a boundary point because the boundary ends both covers at once. So
side exit is not a separate escape route; every termination satisfies the same condition. With
$c=b+1$ that condition reads: $T-b$ and $T-c$ are consecutive members of
$S=\langle a,b,c\rangle$. Consecutive members of $S$ begin only at $f^2$
(\texttt{consecutive\_gap}), so
\[
  T \;\ge\; b+c \;=\; 2f^2-1 .
\]
\end{lemma}
\begin{proof}
Write $m+1=T-b$ and $m=T-c$, both in $S$. Reducing $xf+y(f^2-1)+zf^2$ modulo $f$ shows an element
of $S$ congruent to $-y$; if $m+1<f^2$ then both representations have $z=0$ and $y\le1$, and the
four cases give $f\mid1$, $f\mid2$, or a direct size contradiction. Hence $m+1\ge f^2$ and
$T=b+(m+1)\ge b+c$.
\end{proof}

\begin{theorem}[Straddlers need three strips, so rows $1,2,3$ are clean]\label{thm:strippbound}
Let $f\ge3$. The target is exactly $f$ strips tall ($f^3\sin\beta$ over $c\sin\beta$), and a
ladder descends exactly $b$ along itself per strip (Lemma~\ref{lem:ladder}). By
Lemma~\ref{lem:termination} its length satisfies $T\ge b+c>2b$, so it crosses strictly more than
two strips; since it cannot terminate on the base (Lemma~\ref{lem:ladder}(iii)), a straddler whose
ladder starts at the floor of row $j$ requires $j-1>2$. Hence no straddler occurs at rows $1$,
$2$ or $3$, the row-$3$ fork admits only the two branches of Lemma~\ref{lem:eastfan}. The
$A$-branch's exclusion at row $3$ invokes Proposition~\ref{prop:a2branch} whose south-cover feet
land on the row-$2$ floor, edged only as far as the staircase has reached: that step is
incomplete as written, and the structural claim of reach $4$ is open pending its repair. The
$(1,6)$ conclusion itself is unaffected: it is certified by the independent exhaustion
\texttt{w107}.
\end{theorem}

\begin{corollary}\label{cor:walls16}
Hypothesis~\ref{hyp:walls} holds at $(1,6)$, hence for every $f\le6$: with reach $4$ the four
kills cover depth $\le5$ and no admissible $(bp,cp)$ escapes (\texttt{pincer\_window\_four}).
The $(1,5)$ and $(1,6)$ cases are additionally certified by independent exhaustions of the full
$m=1$ targets ($132{,}519$ and $1{,}251{,}382$ nodes, instances \texttt{w74}, \texttt{w107});
the $(1,7)$ exhaustion is running.
\end{corollary}

\begin{lemma}[Corner lines are column lines]\label{lem:columnlines}
Let $f\ge3$, $m=1$. Every $b$-direction line through two vertices of the corner lattice
$\mathbb Z(a,0)+\mathbb Z\,cu$ meets the base at $(kf,0)$ and the left side at $k\,cu$ for an
integer $k$, and its horizontal distance to the side at row $r$ is exactly $(k-r)a$. Its length is
$kb$ and both ends lie on the boundary, so for $k<f$ Lemma~\ref{lem:jbline} splits both of its
sides into exactly $k$ whole $b$-edges, breaking at every strip junction on both sides at once:
the line is \emph{unstaggered}. Since $k\,c\le f^3$ forces $k\le f$, the only $b$-direction line
of the left corner that escapes is $k=f$ --- the line from $(f^2,0)$ to $f\,cu$, which is the
apex. A ladder is staggered by $c-b=1$, so no straddler sits on a corner line with $k<f$.
\end{lemma}

\begin{lemma}[The chain never needs the apex line]\label{lem:noapexline}
The base $b$-letter sits at a position $bp\in[3,f]$. With the preceding letters forced to $a$ it
begins at $\bigl((bp-1)f,0\bigr)$, and the collision figure there consumes only the brick and
mate of column $bp-2$, hence only the lines $L_k$ with $k\le bp-1\le f-1$. Every line the chain
requires therefore has $k<f$, where Lemma~\ref{lem:columnlines} applies
(\texttt{chain\_needs\_small\_lines}); the apex line $L_f$, the one line escaping
Lemma~\ref{lem:jbline}, is never required.
\end{lemma}

\begin{remark}[The residue, as one statement, and why it must be $m=1$]\label{rem:onegap}
Combining Lemmas~\ref{lem:columnlines} and~\ref{lem:noapexline}: no straddler can sit on any line
the chain needs, \emph{provided no tile crosses that line}. Since $L_k$ runs boundary to boundary
and every needed $k$ satisfies $k<f$, this crossing statement is the only remaining hypothesis at
$e=1$; with it, every fork closes, the chain reaches the $b$, and Hypothesis~\ref{hyp:walls}
follows at $e=1$ for every $f$.

The statement is false without the hypothesis $m=1$. Tested against the five kernel-checked
tilings, all of which have $m\ge2$, tiles cross boundary-to-boundary $b$-lines freely: $42$, $15$,
$24$, $10$ and $98$ crossing incidences respectively. This is consistent with
Remark~\ref{rem:blockbreaks} --- at $m=2$ the corner block can break --- and it settles the shape
of any proof: the crossing statement admits no $m$-independent argument, and must consume the
$m=1$ hypothesis exactly where the corner block is rigid. That input is not supplied here.
\end{remark}

\begin{lemma}[The vertex census]\label{lem:census}
In every tiling of the base-$\beta$ target, write $n_1,n_2$ for the numbers of straight figures
$\{\gamma,\alpha,\beta\}$ and $\{3\alpha,2\beta\}$, and $v_1,\dots,v_4$ for the interior figures
$\{\beta,3\gamma\}$, $\{2\alpha,2\beta,2\gamma\}$, $\{4\alpha,3\beta,\gamma\}$, $\{6\alpha,4\beta\}$.
The base corners fill uniquely as $\{\beta\}$ and the apex as $\{3\alpha\}$, and balancing the
three corner types across the $N$ tiles gives
\[ v_1 \;=\; 1 + n_2 + v_3 + 2v_4 \qquad(\texttt{vertex\_census}). \]
In particular $v_1\ge1$: a $\{\beta,3\gamma\}$ vertex is mandatory, and each $\gamma$-poor figure
demands one more. The identity is verified exactly on all five kernel-checked tilings. It is also
the \emph{entire} content of corner counting: the three balance equations admit exactly one
relation, and it does not involve the side structure --- corner counting alone cannot see $p$,
consistent with the census-based routes logged as dead.
\end{lemma}

\begin{definition}[The frontier invariant]\label{def:frontier}
A \emph{frontier} is a finite family of disjoint uncovered polygons inside the target together
with, for each target side, the multiset of tile-edge types already laid on it. Say a frontier is
\emph{admissible} if all five of the following hold.
\begin{enumerate}
\item[(P1)] Each polygon's area is a positive integer multiple of the tile's area, and the
multiples sum to the number of tiles remaining.
\item[(P2)] Every maximal straight run of a polygon's boundary joining two convex corners has
integral length lying in $\langle a,b,c\rangle$.
\item[(P4)] Every convex corner of every polygon has angle at least $\alpha$.
\item[(P5)] The per-side edge-type multiset is dominated componentwise by an admissible walk
profile.
\item[(P6)] An edge lying in a target side and meeting a corner of that side has the type forced
at that corner.
\end{enumerate}
Admissibility is decidable, and it is preserved by no tiling step: a tiling of the target is
exactly a sequence of placements through admissible frontiers ending at the empty frontier.
\end{definition}

\begin{lemma}[The angle threshold]\label{lem:anglethreshold}
For the tile $(a,b,c)=(f,f^2-1,f^2)$ one has $\cos\alpha=(2f^2-1)/(2f^2)$
(\texttt{cos\_alpha\_closed}), alongside $\cos\gamma=-1/(2f)$ and
$\cos\beta=(3f^2-1)/(2f^3)$. Condition (P4) is therefore the statement that the uncovered region
never acquires a corner sharper than the tile's own sharpest angle.
\end{lemma}

\begin{remark}[The exhaustion invariant, and the shape of a uniform proof]\label{rem:enginelaw}
Definition~\ref{def:frontier} is the invariant maintained by the exhaustive search whose verdicts
certify the small members. Its content relative to the hand arguments of this paper is that (P2)
generalises the semigroup gap lemmas used throughout, while (P4) is a constraint on the uncovered
region rather than on completed vertices, and is not used anywhere in the arguments above. In the
certified cases the searches reach a contradiction long before the tiling is complete. For
$f=2,3,5,6,7$ the maximum refutation depth is $6,13,24,29,35$ against $N=3f^2-1$ placements, so
the fraction of the tiling ever built is $0.55,\,0.50,\,0.32,\,0.27,\,0.24$ and falling, while the
share of (P4) violations rises from $50\%$ to $98\%$. The vanishing fraction is the robust part of
this observation. The linear form is not: the ratio depth$/f$ reads $3.00,\,4.33,\,4.80,\,4.83,\,
5.00$ and is still increasing at $f=7$, so five instances do not separate a bound $Cf$ from a
slowly superlinear law. A proof of Hypothesis~\ref{hyp:walls} for all $f$ would follow from a
uniform bound of the form: from a corner-anchored start with $p\ge1$, no admissible frontier
survives past depth $Cf$. That statement is open and the figures quoted are measurements, not
evidence for the constant.
\end{remark}

\begin{lemma}[The thick-member monochotomy]\label{lem:monochotomy}
Let $e\ge2$, $e<f$, $\gcd(e,f)=1$, tile $(a,b,c)=(ef,\,f^2-e^2,\,f^2)$. The only decomposition of
$c$ as $xa+yb+zc$ is the single $c$: writing $xef+y(f^2-e^2)=f^2$, coprimality forces
$e\mid y-1$, and both $y=1$ (then $xf=e<f$) and $y\ge e+1$ (then the sum exceeds $f^2$ since
$f^2>e^2+e$) are impossible (\texttt{c\_chord\_unique\_thick}). Consequently the $f$-$a$ branch
of Lemma~\ref{lem:cchord} is an $e=1$ phenomenon: at thick members every $c$-chord fork of the
corner chain is forced, and the chain machinery, where it applies, is stiffer than at $e=1$.
\end{lemma}

\begin{lemma}[Thick-side quantization]\label{lem:sidequant}
At any member $(e,f)$, $e\ge1$, a side's edge counts satisfy $xe+zf=f^2$ with no $b$
(\texttt{side\_no\_b\_uncond}), so $f\mid x$ by coprimality (\texttt{side\_quantized}): the number
of $a$-edges on a side is a multiple of $f$. For $e\ge2$ this quantizes the walls violation: any
bound $p<f$ forces $p=0$, the side condition of Hypothesis~\ref{hyp:walls}.
\end{lemma}

\begin{lemma}[The thick base trichotomy]\label{lem:basetri}
For a separated member ($f^2>2ef+e^2$) the base admits exactly three edge decompositions:
$(x,y,z)\in\{(0,e,2e),\,(f,e,e),\,(2f,e,0)\}$ (\texttt{base\_trichotomy}); the walls form is the
middle one, $f$ $a$-feet, $e$ $b$-edges, $e$ $c$-feet. Close pairs ($f^2\le2ef+e^2$) admit in
addition finitely many decompositions with $y>e$, computable per member.
\end{lemma}

\begin{lemma}[The shadow at a $c$-corner]\label{lem:shadow}
Suppose a base corner tile lays $c$ on the base (then it is mirrored with its $a$-edge first on
the side), and the first side run has length at least $2$. The run is $\gamma$-upper by
absorption, so its second tile lays a horizontal $c$-edge at height $a\sin\beta$ --- exactly one
tile-height (\texttt{strips\_tall} family identities). Every tile in the strip below spans it and
carries its $c$ on the floor or the ceiling; a tile with $b$ on the floor and $c$ on the ceiling
would have two parallel sides. Hence every base edge meeting the shadow $(0,\,a\cos\beta+c)$ is a
$c$-edge. Consequences: at such a corner either the first side run has length $1$, or the second
base edge is a $c$ and no $b$-edge meets the shadow. At $e=1$ the two-sided shadow window is
always smaller than the $b$-edge (\texttt{shadow\_footage\_e1}: the comparison reduces to
$3f^2>1$), an independent exclusion of $c$-corners; at $e\ge2$ the comparison reverses, and the
shadow is a constraint rather than a kill.
\end{lemma}

\begin{lemma}[The thick base dichotomy]\label{lem:basedi}
The $\gamma$-trap (Proposition~\ref{M-prop:gammatrap}, every side of a base-$\beta$ target carries
a $c$-edge, general in $(e,f)$) forces $z\ge1$, deleting the third column of
Lemma~\ref{lem:basetri}. At a separated thick member only two base decompositions survive:
$(0,e,2e)$ and the walls form $(f,e,e)$ (\texttt{base\_dichotomy\_thick}).
\end{lemma}

\begin{lemma}[The $c$-corner is rigid]\label{lem:ccorner}
Suppose a base corner tile lays $c$ on the base. Since the corner angle is $\beta$, with flanks
$\{a,c\}$, that tile lays $a$ on the side. Its $b$-edge is a chord joining the base point $(c,0)$
to the side point $a\,u$, so both of its ends lie on the boundary and its far side carries a
single whole $b$ (\texttt{partner\_unique}, \texttt{lean/PentagonLemma.lean}). The corner tile presents $\alpha$ at the base end of
that chord and $\gamma$ at the side end; the reflected partner would repeat $\gamma$ at the side
end, where the junction is straight and $2\gamma>\pi$ because $\cos\gamma=-e/2f<0$. Hence the
direct partner is forced, contributing $\gamma$ at the base end and $\alpha$ at the side end, and
each end closes with exactly one $\beta$ whose flanks are $\{a,c\}$: the second base edge and the
second side edge are each an $a$ or a $c$.

In the base column $(0,e,2e)$ of Lemma~\ref{lem:basedi} the base carries no $a$-edge, so both
corner tiles lay $c$ on the base and the hypothesis holds at both corners: both sides begin with
an $a$-edge, and by Lemma~\ref{lem:sidequant} each side then carries at least $f$ of them. The
$\gamma$-trap bounds the side profiles to $a$-count $fj$ with $je\le f-1$.

This is the unconditional entry point that Lemma~\ref{lem:pentagon} lacks: that lemma disposes of
the walk $(0,e,2e)$ only after both dual blocks are complete, and its stubs are region-boundary
segments only under that hypothesis, whereas the forcing here uses nothing but the boundary
position of the corner chord and $2\gamma>\pi$. What remains for the column is to run the chain
from the rigid $c$-corner far enough to create the stubs of Lemma~\ref{lem:pentagon}
unconditionally; the stub arithmetic itself is already general and machine-checked
(Lemma~\ref{lem:stubgap}).
\end{lemma}

\begin{theorem}[The column $(0,e,2e)$ admits no tiling]\label{thm:n1}
Let $(e,f)$ be a separated member with $e\ge1$, $\gcd(e,f)=1$, $f\ge2$, and suppose the base walk
is $(0,e,2e)$: no $a$-edge on the base. Then no tiling exists.
\end{theorem}
\begin{proof}
Write $\gamma=2\alpha+\beta$; a wedge $X\alpha+Y\beta$ filled by corner counts $(x,y,z)$ satisfies
$x+2z=X$, $y+z=Y$, so the fills of $\beta$, of $\alpha$ and of $\alpha+\beta$ are
$\{\beta\}$, $\{\alpha\}$ and $\{\alpha,\beta\}$ respectively (\texttt{fill\_beta},
\texttt{fill\_alpha}, \texttt{fill\_alpha\_beta}).

The corner angle is $\beta$, so a single tile $T_0$ occupies the corner; its flanks are $\{a,c\}$
and the base carries no $a$-edge, so $T_0$ lays $c$ on the base and $a$ on the side, and its
$b$-edge joins $(c,0)$ to $a\,u$. At the side vertex $a\,u$ the junction is straight and $T_0$
presents $\gamma$, leaving $\alpha+\beta$, so exactly two further tiles meet there, one with
$\alpha$ and one with $\beta$. Their order is forced. The edge $T_0$ presents along the chord is
an edge of a tile, so no tile crosses it, and both of its ends lie on the boundary, so no edge
overruns them; hence the far side of that chord is a union of whole edges of total length $b$,
which by \texttt{partner\_unique} is a single $b$. The tile carrying it meets $a\,u$ at an
endpoint of a $b$-edge, so presents $\alpha$ or $\gamma$ there, and $\gamma$ is already taken. It
is therefore the $\alpha$-tile, and it is the partner $P_0$, forced with no hypothesis on corner
blocks.

Inductively, suppose $T_k$ (a translate of $T_0$ by $(kc,0)$) and $P_{k-1}$ are in place. At the
base vertex $((k+1)c,0)$, $T_k$ presents $\alpha$ and $P_k$ will present $\gamma$; at the vertex
$V_k=(kc,0)+a\,u$, $T_k$ presents $\gamma$ and $P_{k-1}$ presents $\beta$, leaving $\alpha$, whose
unique filler has flanks $\{b,c\}$; a $c$ laid along $T_k$'s $b$-ray would run from $V_k$ past
$((k+1)c,0)$, a point of the base, and so leave the target, so the filler lays a whole $b$ there
and is the partner $P_k$. At the
base vertex $(kc,0)$, $T_{k-1}$'s $\alpha$ and $P_{k-1}$'s $\gamma$ leave exactly $\beta$, whose
filler has flanks $\{a,c\}$: the next base edge is an $a$ or a $c$, hence a $c$, and its tile is
$T_k$. The induction therefore forces every base edge to be a $c$-edge.

The base has length $e(3f^2-e^2)$, so this requires $f^2\mid e(3f^2-e^2)$, hence $f^2\mid e^3$,
hence $f=1$ by coprimality (\texttt{base\_not\_all\_c}), contradicting $f\ge2$.
\end{proof}

\begin{corollary}\label{cor:basewalls}
At a separated thick member the base walk is the walls form $(f,e,e)$: Lemma~\ref{lem:basetri}
gives three columns, the $\gamma$-trap deletes $(2f,e,0)$ (Lemma~\ref{lem:basedi}), and
Theorem~\ref{thm:n1} deletes $(0,e,2e)$.
\end{corollary}

Separation enters that chain only through Lemma~\ref{lem:basetri}, to rule out columns with $y>e$.
That can be done under the weaker hypothesis $f>2e$, so the trichotomy and the dichotomy both
extend. Proposition~\ref{prop:sharpcolumn} is not available for this purpose, because it assumes
$x,z\ge1$ while Lemma~\ref{lem:basetri} allows $x=0$ and $z=0$; the argument below is run with
$x,z\ge0$ throughout.

\begin{corollary}[The base trichotomy and dichotomy without separation]\label{cor:basedi2e}
Let $\gcd(e,f)=1$ and $f>2e$. Then the base admits exactly the three decompositions
$(0,e,2e)$, $(f,e,e)$, $(2f,e,0)$ of Lemma~\ref{lem:basetri}, and after the $\gamma$-trap exactly
the two of Lemma~\ref{lem:basedi}, namely $(0,e,2e)$ and the walls form $(f,e,e)$.
\end{corollary}

\begin{proof}
Reducing $xa+yb+zc=e(3f^2-e^2)$ modulo $f$ kills $a=ef$ and $c=f^2$ and leaves
$-ye^2\equiv-e^3$, so $f\mid e^2(y-e)$ and hence $f\mid y-e$ by coprimality. Write $y=e+kf$; since
$0\le y$ and $e<f$, $k\ge0$.

Suppose $k\ge1$. The base equation is then $xe+zf=2ef-kb$ (Proposition~\ref{prop:closepaircolumns}(i),
whose derivation uses no positivity), and $f\mid x-ke$, so $x=ke-mf$ for an integer $m$; from $x\ge0$,
$mf\le ke$. The identity
\[
(ke-mf)e+k(f^2-e^2)-2ef \;=\; f\bigl(kf-(m+2)e\bigr)
\]
together with $z\ge0$, i.e.\ $xe+kb-2ef\le0$, gives $kf\le(m+2)e$ after dividing by $f>0$. Now
$2m<k$: for $m\le0$ this is clear from $k\ge1$, and for $m\ge1$ it follows from
$2me<mf\le ke$. Finally $f>2e$ and $k\ge1$ give $2ke<kf\le(m+2)e$, hence $2k<m+2$, and with
$2m\le k-1$ this yields $4k<k+3$, i.e.\ $k<1$, a contradiction.

So $y=e$. Then $f\mid x$, say $x=jf$, and the base equation becomes $je+z=2e$, i.e.\
$z=e(2-j)$, so $z\ge0$ gives $j\in\{0,1,2\}$ and the three columns are exactly those listed. The
$\gamma$-trap forces $z\ge1$, deleting $j=2$. We cite the $\gamma$-trap
(Proposition~\ref{M-prop:gammatrap}) rather than Lemma~\ref{lem:basedi}: the latter's formal
counterpart \texttt{base\_dichotomy\_thick} carries a separation hypothesis, whereas the
$\gamma$-trap is proved for an arbitrary side of a base-$\beta$ target and its corner input
(Proposition~\ref{M-prop:cornerfig}) puts no $\gamma$ at a base corner or at the apex, so it applies
here unchanged.
\textup{(\texttt{no\_extra\_column\_of\_f\_gt\_two\_e}.)}
\end{proof}

\begin{remark}[$f>2e$ is sharp, and $(1,2)$ is the unique exception]\label{rem:f2esharp}
The inequality cannot be weakened to $f\ge2e$. Coprimality makes $f=2e$ possible only at $e=1$,
$f=2$, and that member does carry a fourth decomposition: $1\cdot a+3\cdot b+0\cdot c=2+9=11=e(3f^2-e^2)$,
i.e.\ $(x,y,z)=(1,3,0)$ with $y=3>e$. It escapes the argument above exactly at $z=0$. Checked by
structured enumeration of the base equation on all $165038$ members with $f>2e$, $e\le300$: no
decomposition outside the three, and the same enumerator does find $(1,3,0)$ at $(1,2)$.
\end{remark}

\begin{remark}[The band this opens, and what still blocks the walls form there]\label{rem:band2e}
Separation is $f>(1+\sqrt2)e$, so Corollary~\ref{cor:basedi2e} is strictly weaker in hypothesis and
covers the band $2e<f\le(1+\sqrt2)e$ that Lemma~\ref{lem:basetri} left to a per-member
computation: $5064$ of the $17296$ close pairs with $e\le200$ lie in it. For those members the base
dichotomy now holds for the same reason it holds at a separated member.

To reach the walls form one further needs Theorem~\ref{thm:n1}. Its separation hypothesis is indeed
inert --- no step of its proof invokes $f^2>2ef+e^2$, and its arithmetic inputs
(\texttt{partner\_unique}, $c>b$, \texttt{base\_not\_all\_c}) were checked on all $36896$ coprime
members with $e\le200$, $e<f<4e+2$, of which $17296$ are close pairs, with no failure. But deleting
the hypothesis does not deliver the conclusion, because the proof of Theorem~\ref{thm:n1} has a gap
at $e\ge2$, recorded next.
\end{remark}

\begin{remark}[A gap in the induction of Theorem~\ref{thm:n1}]\label{rem:n1gap}
The induction step of Theorem~\ref{thm:n1} argues that at $V_k=(kc,0)+a\,u$ the tiles $T_k$ and
$P_{k-1}$ present $\gamma$ and $\beta$, ``leaving $\alpha$''. That is the straight-junction identity
$\pi-\gamma-\beta=\alpha$, which follows from $3\alpha+2\beta=\pi$ and $\gamma=2\alpha+\beta$. It is
legitimate at $k=0$, where $V_0=a\,u$ lies on an equal side and the junction is genuinely straight.
For $k\ge1$ the point $V_k$ is \emph{strictly interior}: the angle around it is $2\pi$, and
\[
2\pi-\gamma-\beta \;=\; 4\alpha+2\beta \;=\; \pi+\alpha \;\ne\;\alpha .
\]
The discrepancy is exactly $\pi$. For instance at $(e,f)=(2,5)$ one has
$V_1=(30.68,8.23)$, interior to the target $(0,0),(142,0),f^3u$; the residue there is $203.07^\circ$,
not the $\alpha=23.07^\circ$ the step requires. The same holds at every member we computed,
$(2,3)$, $(3,8)$, $(1,4)$ and $(3,7)$ among them. With the residue $\pi+\alpha$ the partner $P_k$ is
not forced, and the chain stops after $T_0,P_0,T_1$; that much is exactly
Proposition~\ref{M-prop:cornerpara}.

At $e=1$ the theorem is nevertheless true, by a different and much shorter route: the column
$(0,e,2e)=(0,1,2)$ has a base word of only three letters, and
Proposition~\ref{M-prop:cornerpara} places the first two and the last two base edges in $\{a,c\}$,
which for a three-letter word is all of them; the word then carries no $b$-edge, contradicting
$y=1$. That argument uses no induction and no separation.

At $e\ge2$ the word has $3e\ge6$ letters, the four pinned positions leave $3e-4\ge e$ interior slots,
and the $b$-edges fit inside them: at $(2,5)$ the word $c\,c\,b\,b\,c\,c$ satisfies
$25+25+21+21+25+25=142=e(3f^2-e^2)$ and survives every filter we have. So Theorem~\ref{thm:n1} is
established at $e=1$ only, and its conclusion at $e\ge2$ --- on which
Corollary~\ref{cor:basewalls} depends --- must be regarded as \textbf{open}. We have not refuted the
conclusion; we refute only the proof. The missing ingredient is sharp: force the $\beta$-tile at
$V_k$, whose flanks are $\{a,c\}$, to lay its $c$ rather than its $a$ along $[V_k,V_{k+1}]$. This is
an instance of the crossing question of Remark~\ref{rem:straddle}, which
Remark~\ref{rem:chainobstruction} already identifies as the obstruction, and
Remark~\ref{rem:whyf3} already warns that a pinning argument must say against what the ends are
pinned. Here the would-be blocking edge at the upper end is $P_k$'s own $c$-edge, the very thing
being derived.
\end{remark}

\begin{remark}[The gap is exactly one crossing, and there is no cheaper repair]\label{rem:n1gapexact}
The step admits a sharp reduction. Let $Q$ be the tile across $T_k$'s $b$-edge
$[V_k,\,((k+1)c,0)]$, whose base end lies on $\partial ABC$ and whose end $V_k$ is interior. Measured
from the anchored base end, $Q$ lays $a$, $b$ or $c$ along that edge:
\begin{itemize}
\item $b$: an exact match, so $Q$ is the partner $P_k$ --- the case the induction wants;
\item $c$: overruns $V_k$ by $c-b=e^2$;
\item $a$: falls short of $V_k$ by $b-a$, leaving a run of that length still to be covered.
\end{itemize}
The third case does not close either, because $b-a$ is representable in $\langle a,b,c\rangle$ for no
member whatever. Indeed if $xa+yb+zc=b-a$ then $y=0$, since $y\ge1$ would give $yb\ge b>b-a$;
reducing $x\,ef+z\,f^2=f^2-e^2-ef$ modulo $f$ then gives $f\mid e^2$, hence $f=1$ by coprimality.
(\texttt{gap\_b\_sub\_a}; checked on all $29435$ coprime members with $e\le200$, $e<f<4e+3$, with no
representation in any of them.)

So both surviving branches force a tile edge past $V_k$, and the induction step of
Theorem~\ref{thm:n1} is \emph{equivalent} to the assertion that nothing overruns $V_k$. That is the
one-end-anchored pinning question in its purest form: the base end cannot be overrun, because the
continuation would leave the target, while $V_k$ is interior and, as computed in
Remark~\ref{rem:n1gap}, is not blocked by any tile already forced. Repairing Theorem~\ref{thm:n1} at
$e\ge2$ is therefore neither more nor less than settling that question, and the same question is what
stops the corner cascade of Remark~\ref{rem:route1uniform} and the filler dichotomy of
Section~\ref{sub:sidep}. Three routes, one wall.
\end{remark}

\begin{remark}[The overrun configuration at $V_k$, determined]\label{rem:overrunshape}
The crossing that has to be excluded is more rigid than it looks. By
Remark~\ref{rem:noedgetoedge}'s accounting, a tile whose edge passes through $V_k$ contributes
exactly $\pi$ there, so an overrun forces the tiles with a \emph{vertex} at $V_k$ to form a
$\pi$-figure. The forced ones are $T_k$, presenting $\gamma$, and $P_{k-1}$, presenting $\beta$, and
the only $\pi$-figure containing $\{\gamma,\beta\}$ is $\{\gamma,\alpha,\beta\}$. So an overrun
admits \emph{exactly one} further tile at $V_k$, and it presents $\alpha$.

That tile has room. The chord leaves the base at angle $\alpha$, so its extension at $V_k$ lies at
$\pi-\alpha$, while $P_{k-1}$'s $c$-edge lies along $\pi$; the wedge between them is exactly
$\alpha$. The configuration is therefore angularly consistent and cannot be refuted by counting
angles, which is why it has survived.

Its two edges at $V_k$ are $b$ and $c$, one along the chord's extension and one along $P_{k-1}$'s
$c$-edge. Along the latter, of length $c$:
\begin{itemize}
\item laying $c$ matches exactly;
\item laying $b$ falls short of $V_{k-1}$ by $c-b=e^2$, and $e^2$ lies in $\langle a,b,c\rangle$ for
no member (\texttt{gap\_e\_squared}), so that residue admits no whole-edge closure either.
\end{itemize}
The overrun therefore forces either a second crossing further along the same horizontal, or the
$\alpha$-tile to share $P_{k-1}$'s $c$-edge exactly and lay its $b$ along the chord's extension. We
record this because it reduces the open question to a single configuration, not because it settles
it: the second branch remains open, and a descent on the first is blocked by the fact that the
$e^2$-residue may be closed by an edge overrunning its \emph{other}, interior end rather than
$V_{k-1}$.

Where that surviving branch lands is worth recording, because it is not an accident. The chord's
direction vector is $a\,u-(c,0)$, which is $T_0$'s own $b$-edge, so
\[
W \;=\; V_k+\bigl(a\,u-(c,0)\bigr) \;=\; \bigl((k-1)c,\,0\bigr)+2a\,u .
\]
At $k=1$ this is $W=2a\,u$, a point of the equal side, and $R$ then lays an $a$-edge on that side
from $a\,u$ to $2a\,u$. Far from being absurd, that is exactly the configuration
Lemma~\ref{lem:ccorner} predicts for this column. So the overrun is geometrically coherent, not
merely angularly so, and it will not be refuted locally.
\end{remark}

\begin{remark}[Theorem~\ref{thm:n1} at $e\ge2$, repaired conditionally]\label{rem:n1conditional}
The previous remark also supplies the honest repair. Lemma~\ref{lem:ccorner} already observes that in
the column $(0,e,2e)$ the base carries no $a$-edge, so both corner tiles lay $c$ on the base, both
equal sides begin with an $a$-edge, and by Lemma~\ref{lem:sidequant} each carries at least $f$ of
them. In the notation of \S\ref{sub:sidep} that is $p\ge1$ on both sides. Hence

\begin{proposition}\label{prop:n1fromwalls}
Assume Hypothesis~\ref{hyp:walls} at the member $(e,f)$, i.e.\ $p=0$. Then no tiling has base walk
$(0,e,2e)$, for every $e\ge1$.
\end{proposition}

\begin{proof}
By Lemma~\ref{lem:ccorner} the column $(0,e,2e)$ forces $p\ge1$ on each equal side, contradicting
$p=0$.
\end{proof}

This is two lines and uses no induction, no chord pinning, and no crossing question. It is weaker
than Theorem~\ref{thm:n1} as printed, being conditional; but Hypothesis~\ref{hyp:walls} is assumed
throughout this programme anyway, so for every downstream use the conditional form is as good.
Corollary~\ref{cor:basewalls} and Corollary~\ref{cor:wallsf2e} may accordingly be read as holding
under Hypothesis~\ref{hyp:walls} at every member, and unconditionally at $e=1$, where
Remark~\ref{rem:n1gap}'s three-letter argument applies.

The moral is worth stating. Theorem~\ref{thm:n1}'s induction was attempting to prove, unconditionally
and by a long chord cascade, a statement that follows from Hypothesis~\ref{hyp:walls} immediately.
Its gap is therefore not an accident of that particular proof: excluding the overrun at $V_1$ means
exactly bounding the side's run of $a$-edges, which is Hypothesis~\ref{hyp:walls} again. The two are
one problem, and the cascade was the long way round.
\end{remark}

\begin{remark}[Counting cannot decide the crossing question]\label{rem:censusblind}
Since an overrun at $Y$ is equivalent to the tiles with a \emph{vertex} at $Y$ forming a
$\pi$-figure, it is natural to attack it with the vertex census of
Lemma~\ref{lem:census}, which counts exactly those figures. That cannot work, for a structural
reason worth recording so it is not attempted again.

Write $n_1$ for the number of $\{\gamma,\alpha,\beta\}$ figures and $n_2$ for the number of
$\{3\alpha,2\beta\}$ figures, and $v_1,\dots,v_4$ for the four interior types. The census consists of
one equation per angle, and its identity is obtained by subtracting the $\beta$ equation from the
$\alpha$ equation:
\[
v_1 \;=\; 1+n_2+v_3+2v_4 ,
\]
in which $n_1$ \emph{cancels}. The third equation then pins $n_1=N-3v_1-2v_2-v_3$, but supplies no
bound of its own. So the census constrains $\{3\alpha,2\beta\}$ and says nothing bounding about
$\{\gamma,\alpha,\beta\}$.

That is precisely the wrong way round for us: at $V_k$ the forced tiles are $T_k$, presenting
$\gamma$, and $P_{k-1}$, presenting $\beta$, and $\{\gamma,\alpha,\beta\}$ is the only $\pi$-figure
containing both. The census is therefore blind to exactly the configuration in question, and no
counting argument built on it can decide the crossing.

The kernel-checked $99$-tiling illustrates both points. Its vertex figures are: on the boundary,
$22$ of type $\{\gamma,\alpha,\beta\}$, with the corners contributing $\{\beta\}$ twice and
$\{3\alpha\}$ once; in the interior, $50$ of type $\{\gamma,\alpha,\beta\}$, $2$ of type
$\{3\alpha,2\beta\}$, and $(v_1,v_2,v_3,v_4)=(4,7,1,0)$. The identity reads $4=1+2+1+0$, exactly. And
the $50+2=52$ interior $\pi$-figures agree exactly with the $52$ edges carrying an interior vertex
found in Remark~\ref{rem:noedgetoedge} by an unrelated computation --- a useful cross-check on both.
Note the ratio: $50$ of the $52$ crossings are of the type the census cannot see.
\end{remark}

\begin{remark}[Anchoring one end is not enough, and the data says which edges get crossed]\label{rem:onlyanchor}
It is tempting to hope for a lemma of the shape ``a tile edge with one end on $\partial ABC$ is not
crossed''. No such lemma exists. Across the four kernel-checked tilings, the split edges with exactly
one endpoint on the boundary number $19$, $8$, $6$ and $5$ respectively, so crossings of that shape
are routine.

Sorting them by edge type is more informative. Writing crossed/total for edges with exactly one end
on the boundary:
\[
\begin{array}{lccc}
 & a & b & c\\
\text{99-tiling} & 0/27 & 4/12 & 15/22\\
\text{77-tiling} & 0/9 & 1/5 & 7/13\\
\text{44-tiling} & 0/10 & 1/13 & 5/9\\
\text{parallelogram 52} & 0/8 & 3/14 & 2/13
\end{array}
\]

Two things follow. First, $b$-edges with one anchored end \emph{are} crossed, in every one of the
four. The chord at $V_k$ in Theorem~\ref{thm:n1} is exactly such an edge, so the crossing there
cannot be excluded by any statement about anchoring alone: the exclusion has to consume the cascade's
forced tiles, and that is precisely what has resisted.

Second, $a$-edges with one anchored end are crossed in none of them. That one is not a coincidence,
and the next lemma explains it with a sharp constant.
\end{remark}

\begin{lemma}[Blocked-end quantization]\label{lem:anchorclear}
Let $[X,Y]$ be an edge of a tile of the tiling whose extension beyond $X$ is blocked --- either
because $X\in\partial ABC$, or because the continuation enters the interior of another tile
(Remark~\ref{rem:tib}). Then every vertex of the tiling lying in the relative interior of $[X,Y]$ is
at distance from $X$ belonging to the numerical semigroup $\langle a,b,c\rangle$. Consequently
$[X,Y]$ carries no interior vertex whenever
\[
\langle a,b,c\rangle\cap\bigl(0,\,|XY|\bigr)=\varnothing ,
\]
and in particular whenever $|XY|=\min(a,b,c)$.
\end{lemma}

\begin{proof}
Let $T$ be the tile having $[X,Y]$ as an edge. On $T$'s side of the segment nothing can contribute an
interior vertex, since $[X,Y]$ is a single edge of the triangle $T$. So any vertex in the relative
interior of $[X,Y]$ comes from the opposite side.

That side is covered by tiles abutting $[X,Y]$, each contributing a subsegment of one of its own
edges, and consecutive contributions meet at points of $[X,Y]$; a vertex in the relative interior is
exactly such a meeting point. By hypothesis no covering edge overruns $X$: its continuation would
leave $ABC$, or would enter the interior of a tile. Hence the first covering edge begins exactly at
$X$, and the covering proceeds by whole tile edges laid end to end. The meeting points are therefore
the partial sums of a sequence of tile side lengths, so each lies at distance from $X$ in
$\langle a,b,c\rangle$. (The last covering edge may overrun $Y$; that affects no meeting point
interior to $[X,Y]$.) The two consequences are immediate, the second because
$\langle a,b,c\rangle$ has no element strictly between $0$ and its smallest generator.
\end{proof}

\begin{corollary}[One-end-blocked chord dichotomy]\label{cor:onebloc}
Let $[X,Y]$ be a tile edge with the extension beyond $X$ blocked. Then the far side of $[X,Y]$ is
covered by whole tile edges laid end to end from $X$, and exactly one of the following holds:
\begin{enumerate}
\item[\textup{(i)}] they partition $[X,Y]$ exactly, so $|XY|$ is a sum of tile sides and the
arithmetic of Lemma~\ref{lem:cchord} applies as it stands;
\item[\textup{(ii)}] the last one overruns $Y$, and those before it sum to some
$s\in\langle a,b,c\rangle\cap\bigl(0,|XY|\bigr)$, with $s=0$ permitted.
\end{enumerate}
In particular, if $|XY|=\min(a,b,c)$ then that intersection is empty, so $s=0$ and the far side of
$[X,Y]$ is a \emph{single} tile edge, either matching $[X,Y]$ exactly or overrunning $Y$.
\end{corollary}

\begin{proof}
Immediate from Lemma~\ref{lem:anchorclear}: the covering begins at $X$ and proceeds by whole edges,
so the partial sums before the final edge lie in $\langle a,b,c\rangle$, and the final edge either
ends at $Y$ or passes it.
\end{proof}

This is the common form of three arguments that were being made separately. Lemma~\ref{lem:cchord} is
the case where \emph{both} ends are blocked, which deletes branch (ii) outright. The chord at $V_k$
in Remark~\ref{rem:n1gapexact} is the case $|XY|=b$ with the base end blocked, where
$\langle a,b,c\rangle\cap(0,b)$ consists of the multiples of $a$ below $b$ --- which is exactly the
``$j$ $a$-edges, then an overrun'' description obtained there by hand. And the filler dichotomy of
\S\ref{sub:sidep} is the same statement at a side junction. Branch (ii) is the crossing that
Remark~\ref{rem:onlyanchor} shows cannot be removed by anchoring alone; the corollary does not close
it, it states it once instead of three times.

\begin{remark}[Branch \textup{(ii)} occurs, so no general lemma will delete it]\label{rem:branchii}
It is worth being explicit that branch (ii) is not a formal possibility awaiting exclusion: it is
realised in a kernel-checked tiling. In the $99$-tiling, whose tile scaled by $16$ is
$(a,b,c)=(32,48,64)$, the $b$-edge from $(88,24)$ --- a boundary point --- to the interior point
$(136,24)$ carries one interior vertex, at distance $32$ from the blocked end. The remaining tail is
$48-32=16$, and $16\notin\langle32,48,64\rangle$, so the covering edge that meets $Y$ must overrun
it. Four $b$-edges of that shape occur in that tiling alone.

Together with Remark~\ref{rem:onlyanchor} this settles the shape of the obstruction. Three successive
generalisations have now been refuted by explicit certified witnesses: that an edge with one end on
the boundary is uncrossed (false, $19+8+6+5$ counterexamples); that anchoring plus a length condition
suffices (false, Remark~\ref{rem:noedgetoedge}); and that branch (ii) can be deleted (false, above).
What survives is Corollary~\ref{cor:onebloc} itself, which is sharp, and the conclusion that the
crossing at a given chord can only be excluded using the tiles the cascade has already forced there.
No lemma quantified over chords will do it. That is a constraint on how the remaining work must be
organised, and it is the reason the open steps in this note are stated as facts about particular
configurations rather than as general pinning principles.
\end{remark}

\begin{remark}[The quantization is exactly tight]\label{lem:anchorclearsharp}
The lemma is sharp in the strongest sense: not only does every crossing sit at a semigroup position,
but \emph{every} semigroup position below the edge length is realised. Across the eight
kernel-checked tilings, all $52$ crossings on edges with exactly one anchored end lie at a distance
in $\langle a,b,c\rangle$ from that end, with no exception; and edges of minimal length are crossed
$0$ times out of $74$. Comparing permitted with observed positions, by edge length:
\[
\begin{array}{llll}
\text{$44$-, $77$-, $99$-tilings} & (a,b,c)=(16,24,32)\ \text{or}\ (32,48,64)\\
\qquad\text{length }a & \text{permitted } \varnothing & \text{observed } \varnothing\\
\qquad\text{length }b & \text{permitted }\{a\} & \text{observed }\{a\}\\
\qquad\text{length }c & \text{permitted }\{a,b\} & \text{observed }\{a,b\}\\[2pt]
\text{parallelogram }52 & (a,b,c)=(18,48,54)\\
\qquad\text{length }a & \text{permitted } \varnothing & \text{observed } \varnothing\\
\qquad\text{length }b & \text{permitted }\{18,36\} & \text{observed }\{18,36\}\\
\qquad\text{length }c & \text{permitted }\{18,36,48\} & \text{observed }\{18,36,48\}
\end{array}
\]
So no strengthening of the conclusion is possible: the permitted set is attained.

Which side is minimal depends on the member: $a<b$ exactly when $f>\varphi e$ with $\varphi$ the
golden ratio, since $a<b$ reads $f^2-ef-e^2>0$. At the members carrying the certified tilings ---
$(1,2)$ and $(1,3)$ --- one has $a=m$, which is why the observation appeared as a statement about
$a$-edges.

The lemma is elementary and geometric, and its arithmetic content is only that a tile side is at
least the least tile side, so we do not formalise it; it carries the label PROVED. It does not touch
the chord at $V_k$, whose length is $b$ and which is crossed in practice
(Remark~\ref{rem:onlyanchor}). It is recorded because it is the first general statement we have that
excludes crossings outright, and because it fixes the exact radius within which the boundary still
controls the interior.
\end{remark}

\begin{corollary}[The walls form at $e=1$, $f\ge3$]\label{cor:wallsf2e}
Let $e=1$ and $f\ge3$. Then the base walk is the walls form $(f,1,1)$.
\end{corollary}

\begin{proof}
Corollary~\ref{cor:basedi2e} applies, since $f\ge3>2=2e$, and leaves $(0,1,2)$ and $(f,1,1)$; the
first is excluded by the three-letter argument of Remark~\ref{rem:n1gap}.
\end{proof}

\subsection{The side parameter, and where its range is bounded}\label{sub:sidep}

Once the base walk is pinned, what remains is the side. The side question has a single integer
parameter, and on the family of Corollary~\ref{cor:wallsf2e} that parameter takes only two open
values, uniformly in the member.

The parametrisation itself is already in hand and is not new here. By
Lemma~\ref{lem:sidequant} an equal side carries no $b$-edge at any member and its $a$-count is a
multiple of $f$; writing that count as $pf$, the side equation $pf\cdot ef+n_cf^2=f^3$ gives
\[
n_c+pe=f ,
\]
and the $\gamma$-trap's $n_c\ge1$ gives $0\le p\le(f-1)/e$. Hypothesis~\ref{hyp:walls} is exactly
the assertion $p=0$. Lemma~\ref{lem:ccorner} already records this profile in the form ``$a$-count
$fj$ with $je\le f-1$'', and the Lean statements are
\texttt{SideNoB.side\_no\_b\_uncond}, \texttt{SideNoB.side\_quantized} and
\texttt{SideNoB.c\_corner\_forces\_side\_a}. What follows is the one consequence we did not find
recorded: on the close pairs with $f>2e$ the range of $p$ is not merely bounded but equal to
$\{0,1,2\}$, uniformly in the member.

\begin{corollary}[A uniform two-case side problem]\label{cor:ptwo}
In the setting above, suppose in addition that $2e<f$ and that $(e,f)$ is a close pair,
$f^2\le2ef+e^2$. Then
$\lfloor(f-1)/e\rfloor=2$. Hence on that family Hypothesis~\ref{hyp:walls} reduces to excluding
$p\in\{1,2\}$ --- the same two cases for every member.
\textup{(\texttt{side\_p\_le\_two}.)}
\end{corollary}

\begin{proof}
From $f\ge2e+1$ we get $(f-1)/e\ge2$. Conversely $p\ge3$ would give $3e\le pe\le f-1$, so
$f\ge3e+1$ and then $f^2-2ef-e^2\ge(3e+1)(e+1)-e^2=2e^2+4e+1>0$, contradicting the close-pair
inequality.
\end{proof}

\begin{remark}[What this buys, and what it does not]\label{rem:ptwoscope}
Corollary~\ref{cor:ptwo} is independent of Theorem~\ref{thm:n1} and of the base walk: it
constrains the side at any member. On the close pairs
with $f>2e$ --- $203$ of them with $e\le40$, the smallest being $(3,7)$ with $N=138$, then $(4,9)$
with $N=227$ and $(5,11)$ with $N=338$ --- the side parameter is bounded by $2$. Checked on all
$126003$ close pairs with $f>2e$ and $e\le1000$: $\lfloor(f-1)/e\rfloor=2$ in every one.

On that family Corollary~\ref{cor:basedi2e} also reduces the base to two columns. It does not reduce
it to one: that step needs Theorem~\ref{thm:n1} at $e\ge2$, which Remark~\ref{rem:n1gap} leaves open.
So for these members the current state is two base columns and two side parameters, not one of each.

The contrast with $e=1$ is the point. There the range is $1\le p\le f-1$ and grows without bound
with $f$, which is why the corner cascade of Conjecture~\ref{conj:advance} has to be run as an
induction over block lengths. Here the range is $\{1,2\}$ for every member of an infinite family, so
the side question is a fixed finite case analysis rather than an induction. This is the first
sub-family in which that happens.

At $(3,7)$, concretely: $p=1$ gives a side of $7$ $a$-edges and $4$ $c$-edges, and $p=2$ gives $14$
$a$-edges and $1$ $c$-edge. Excluding those two remains open; the reduction bounds the problem, it
does not solve it.
\end{remark}

\subsection{A parity of the straight figures}\label{sub:parity}

Lemma~\ref{lem:census} records three counting identities whose published consequence is
$v_1=1+n_2+v_3+2v_4$. The $\alpha$-identity also carries a parity statement, which we record together
with the reason it does \emph{not} exclude any member.

\begin{lemma}[Straight-figure parity]\label{lem:parity}
Write $S=n_1+n_2$ for the total number of straight ($\pi$) figures. Then $S\equiv N+1\pmod 2$.
\end{lemma}

\begin{proof}
The $\alpha$-identity is $3+n_1+3n_2+2v_2+4v_3+6v_4=N$; substituting $n_1=S-n_2$ gives
$S+2(n_2+v_2+2v_3+3v_4)=N-3$. \textup{(}\texttt{census\_parity}.\textup{)}
\end{proof}

The three identities have rank two and this is the only parity they contain: substituting the
$v_1$ relation into the $\beta$-identity returns the $\alpha$-identity exactly.

\begin{remark}[Why this does not exclude a member, and the trap it sets]\label{rem:parityscope}
It is tempting to combine Lemma~\ref{lem:parity} with the boundary walks. On a member all of whose
walks have odd edge-counts --- $(3,7)$ is one, with $k=4p+7$ on a side and $9$ or $13$ on the base
--- the boundary contributes $\sum k_i-3$ straight figures, an even number, against the odd value
$N+1=139$ that Lemma~\ref{lem:parity} demands. That looks like an exclusion. It is not.

$S$ counts \emph{every} straight figure, and a straight figure need not lie on the boundary: at an
interior $T$-junction one tile has the vertex in the relative interior of a side, contributing a
straight angle and no corner, while the tiles on the other side have corners summing to $\pi$. Such a
vertex is a $\{\gamma,\alpha,\beta\}$ or $\{3\alpha,2\beta\}$ figure and is counted in $n_1+n_2$.
Writing $T$ for their number, $S=\sum k_i-3+T$, so the parity constrains only $T$: at $(3,7)$ every
tiling has an \emph{odd} number of interior $T$-junctions. Since non-edge-to-edge incidences are
exactly what this problem must allow, nothing is excluded.

This is the failure mode Lemma~\ref{lem:census} already warns of --- corner counting cannot see the
side structure --- and the boundary walks are side structure. It is recorded here because the
temptation is sharp and the arithmetic is otherwise correct.
\end{remark}


\subsection{Apex rigidity, and a bound on the side parameter}\label{sub:apex}

The apex of a base-$\beta$ target is completely rigid, and the rigidity is an identity of the whole
family rather than a coincidence at one member. It forces the covering of the first chord exactly, and
from that it bounds $p$.

Throughout, $A$ is the apex, the equal sides have length $f^3$, the apex angle is $3\alpha$, and
$\sin(\alpha/2)=e/(2f)$, whence
\[
\cos\tfrac{\alpha}{2}=\frac{\sqrt{4f^{2}-e^{2}}}{2f},\qquad
\cos\tfrac{3\alpha}{2}=\frac{(f^{2}-e^{2})\sqrt{4f^{2}-e^{2}}}{2f^{3}},\qquad
\sin\tfrac{3\alpha}{2}=\frac{e(3f^{2}-e^{2})}{2f^{3}} .
\]
By Lemma~\ref{lem:sidenob} an equal side carries no $b$-edge, and the apex figure is $\{3\alpha\}$, so
the tile at $A$ presents $\alpha$ and its two edges there are $b$ and $c$; the one on the side is
therefore a $c$. Hence the first junction $J$ of the side sits at distance $c=f^{2}$ from $A$. Write
$C$ for the chord through $J$ parallel to the base. It cuts off a triangle similar to the target with
ratio $f^{2}/f^{3}=1/f$, so the area above $C$ is $N/f^{2}$ tile areas and $|C|=e(3f^{2}-e^{2})/f$.

\begin{lemma}[The apex identities]\label{lem:apexid}
For every $(e,f)$,
\[
b\cos\tfrac{\alpha}{2}=c\cos\tfrac{3\alpha}{2},
\qquad
c\sin\tfrac{3\alpha}{2}-b\sin\tfrac{\alpha}{2}=a .
\]
\end{lemma}

\begin{proof}
Both sides of the first are $(f^{2}-e^{2})\sqrt{4f^{2}-e^{2}}/(2f)$. For the second,
$c\sin\frac{3\alpha}{2}=e(3f^{2}-e^{2})/(2f)$ and $b\sin\frac{\alpha}{2}=e(f^{2}-e^{2})/(2f)$, whose
difference is $e\cdot 2f^{2}/(2f)=ef=a$. Both are machine-checked
(\texttt{lean/ApexRigidity.lean}: \texttt{apex\_drop\_eq}, \texttt{apex\_edge\_eq}).
\end{proof}

\begin{theorem}[The first chord is covered exactly]\label{thm:apexconfig}
Let $T_1,T_2,T_3$ be the three tiles at $A$, with $T_1,T_3$ the outer ones. Then:
\begin{enumerate}
\item[(a)] $T_1$'s remaining apex edge is its $b$, ending at a point $P$ lying \emph{exactly} on $C$;
\item[(b)] $T_1$'s third edge $JP$ lies \emph{along} $C$ and has length exactly $a$;
\item[(c)] $T_2$ has exactly the fraction $b/c$ of its area above $C$, and its trace on $C$ is the
segment $PP'$ joining $P$ to the mirror point $P'$ of $T_3$, of length $e(f^{2}-e^{2})/f$;
\item[(d)] consequently $a+|PP'|+a=|C|$ and the area above $C$ is $2+b/c=N/f^{2}$ --- so $T_1$, $T_3$
and that fraction of $T_2$ fill the region above $C$ \emph{exactly}, and no other tile has any area
above $C$.
\end{enumerate}
\end{theorem}

\begin{proof}
(a) The drop of $J$ below $A$ is $c\cos\frac{3\alpha}{2}$ and that of $P$ is $b\cos\frac{\alpha}{2}$;
they agree by Lemma~\ref{lem:apexid}. (b) The horizontal offsets of $J$ and $P$ from the axis of
symmetry are $c\sin\frac{3\alpha}{2}$ and $b\sin\frac{\alpha}{2}$, differing by $a$.

(c) Note first that $P$ and $P'$ are defined by $T_1$ and $T_3$, at distance $b$ along the two inner
rays from $A$, independently of $T_2$. Now $T_2$ presents $\alpha$ at $A$, so its two edges there are a
$b$ and a $c$, laid on those same two rays --- and \emph{both orientations occur}, so we treat them
together. Its $b$-edge terminates at drop $b\cos\frac{\alpha}{2}$, which is the chord by
Lemma~\ref{lem:apexid}, hence at whichever of $P,P'$ lies on that ray. Its $c$-edge terminates at drop
$c\cos\frac{\alpha}{2}$, strictly below $C$, and the chord meets it at distance $b$ from $A$, since
$\cos\frac{3\alpha}{2}/\cos\frac{\alpha}{2}=b/c$ --- that is, at the other of $P,P'$. So in either
orientation $T_2\cap C=PP'$, and $T_2$ is cut by a line through one of its vertices, giving area ratio
$b/c$ (\texttt{middle\_fraction}). The length of $PP'$ is $2b\sin\frac{\alpha}{2}=e(f^{2}-e^{2})/f$.

(d) $2ef+e(f^{2}-e^{2})/f=e(3f^{2}-e^{2})/f=|C|$, and $2+b/c=(3f^{2}-e^{2})/f^{2}=N/f^{2}$
(\texttt{area\_above\_chord}), which is the area above $C$. Since the three tiles already supply all of
it, nothing else can.
\end{proof}

\begin{corollary}[The tile below $T_1$'s $a$-edge has a vertex at $P$]\label{cor:noTP}
$T_1$ presents $\gamma$ at $P$, and either $T_2$ presents $\gamma$ at $P$ as well, or $P$ lies interior
to an edge of $T_2$. In both cases no \emph{further} tile has $P$ interior to an edge; in particular
the tile below $T_1$'s $a$-edge has a vertex at $P$, and the residual figure there is
$\{2\alpha,2\beta\}$, $\{\beta,\gamma\}$ or $\{\alpha,\beta\}$.
\end{corollary}

\begin{proof}
$T_1$ has sides $c,b,a$ at $A,P,J$ by Theorem~\ref{thm:apexconfig}, so its angle at $P$ is opposite
$c$, namely $\gamma$. By Theorem~\ref{thm:apexconfig}(c) the trace of $T_2$ on $C$ is $PP'$, so $P$
lies on $\partial T_2$: it is either a vertex of $T_2$ or interior to one of its edges.

Suppose it is a vertex. The edge of $T_2$ at $P$ along the ray towards $A$ terminates at $P$, at
distance $b$ from $A$; $T_2$'s flanks at $A$ are $b$ and $c$, and a $c$ there would terminate at
distance $c\ne b$, so that edge is $T_2$'s $b$. Hence $T_2$'s angle at $P$ has $b$ as a flank and is
$\alpha$ or $\gamma$. If it were $\alpha$, the other flank at $P$ would be $c$ and $T_2$ would be the
triangle $APR$ with $AP=b$, $PR=c$, $AR=a$, whose angle at $A$ is opposite $c$, that is $\gamma$ ---
contradicting that $T_2$ presents $\alpha$ at $A$. So $T_2$ presents $\gamma$.

For the count: two $\gamma$'s total $\pi+\alpha$, leaving $\pi-\alpha<\pi$; a $\gamma$ and a straight
angle total $\pi+\gamma$, leaving $\pi-\gamma<\pi$. In neither case does a further straight angle fit.
Finally the tile below $T_1$'s $a$-edge is neither $T_1$ nor $T_2$, since at $P$ those occupy the
wedge towards $A$ and the wedge towards $P'$ respectively, so it has a vertex at $P$.
\end{proof}

\begin{remark}[The two equal sides differ]\label{rem:twosides}
The dichotomy in Corollary~\ref{cor:noTP} is not an artefact. $T_2$ lays $b$ on one of the two rays it
shares with $T_1$ and $T_3$ and $c$ on the other, so exactly one of $P$, $P'$ is a genuine vertex
point and the other carries a $T$-junction. This is the asymmetry between the two equal sides already
visible in Remark~\ref{rem:dichotomy}. It does not affect what follows: Corollary~\ref{cor:noTP}
delivers the same conclusion on both sides.
\end{remark}

\begin{theorem}[The second edge of an equal side is a $c$]\label{thm:secondc}
Assume $e^{2}+ef<f^{2}$, equivalently $a<b$. Then on every equal side the tile edge immediately below
the first junction $J$ is a $c$.
\end{theorem}

\begin{proof}
Let $Y$ be the tile below $T_1$'s $a$-edge, adjacent to it at $P$. By Corollary~\ref{cor:noTP}, $Y$ has
a vertex at $P$ and an edge running from $P$ along $C$ towards $J$. Since $J$ lies on the boundary at
distance $a$, that edge has length at most $a$; as $a<b<c$ it is exactly $a$, so $Y$'s $a$-edge is
$JP$, matching $T_1$'s. In particular $Y$ does not present $\alpha$ at $P$, whose flanks are $b$ and
$c$.

At $J$ the figure is straight, hence $\{\alpha,\beta,\gamma\}$ or $\{3\alpha,2\beta\}$. $T_1$ presents
$\beta$ there (its angle opposite $b$). The angles at the two ends of $Y$'s $a$-edge are $\beta$ and
$\gamma$, so $Y$ presents $\gamma$ or $\beta$ at $J$. In the first case the figure must be
$\{\alpha,\beta,\gamma\}$ and the third tile presents $\alpha$; in the second it must be
$\{3\alpha,2\beta\}$ and the remaining three tiles all present $\alpha$. Either way, $T_1$ occupies the
extreme position against the ascending side and $Y$ the next, so the tile against the \emph{descending}
side presents $\alpha$. Its flanks are $b$ and $c$, and the side edge below $J$ is that tile's edge; by
Lemma~\ref{lem:sidenob} it is not a $b$, hence it is a $c$.
\end{proof}

\begin{corollary}[A bound on the side parameter]\label{cor:pbound}
If $e^{2}+ef<f^{2}$ then $n_c\ge2$, hence
\[
pe+2\le f .
\]
In particular $p=2$ is impossible whenever $2e+2\le f$ fails, and on the tight subfamily $f=2e+1$ the
bound reads $pe\le 2e-1$, so $p\le1$.
\end{corollary}

\begin{proof}
The edge at the apex and the edge below $J$ are both $c$ by Theorem~\ref{thm:secondc}, so $n_c\ge2$;
with $n_c+pe=f$ this is $pe+2\le f$. For $f=2e+1$, $p\ge2$ would give $2e\le pe\le2e-1$
(\texttt{side\_p\_bound}, \texttt{p\_le\_one\_of\_tight}).
\end{proof}

\begin{remark}[What this settles, and exactly where it stops]\label{rem:apexscope}
Corollary~\ref{cor:pbound} is a second and independent proof of Theorem~\ref{thm:ptwodead} on the
tight subfamily, by a route that uses no chord bound and no height identity --- only the apex
identities and the boundary at $J$. It also completes the first bullet of
Remark~\ref{rem:dichotomy}: there the $a\,c$ ending was excluded only in the figure
$\{\gamma,\alpha,\beta\}$ and only with the $\alpha$ adjacent to $C$, whereas
Theorem~\ref{thm:secondc} excludes it in \emph{both} straight figures and for every $p$. Hence at
$(3,7)$ the side word of an equal side must end $c\,c$; the two cases $X=\beta$ and $X=\gamma$ of
Remark~\ref{rem:dichotomy} both persist, but now only inside that ending.

It does not reach $p=1$ at any new member, and the reason is arithmetic rather than incidental.
Killing $p=1$ by this bound needs $e\ge f-1$; with $e^{2}+ef<f^{2}$ that forces $f=e+1$ and then
$e^{2}<e+1$, so $e=1$ and $f=2$ --- the member already closed by
Theorem~\ref{thm:basebeta-e1} (\texttt{scope\_limitation}). The hypothesis $e^{2}+ef<f^{2}$ says
$f/e$ exceeds the golden ratio, and $e\ge f-1$ says it is close to $1$; the two cannot hold together
beyond $(1,2)$.

The obstruction to iterating is equally precise. Corollary~\ref{cor:noTP} needed \emph{two} tiles
presenting $\gamma$ at $P$, and the second, $T_2$, exists only because the apex figure is $3\alpha$.
At the analogous point one level down --- the far end of the next $c$-tile's $a$-edge, which does lie
exactly on the next chord, by the same identities --- only one $\gamma$ is forced. Supplying a second
is what an induction would require, and we do not have it.
\end{remark}

\begin{proposition}[The descent is self-similar]\label{prop:selfsim}
Suppose a tile $Z$ lays a $c$-edge of an equal side between consecutive junctions $J'$ (nearer the
apex) and $J$, and presents $\alpha$ at $J'$. Let $W$ be its third vertex. Then $W$ lies \emph{exactly}
on the chord $C$ through $J$; the segment $JW$ is $Z$'s $a$-edge, of length $a$, lying along $C$; and
$Z$ presents $\beta$ at $J$ and $\gamma$ at $W$.
\end{proposition}

\begin{proof}
The two edges of $Z$ at $J'$ are its flanks at $\alpha$, namely $b$ and $c$; the $c$ runs down the side
to $J$ and the $b$ into the interior, at angle $\alpha$ from the descending side, hence at
$\frac{3\alpha}{2}-\alpha=\frac{\alpha}{2}$ from the downward vertical --- the same inclination as the
apex tile's $b$-edge in Theorem~\ref{thm:apexconfig}. So the drop from $J'$ to $W$ is
$b\cos\frac{\alpha}{2}$ and that from $J'$ to $J$ is $c\cos\frac{3\alpha}{2}$, equal by
Lemma~\ref{lem:apexid}; and the horizontal offsets differ by
$c\sin\frac{3\alpha}{2}-b\sin\frac{\alpha}{2}=a$. The angles at $J$ and $W$ are those opposite $b$ and
$c$, namely $\beta$ and $\gamma$.
\end{proof}

\begin{remark}[Why the descent does not close, and it is not for want of trying]\label{rem:selfsim}
Proposition~\ref{prop:selfsim} says the configuration at $J$ is a verbatim copy of the configuration
at the first junction, with $(A,J_1,P)$ replaced by $(J',J,W)$. If the argument of
Theorem~\ref{thm:secondc} could be run at $J$ as it was at $J_1$, it would force the side edge below
$J$ to be a $c$ as well, and induction down the side would give $n_a=0$, that is $p=0$ --- exactly
Hypothesis~\ref{hyp:walls}. It cannot, and the obstruction is sharp.

The argument at $J_1$ turns on Corollary~\ref{cor:noTP}: \emph{two} tiles present $\gamma$ at $P$, so
the residual is $\pi-\alpha<\pi$ and no tile can have $P$ interior to an edge. The second $\gamma$
comes from $T_2$, which exists only because the apex figure is $\{3\alpha\}$. At $W$ only $Z$'s
$\gamma$ is forced, and no counting argument can supply the other:

\begin{lemma}[A single $\gamma$ never excludes a $T$-junction]\label{lem:onegamma}
Let $x$ be an interior point at which one tile presents $\gamma$ and one tile has $x$ interior to an
edge. Then the remaining tiles at $x$ present angles totalling exactly $\alpha+\beta$, realised by one
$\alpha$ and one $\beta$. The angle count is therefore always consistent.
\end{lemma}

\begin{proof}
$2\pi=\gamma+\pi+R$ gives $R=\pi-\gamma=\alpha+\beta$, since $\alpha+\beta+\gamma=\pi$.
\end{proof}

So the obstruction is not a missing computation but the tile's own angle sum: a lone $\gamma$ beside a
straight edge balances at every member and at every point. Exclusion by angle count needs two
$\gamma$'s, and the apex is the only place that supplies two.

We also record a near miss, since it shows how narrow the gap is. At $W$ the $T$-junction \emph{can}
be excluded when it would sit on the tile $V$ angularly adjacent to $Z$ at $J'$: if $V$ lays its $c$
along $Z$'s $b$-ray, the budget is $\gamma+\pi$ with residual $\pi-\gamma<\pi$, so no \emph{second}
straight angle fits, and the tile below $Z$'s $a$-edge is forced to have a vertex at $W$ and to lay
exactly $a$ --- which is all Theorem~\ref{thm:secondc} needs. What defeats this is that the
$T$-junction need not sit on $V$ at all: the tile \emph{below} $C$ may carry an edge running along $C$
from $J$ through $W$, of length $b$ or $c$, with its vertex on the boundary at $J$. That is precisely
the tile the argument was trying to pin, so it closes on itself. Settling whether that placement can
occur is, as far as we can see, the whole of what remains at $p=1$.
\end{remark}

\begin{definition}[The points $W(J)$]\label{def:WJ}
Let $J$ be a junction of an equal side whose edge \emph{above} it is a $c$, laid by a tile presenting
$\alpha$ at the junction above. By Proposition~\ref{prop:selfsim} that tile's third vertex lies on the
chord through $J$ at distance exactly $a$ from $J$; call it $W(J)$.
\end{definition}

\begin{theorem}[An $a$-edge forces a $T$-junction]\label{thm:aforcesT}
Assume $e^{2}+ef<f^{2}$. Let $J$ be as in Definition~\ref{def:WJ} and suppose the side edge
\emph{below} $J$ is an $a$. Then $W(J)$ lies in the \emph{interior} of an edge of some tile.
\end{theorem}

\begin{proof}
Every tile meeting $J$ has a vertex there. A tile with $J$ interior to an edge would contribute $\pi$
and so exhaust the straight figure by itself, forcing the two side edges at $J$ to be collinear parts of
a single tile edge, in which case $J$ would not be a junction; and an edge through $J$ transversal to the
side would carry that tile outside the target.

Write $Z$ for the $c$-tile above $J$; by Proposition~\ref{prop:selfsim} it presents $\beta$ at $J$ and
its $a$-edge is $JW$ with $W=W(J)$. Let $M$ be the tile below sharing the chord ray with $Z$, and let
$X$ be the last tile in angular order, standing against the descending side; $X$ lays the next side
edge on one of its flanks, and that flank is not a $b$ by Lemma~\ref{lem:sidenob}.

The figure at $J$ is $\{\alpha,\beta,\gamma\}$ or $\{3\alpha,2\beta\}$, with $Z$ first. Exhausting the
orders with first entry $\beta$:
\[
\begin{array}{llll}
\text{figure} & \text{order} & M & \text{side edge below }J\\\hline
\{\alpha,\beta,\gamma\} & (\beta,\alpha,\gamma) & \alpha & a\\
\{\alpha,\beta,\gamma\} & (\beta,\gamma,\alpha) & \gamma & c\\
\{3\alpha,2\beta\} & (\beta,\alpha,\alpha,\alpha,\beta) & \alpha & a \text{ or } c\\
\{3\alpha,2\beta\} & (\beta,\alpha,\alpha,\beta,\alpha) & \alpha & c\\
\{3\alpha,2\beta\} & (\beta,\alpha,\beta,\alpha,\alpha) & \alpha & c\\
\{3\alpha,2\beta\} & (\beta,\beta,\alpha,\alpha,\alpha) & \beta & c
\end{array}
\]
(the side edge is read off $X$'s flanks with $b$ removed: $\alpha\mapsto c$, $\beta\mapsto a$ or $c$,
$\gamma\mapsto a$). In both rows permitting an $a$, $M$ presents $\alpha$. Its flanks are then $b$ and
$c$, so the edge $M$ lays along the chord ray has length $b$ or $c$, and $e^{2}+ef<f^{2}$ gives
$a<b<c$. Since $|JW|=a$, the point $W$ lies strictly inside that edge.
\end{proof}

\begin{corollary}[A sufficient condition for Hypothesis~\ref{hyp:walls}]\label{cor:walls-from-T}
Assume $e^{2}+ef<f^{2}$. If for every junction $J$ as in Definition~\ref{def:WJ} the point $W(J)$ is a
\emph{vertex} of every tile meeting it, then $p=0$ on every equal side, i.e.\
Hypothesis~\ref{hyp:walls} holds at that member.
\end{corollary}

\begin{proof}
Suppose some equal side carries an $a$-edge, and take the first one counted from the apex. Every edge
above it is a $c$, so the junction $J$ immediately above it satisfies Definition~\ref{def:WJ}, and
Theorem~\ref{thm:aforcesT} puts $W(J)$ interior to an edge --- contrary to hypothesis. Hence
$n_a=0$, that is $p=0$.
\end{proof}

\begin{remark}[The crux, in its sharpest form]\label{rem:cruxsharp}
Corollary~\ref{cor:walls-from-T} replaces ``exclude the surviving configuration'' by a single
question about an explicitly located finite set of points. The exclusion succeeds at the first
junction: there $W(J_1)=P$, and Corollary~\ref{cor:noTP} rules out a $T$-junction because two tiles
present $\gamma$ at $P$. It is unavailable at every later junction, and not by accident:
Lemma~\ref{lem:onegamma} shows a single $\gamma$ beside a straight edge always balances, so no angle
count can do it, and Proposition~\ref{prop:selfsim} shows the configuration one level down is a
verbatim copy, so no new local information appears. What is needed is a second $\gamma$ at $W(J)$ from
geometry rather than from counting.

We note the numerical consequence is weak and should not be oversold: only a $c\to a$ transition
forces a $T$-junction, and a side word such as $c\,c\,c\,c\,a^{7}$ has exactly one. The theorem gives
one forced $T$-junction per equal side carrying an $a$, not one per $a$-edge, and in particular it
does not revive the counting argument retracted in Remark~\ref{rem:parityscope}.
\end{remark}



\subsection{The tight subfamily \texorpdfstring{$f=2e+1$}{f=2e+1}}\label{sub:tight}

Corollary~\ref{cor:ptwo} leaves the two cases $p\in\{1,2\}$. The case $p=2$ is not uniformly hard:
on one infinite subfamily it is completely rigid, and the rigidity is visible in a counting budget.

\begin{lemma}[The $\gamma$-injection budget at $p=2$]\label{lem:tight}
A side of parameter $p$ carries $n_a=pf$, $n_b=0$, $n_c=f-pe$, hence $k=pf+(f-pe)$ edges and $k-1$
junctions, so the $\gamma$-injection bound $\#a+\#b\le k-1$ has slack $f-pe-1$. At $p=1$ the slack is
$f-e-1>0$ for every member. At $p=2$ it is $f-2e-1$, which vanishes exactly when $f=2e+1$.
\textup{(}\texttt{gamma\_slack}, \texttt{p\_two\_tight\_iff}.\textup{)}
\end{lemma}

The subfamily $f=2e+1$ is infinite and automatically coprime, with $N=3f^2-e^2=11e^2+12e+3$: the
members $(1,3)$ at $N=26$, $(2,5)$ at $71$, $(3,7)$ at $138$, $(4,9)$ at $227$, $(5,11)$ at $338$.
The first is closed by Theorem~\ref{thm:walls13} and the second by search, so $(3,7)$ is the smallest
tight member not already settled.

\begin{proposition}[The $p=2$ side is forced]\label{prop:tightside}
Let $f=2e+1$ and suppose an equal side has $p=2$. Then $n_c=f-2e=1$, and the side word is
$a^{2f}c$, read from the base corner to the apex. Every one of its $2f$ junctions carries exactly one
$\gamma$ and is the figure $\{\gamma,\alpha,\beta\}$; each $a$-tile presents $\beta$ at its left
junction and $\gamma$ at its right; the corner tile lays $a$ on the side and $c$ on the base.
\end{proposition}

\begin{proof}
$n_c=f-2e=1$ and $2f\cdot ef+f^2=f^2(2e+1)=f^3$, so the profile is consistent
(\texttt{p\_two\_single\_c}). Tightness makes the $\gamma$-injection a bijection, so every junction
carries exactly one $\gamma$; a $\pi$-vertex carrying a $\gamma$ is $\{\gamma,\alpha,\beta\}$
(\texttt{pi\_vertex\_with\_gamma}). The apex figure is $3\alpha$ and the sides carry no $b$-edge
(Lemma~\ref{lem:sidenob}), so the apex-end edge is a $c$; with only one $c$ available it sits there,
fixing the word. The corner figure is a single $\beta$ with flanks $\{a,c\}$, and the base-corner end
being an $a$-edge puts the $c$ on the base. The alternation follows: the corner tile presents
$\gamma$ at $J_1$, and each junction having a unique $\gamma$ forces the next $a$-tile to present
$\beta$ there.
\end{proof}

\begin{remark}[The last junction, and what is still missing]\label{rem:lastjunction}
At $J_{2f}$ the two boundary neighbours are the $a$-tile $A_{2f}$, presenting $\gamma$ with a
$b$-ray inward, and the apex $c$-tile, presenting $\beta$ with an \emph{$a$}-ray inward; the filler's
rays are $b$ and $c$. Since neither $b$ nor $c$ equals $a$, \emph{$J_{2f}$ cannot be edge-to-edge},
and each placement leaves a residual run: $c-a$, or $b-a$ together with $c-b=e^2$. All three are gaps
of $\langle a,b,c\rangle$ --- $b-a$ by \texttt{gap\_b\_sub\_a}, $e^2$ by \texttt{gap\_e\_squared},
and $c-a$ by \texttt{gap\_c\_sub\_a\_tight}, which holds precisely for $e\ge2$ since $c-a=2a$ at
$(1,3)$. Moreover $e^2<a$, so no tile edge fits inside that run at all. Hence the covering
\emph{must} overrun a filler vertex.

Each placement therefore forces the covering to overrun a filler vertex. That alone is not an
exclusion, and the exclusion comes instead from a chord bound, which we now state.
\end{remark}

\begin{lemma}[The chord at the last junction]\label{lem:chord}
Let $f=2e+1$ and let a side have $p=2$. The apex $c$-tile presents $\beta$ at $J_{2f}$ with its
$a$-ray pointing along the horizontal, and $J_{2f}$ lies at height fraction $2e/f$, so the horizontal
chord from $J_{2f}$ to the opposite side has length exactly $e(3f^2-e^2)(f-2e)/f=eN/f$. At $(3,7)$
this is $414/7$.

Moreover the tiles meeting that chord from either side meet it in \emph{whole} edges: a tile whose
interior misses the chord line lies in a closed half-plane of it, so its contact is a face
(\texttt{tile\_contact\_face}), and at a non-vertex point that face is an edge
(\texttt{contact\_is\_edge}). Disjointness of tile interiors then makes the covering a chain starting
at the end of the $c$-tile's $a$-edge, and convexity of the target keeps every endpoint within the
chord.
\end{lemma}

\begin{theorem}[$p=2$ is excluded on the tight subfamily]\label{thm:ptwodead}
Let $f=2e+1$. Then no equal side of the base-$\beta$ target at $m=1$ has $p=2$.
\end{theorem}

\begin{proof}
By Proposition~\ref{prop:tightside} the side word is $a^{2f}c$ and the last junction carries
$\{\gamma,\alpha,\beta\}$, so by Remark~\ref{rem:lastjunction} a tile must sit above the chord of
Lemma~\ref{lem:chord} with an edge beginning at chord position $a$. Enumerate those placements: the
edge on the chord is $a$, $b$ or $c$, each with two chiralities.

If the edge is $c$ it leaves the target, since $(a+c)f-eN=f^3-2ef^2+e^3=(2e+1)^2+e^3>0$.

If the edge is $a$, the height of the tile over its $a$-edge is $2\,\mathrm{Area}/a$, and the apex
stands $H(f-2e)/f$ above the chord where $H$ is the target height. Writing $A$ for the target area,
these are $2A/(Nef)$ and $2A(f-2e)/(eNf)$, equal precisely when $f-2e=1$. So on this subfamily the
third vertex lies at exactly apex height, and the target meets that height in one point: the apex.
The two chiralities then demand $N(f-2e-1)=2f^2$ and $N(f-2e+1)=4f^2$, i.e.\ $0=2f^2$ and $N=2f^2$,
the latter giving $e=f$. Both are impossible. \textup{(}At $(3,7)$ the tile would be $(21,40,41)$.\textup{)}

If the edge is $b$, one chirality overruns the chord when $-e^3+3e^2+4e+1>0$, i.e.\ for $e\le4$, and
for $e\ge5$ its third vertex falls outside on the left, since $(f^2-e^2)(2f-e)\ge e(3f^2-e^2)$
reduces to $e^3\le3e^2+4e+1$, false there; the ranges are complementary. The other chirality places
the third vertex at chord position $a+(2f^2-e^2)/2$, exceeding the chord by $(7e^2+8e+2)/(2f)>0$ for
every $e$.

No placement survives.
\end{proof}

\begin{proof}[Second proof, by the angle count at $P$]
By Proposition~\ref{prop:tightside} the last junction carries $\{\gamma,\alpha,\beta\}$ with the
$\alpha$-filler adjacent to the apex $c$-tile, so the filler lays a $b$ or a $c$ along that tile's
$a$-ray. Both exceed $a=|JP|$, so $P$ lies interior to the filler's edge and it contributes $\pi$
there. But $T_1$ presents $\gamma$ at $P$ and $T_2$'s chord trace ends at $P$, contributing at least
$\gamma$; the total is at least $2\gamma+\pi=2\pi+\alpha>2\pi$
(\texttt{SecondEdge.filler\_excluded}). No chord estimate and no height identity is used, and the
argument needs only $e^{2}+ef<f^{2}$ rather than $f=2e+1$ --- see Corollary~\ref{cor:pbound}, which
draws the same conclusion from the side count.
\end{proof}

\begin{remark}[The audit of Theorem~\ref{thm:ptwodead} against the apex rigidity]\label{rem:ptwoaudit}
Theorem~\ref{thm:ptwodead} predates \S\ref{sub:apex}, and the two overlap. Reconciling them shows the
earlier enumeration was doing more work than necessary, and we record this rather than leave two
arguments whose relation is unclear.

Remark~\ref{rem:lastjunction} has the $\alpha$-filler standing between the apex $c$-tile and the
$a$-tile, so it lays one of its flanks --- a $b$ or a $c$ --- along the apex tile's $a$-ray. Both
exceed $a=|JP|$, so $P$ would lie \emph{interior} to that edge, contributing $\pi$ at $P$. But
Theorem~\ref{thm:apexconfig} puts $T_1$ at $\gamma$ there, and $T_2$'s trace on the chord ends at $P$,
so $T_2$ contributes at least $\gamma$ as well. The total is at least $2\gamma+\pi=2\pi+\alpha>2\pi$.

So that configuration does not occur at all, and the placements enumerated in the proof of
Theorem~\ref{thm:ptwodead} are placements of a tile that cannot be there. The conclusion is unharmed
--- indeed Corollary~\ref{cor:pbound} re-derives it, on a wider range of members and without any chord
estimate --- and Theorem~\ref{thm:lastjunction}'s conclusion, that the adjacent tile lays an $a$-edge,
is likewise immediate from the same count. What the audit removes is the impression that the tight
subfamily needs the height identity: it does not.
\end{remark}


\begin{theorem}[The last-junction dichotomy]\label{thm:lastjunction}
Let $(e,f)$ be any member with $f=2e+1$, and consider the last junction $J$ of an equal side --- the
junction where the final $c$-edge begins, at distance $f^3-f^2$ from the base corner, hence at height
fraction $(f-1)/f$ for every value of $p$. Let $X$ be the tile angularly adjacent to the apex
$c$-tile $C$ on the interior side of $C$'s $a$-ray. Then $X$ lays an $a$-edge along that ray,
coinciding with $C$'s. Equivalently, \emph{$X$ does not present $\alpha$ at $J$}.
\end{theorem}

\begin{proof}
$X$ lays one of its own flanks along $C$'s $a$-ray. If it lays $c$, the run $c-a$ must be covered
inside the chord of Lemma~\ref{lem:chord} and that window is empty. If it lays $b$, the window pins
the covering to a single $a$-edge, and the height identity of Theorem~\ref{thm:ptwodead} then places
that tile's third vertex at the apex, where its remaining sides fail to be tile sides. So $X$ lays
$a$. The flanks of $\alpha$ are $b$ and $c$, so $X$ is not an $\alpha$-vertex.
\textup{(}\texttt{alpha\_cannot\_lay\_a}, \texttt{adjacent\_angle\_not\_alpha}.\textup{)}
\end{proof}

\begin{remark}[What the dichotomy settles, and what it leaves]\label{rem:dichotomy}
Theorem~\ref{thm:ptwodead} is the case $p=2$: there the tightness of Proposition~\ref{prop:tightside}
forces $X$ to be the $\alpha$-filler, which Theorem~\ref{thm:lastjunction} forbids. The dichotomy is
strictly more general --- it is stated for every $p$ and every vertex figure --- and at $p=1$ it
forces structure without closing the case:

\begin{itemize}
\item[] in the figure $\{\gamma,\alpha,\beta\}$ it forces the angular order $C(\beta),\gamma,\alpha$,
so the side word cannot end $a\,c$ with the $\alpha$ adjacent to $C$;
\item[] in the figure $\{3\alpha,2\beta\}$ it forces $X$ to be the second $\beta$. That $\beta$ shares
$C$'s $a$-edge and is the reflection of $C$ across it, so both present $\gamma$ at the far end $Q$ of
that edge. Since $2\gamma=\pi+\alpha<2\pi$, the point $Q$ is interior. Its figure is in fact
\emph{exactly} $\{\beta,3\gamma\}$: by Theorem~\ref{thm:apexconfig} the middle apex tile $T_2$ also
meets $Q$, its chord trace ending there, and a straight angle from $T_2$ would give
$2\gamma+\pi>2\pi$; so $T_2$ contributes a third $\gamma$, leaving $2\pi-3\gamma=\beta$. The
alternative $\{2\alpha,2\beta,2\gamma\}$, which carries only two $\gamma$'s, is excluded.
\end{itemize}

At $(3,7)$ the reflected tile is $(J,Q,R)$ with $R=(207,100\sqrt{187}/7)$, a genuine $(21,49,40)$
placement lying inside the target, and both figures at $Q$ are consistent. So the two surviving
configurations at $p=1$ --- the mirror partner and the $\gamma$ --- are a single phenomenon, and
neither is excluded by the chord or the height identity. Closing $p=1$, and with it
Hypothesis~\ref{hyp:walls} on the whole subfamily $f=2e+1$, requires a further mechanism ---
and \S\ref{sub:nogo} rules out the two most natural candidates, so that mechanism must be metric.
\end{remark}

\begin{theorem}[The side carrying $T_2$'s $c$-edge is not mirrored]\label{thm:nobothmirror}
Assume $e^{2}+ef<f^{2}$. The middle apex tile $T_2$ lays its $b$ on one of the two inner rays at $A$
and its $c$ on the other. Then the equal side on the ray carrying the $c$ has figure
$\{\gamma,\alpha,\beta\}$ at its last junction, not $\{3\alpha,2\beta\}$.

In particular the two equal sides are not both mirrored, and at least one carries
$\{\gamma,\alpha,\beta\}$.
\end{theorem}

\begin{proof}
Suppose both do. By Remark~\ref{rem:dichotomy} the tile $X$ adjacent to the apex $c$-tile on each side
is that tile's reflection across their common $a$-edge, and by Theorem~\ref{thm:apexconfig} that edge
lies along the chord $C$. So each $X$ is the reflection of the corresponding apex tile across $C$, and
in particular each has the reflection $R$ of the apex $A$ across $C$ as a vertex. Thus \emph{both}
mirror partners meet at the single point $R$, which lies on the axis of symmetry, at height
$H-2H/f$ --- that is, on the chord through the second junction.

Now consider $T_2$. It presents $\alpha$ at $A$, laying its $b$ on one of the two inner rays and its
$c$ on the other; let $V$ be the far endpoint of the $c$, at distance $c$ from $A$. Its horizontal
offset from the axis is $c\sin\frac{\alpha}{2}=ef/2=a/2$, exceeding the offset
$b\sin\frac{\alpha}{2}=e(f^{2}-e^{2})/(2f)$ of $P'$ by $e^{3}/(2f)>0$. Testing $V$ against the three
edges of the mirror partner on that side gives cross products
\[
\frac{e^{3}\sqrt{4f^{2}-e^{2}}}{2},\qquad
\frac{e^{3}\sqrt{4f^{2}-e^{2}}\,(f-e)(f+e)}{2f^{2}},\qquad
\frac{e\sqrt{4f^{2}-e^{2}}\,(e^{2}-ef-f^{2})(e^{2}+ef-f^{2})}{2f^{2}} .
\]
The first two are positive since $e<f$. In the third, $e^{2}-ef-f^{2}<0$ always, and
$e^{2}+ef-f^{2}<0$ is exactly the hypothesis, so the product is positive. All three agree in sign, so
$V$ lies \emph{strictly inside} that mirror partner. But $V$ is a vertex of $T_2$, so the interiors of
$T_2$ and of the mirror partner meet --- contradicting the dissection. Only the mirror partner on the
ray carrying the $c$ enters the test, so one side suffices; the other ray is handled by reflection.
\end{proof}

\begin{corollary}[The figure at $P$ on that side is fully determined]\label{cor:figureP}
On the equal side carrying $T_2$'s $c$-edge, the point $P$ lies at distance $b<c$ from $A$ along that
edge, hence interior to it, so $T_2$ is the straight tile at $P$. The figure at $P$ is then forced:
\[
\underbrace{\gamma}_{T_1}\;+\;\underbrace{\pi}_{T_2}\;+\;\underbrace{\beta}_{Y}\;+\;\underbrace{\alpha}_{U}\;=\;2\pi ,
\]
with $Y$ the tile below $T_1$'s $a$-edge and $U$ a single further tile presenting $\alpha$. No other
distribution is possible.
\end{corollary}

\begin{proof}
$T_1$ presents $\gamma$ at $P$ and $T_2$ contributes $\pi$. By Theorem~\ref{thm:nobothmirror} the
figure at the last junction of that side is $\{\gamma,\alpha,\beta\}$, in which $Y$ presents $\gamma$
at $J$ and therefore $\beta$ at $P$, those being the two angles at the ends of its $a$-edge. That
leaves $2\pi-\gamma-\pi-\beta=\pi-\gamma-\beta=\alpha$, so exactly one further tile meets $P$ and it
presents $\alpha$.
\end{proof}

\begin{proposition}[The two placements of $U$]\label{prop:Uplacements}
On the side carrying $T_2$'s $c$-edge, the tile $U$ of Corollary~\ref{cor:figureP} has $\alpha$ at $P$,
so its flanks there are $b$ and $c$, one on the ray shared with $Y$ and one on the ray shared with
$T_2$. Exactly two placements arise.
\begin{enumerate}
\item[\textup{(i)}] $U$ lays $c$ towards $Y$ and $b$ towards $V$. Then its $c$ matches $Y$'s exactly,
so the two share the vertex $W$; and its third side is an $a$-edge lying \emph{along the chord through
the second junction}, immediately to the right of the $a$-edge that the $c$-tile below $J$ lays there.
At $(3,7)$ the two endpoints are $\bigl(\tfrac{1182}{7},\tfrac{100\sqrt{187}}{7}\bigr)$ and
$\bigl(\tfrac{1329}{7},\tfrac{100\sqrt{187}}{7}\bigr)$, at distance exactly $21$.
\item[\textup{(ii)}] $U$ lays $b$ towards $Y$ and $c$ towards $V$. Then its vertex falls strictly
inside $Y$'s $c$-edge, at distance $b$ from $P$, and $U$ presents $\gamma$ there against $Y$'s straight
angle.
\end{enumerate}
In both placements the vertex $V$ of $T_2$ lies interior to $U$'s edge, since $|PV|=c-b<b<c$.
\end{proposition}

\begin{remark}[What the cascade gives]\label{rem:cascade}
Placement (i) continues the pattern of \S\ref{sub:apex} one level down: two consecutive $a$-edges lie
flush along the chord through the second junction, laid by the $c$-tile below $J$ and by $U$. Neither
placement is contradictory as far as we can determine, so $U$ does not close the case; but the
configuration on the unmirrored side is now determined down to this single binary choice, with every
vertex located exactly.
\end{remark}

\begin{proposition}[The figure at $P'$ on the mirrored side]\label{prop:figurePprime}
On the equal side carrying $T_2$'s $b$-edge, the figure at $P'$ is \emph{exactly} $\{\beta,3\gamma\}$,
whichever straight figure that side carries at its last junction.
\end{proposition}

\begin{proof}
There $T_2$ has its $b$-endpoint at $P'$, so by Corollary~\ref{cor:noTP} it presents $\gamma$; the apex
tile $T_3$ presents $\gamma$ as well. If the side carries $\{3\alpha,2\beta\}$, its adjacent tile is the
mirror partner and presents $\gamma$ at $P'$, giving $3\gamma$ and a residual $2\pi-3\gamma=\beta$. If
it carries $\{\gamma,\alpha,\beta\}$, the adjacent tile presents $\beta$ at $P'$, giving
$2\gamma+\beta$ and a residual $2\pi-2\gamma-\beta=\gamma$. Either way the multiset is
$\{\beta,3\gamma\}$.
\end{proof}

\begin{remark}[Both apex vertices are now determined]\label{rem:bothdetermined}
Corollary~\ref{cor:figureP} and Proposition~\ref{prop:figurePprime} together determine the vertex
figures at $P$ and $P'$ completely: on the side carrying $T_2$'s $c$ the figure is
$\gamma+\pi+\beta+\alpha$ with $T_2$ straight, and on the side carrying its $b$ it is
$\{\beta,3\gamma\}$. The Case~A / Case~B alternative of Remark~\ref{rem:twosides} is therefore not a
branching of the argument but a \emph{labelling} of the two equal sides, and the labelling is decided
by which ray $T_2$ puts its $c$ on.
\end{remark}


\begin{remark}[The first coupling of the two sides]\label{rem:coupling}
Every previous constraint in this branch has been local to one equal side; Theorem~\ref{thm:nobothmirror}
is the first that forbids a combination across both. Its mechanism is that the two mirror partners are
forced to share the vertex $R$ on the axis, which leaves the middle apex tile nowhere to put its far
vertex.

Note also that the sign condition governing the third cross product is $e^{2}+ef<f^{2}$ --- exactly
the hypothesis under which $a$ is the shortest side and the whole of \S\ref{sub:apex} operates. That
the same inequality controls both an edge-length comparison and a containment test is not obvious a
priori, and it suggests the condition is the natural one for this geometry rather than an artefact of
the arguments.
\end{remark}


\begin{remark}[Scope, and where $e\ge2$ is easier]\label{rem:ptwoscopeb}
Theorem~\ref{thm:ptwodead} closes $p=2$ on the whole subfamily $f=2e+1$, hence at $(1,3)$, $(2,5)$,
$(3,7)$, $(4,9)$, $(5,11)$ and onward; verified independently by exact enumeration of the six
placements for every member with $e\le400$. Combined with Corollary~\ref{cor:ptwo} it leaves $p=1$ as
the only obstruction to Hypothesis~\ref{hyp:walls} on that subfamily.

It does not extend beyond $f=2e+1$: the height identity used above holds exactly there.

Two by-products. First, $c-a$ is a gap of $\langle a,b,c\rangle$ precisely for $e\ge2$, failing at
$(1,3)$ where $c-a=2a$; this is the first point in the problem where $e\ge2$ is easier than $e=1$.
Second, at $(3,7)$ the base walk $(0,e,2e)$ carries no $a$-edge, so the corner tile lays $c$ on the
base and $a$ on the side, forcing $p\ge1$ on both equal sides; with $p=2$ excluded,
Hypothesis~\ref{hyp:walls} at $(3,7)$ therefore requires the base walk $(f,e,e)$.\end{remark}

\subsection{Two mechanisms that provably cannot close \texorpdfstring{$p=1$}{p=1}}\label{sub:nogo}

Remark~\ref{rem:dichotomy} ends by saying that closing $p=1$ ``requires a further mechanism''. That
is a hedge, and it can be replaced by theorems. This subsection records two natural mechanisms and
proves that neither can work, and one arithmetic fact that survives and points the other way. The
member is $(e,f)=(3,7)$ throughout: tile $(a,b,c)=(21,40,49)$, target $(343,343,414)$, $N=138$.

Fix an equal side with $p=1$. Its word is a permutation of $a^{7}c^{4}$ ending in $c$
($7\cdot21+4\cdot49=343$), so it has $11$ edges and $10$ interior junctions. Each edge is laid by a
tile, and the two angles of that tile at the ends of the edge are the two angles \emph{not} opposite
it: an $a$-edge is flanked by $\beta$ and $\gamma$, a $c$-edge by $\alpha$ and $\beta$. Call the
choice of which flank sits at the lower end the \emph{orientation} of the edge.

\begin{lemma}[Endpoint angles]\label{lem:endpoints}
The angle at the top of the last edge is $\alpha$, and the angle at the bottom of the first edge is
$\beta$.
\end{lemma}

\begin{proof}
By Lemma~\ref{lem:census} the apex fills uniquely as $\{3\alpha\}$ and each base corner uniquely as
$\{\beta\}$. The last edge of the side is a $c$-edge, whose flanking angles are $\alpha$ and $\beta$;
the tile at the apex presents $\alpha$, so the top angle is $\alpha$. The single tile at the base
corner presents $\beta$, so the bottom angle of the first edge is $\beta$.
\end{proof}

\begin{proposition}[The junction automaton is consistent, and hugely so]\label{prop:nogoauto}
Regard the side as a transfer system: a state is the angle at the top of the current edge, and a
transition is admissible when the pair (top angle of one edge, bottom angle of the next) extends to a
straight figure at their junction. Then the only inadmissible pair is $(\gamma,\gamma)$, which can
occur only at a junction between two $a$-edges. Consequently, over all $120$ words $a^{7}c^{4}$
ending in $c$ and all $2^{11}$ orientations, exactly $10\,560$ configurations satisfy every junction
constraint together with Lemma~\ref{lem:endpoints}.
\end{proposition}

\begin{proof}
A junction of an equal side carries a straight figure, so by Lemma~\ref{lem:census} it is
$\{\gamma,\alpha,\beta\}$ or $\{3\alpha,2\beta\}$. The figure $\{\gamma,\alpha,\beta\}$ has one angle
of each kind, so it admits a pair of side angles precisely when the two are distinct. The figure
$\{3\alpha,2\beta\}$ contains no $\gamma$, so it admits precisely the pairs drawn from
$\{\alpha,\beta\}$, all four of which leave an admissible remainder ($3\alpha$, $2\alpha+\beta$,
$2\alpha+\beta$, $\alpha+2\beta$ respectively). A pair is therefore inadmissible exactly when it is
equal and involves $\gamma$; and $\gamma$ occurs only as a flank of an $a$-edge. The count is a
finite check over $120\cdot2^{11}$ cases.
\end{proof}

\begin{proposition}[The census contributes one relation, and it is already spent]\label{prop:nogocensus}
Let $x$ be the number of junctions of the side carrying $\{\gamma,\alpha,\beta\}$. Then $x\ge7$, and
the $\alpha$- and $\beta$-counts along the side are linearly dependent: they sum to the identity
$13=13$. Both $x=7$ and $x=10$ occur in solutions, so the census does not determine $x$.
\end{proposition}

\begin{proof}
The $7$ $a$-edges contribute one $\beta$ and one $\gamma$ each, the $4$ $c$-edges one $\alpha$ and
one $\beta$ each: $7\gamma$, $11\beta$, $4\alpha$ in all. By Lemma~\ref{lem:endpoints} the apex
absorbs one $\alpha$ and the base corner one $\beta$, leaving $7\gamma$, $10\beta$, $3\alpha$
distributed over the $20$ side-tile slots at the $10$ junctions. A $\{3\alpha,2\beta\}$ junction
contains no $\gamma$ and a $\{\gamma,\alpha,\beta\}$ junction exactly one, so the $7$ $\gamma$'s
occupy $7$ distinct junctions of the latter type, whence $x\ge7$.

Write $u$ for the number of junctions whose $\gamma$ lies on a side tile and whose other side slot is
$\beta$, and $(a_i,b_i)$ with $a_i+b_i=2$ for the side slots at the $10-x$ junctions of type
$\{3\alpha,2\beta\}$. Counting $\alpha$ and $\beta$ separately,
\[
(7-u)+(x-7)+\textstyle\sum_i a_i \;=\;3,
\qquad
u+(x-7)+\textstyle\sum_i b_i \;=\;10 .
\]
Adding and using $\sum_i(a_i+b_i)=2(10-x)$ gives $7+2x-14+20-2x=13$, an identity. The two equations
are therefore dependent. Taking $x=10$, $u=7$ satisfies both, as does $x=7$, $u=4+\sum_i a_i$ for any
$\sum_i a_i\le3$.
\end{proof}

\begin{remark}[What the two no-go results mean]\label{rem:nogoscope}
Proposition~\ref{prop:nogoauto} shows that the local combinatorics of the side is almost unconstrained:
the automaton forbids one adjacency and leaves $10\,560$ configurations. Proposition~\ref{prop:nogocensus}
shows that adjoining the census adds nothing, because the two counting equations are dependent. This
is the third independent confirmation of the warning in Lemma~\ref{lem:census} that corner counting
cannot see $p$, and it is also the exact reason the retraction described in
Remark~\ref{rem:parityscope} was possible: a counting argument at $(3,7)$ has one relation to spend
and it is spent.

The conclusion is directional, not merely negative. Any mechanism that closes $p=1$ must use metric
information --- positions, not multiplicities. Proposition~\ref{prop:straddle} is the first such fact.
\end{remark}

\begin{proposition}[Every junction chord is straddled]\label{prop:straddle}
Let $J$ be a junction of an equal side at distance $d$ from the apex, and let $C$ be the chord of the
target through $J$ parallel to the base. Then some tile of the tiling meets both open sides of $C$.
\end{proposition}

\begin{proof}
Since $d$ is a sum of edge lengths $21$ and $49$, we have $d=21i+49j$ and the height fraction is
$\rho=d/343=(3i+7j)/49$. Write $n=3i+7j$, so $1\le n\le48$ for a junction strictly between the apex
and the base. The chord cuts off a triangle similar to the target with ratio $\rho$, whose area in
units of the tile area is
\[
138\rho^{2}\;=\;\frac{138\,n^{2}}{2401}.
\]
Now $2401=7^{4}$ and $138=2\cdot3\cdot23$, so $\gcd(138,2401)=1$ and the quantity is an integer if
and only if $7^{4}\mid n^{2}$, i.e.\ $49\mid n$ --- impossible for $1\le n\le48$. Hence the area above
$C$ is not an integral number of tiles. If no tile met both open sides of $C$, every tile would lie
weakly on one side and the area above $C$ would be a sum of whole tile areas, a contradiction.
\end{proof}

\begin{remark}[Why this is the right direction]\label{rem:straddlescope}
At the last junction, $n=7$ and the area above the chord is $138/49=2.816\ldots$ tiles, while the
chord has length $414/7=59.14\ldots$ and a tile meets a line in a segment of length at most $c=49$;
so at least two tiles meet that chord, and at least one straddles it. This is exactly the region
below the chord that the enumeration of Theorem~\ref{thm:ptwodead} does not reach --- see the
discussion of the $c\,c$ ending in Remark~\ref{rem:dichotomy} --- and it is the reason the surviving
configurations survive: not that the chord bound fails there, but that it was never applied there.
Whether the straddling constraint can be sharpened into a contradiction is open.
\end{remark}

\begin{proposition}[Chord decomposition at the last junction]\label{prop:chorddecomp}
Let $C$ be the chord of the target through the last junction $J$ of an equal side at $(3,7)$, so
$|C|=414/7$. Let $\mathcal U$ be the set of tiles whose interiors meet the open half-plane above $C$.
Each $T\in\mathcal U$ either \emph{straddles} $C$ (its interior meets both open sides) or is
\emph{flush from above} (it lies weakly above the line of $C$). Then:
\begin{enumerate}
\item[(a)] a flush tile meets the line of $C$ in a full tile edge, of length $21$, $40$ or $49$;
\item[(b)] the traces on $C$ of the tiles of $\mathcal U$ have pairwise disjoint interiors and cover
$C$, so their lengths sum to $414/7$;
\item[(c)] at most two tiles of $\mathcal U$ are flush, and the multiset of flush lengths is one of
$\varnothing,\{21\},\{40\},\{49\},\{21,21\}$;
\item[(d)] at least one tile straddles, and if none is flush then at least two straddle.
\end{enumerate}
\end{proposition}

\begin{proof}
(a) For a tile lying weakly above the line of $C$, that line is a supporting line, so the intersection
is a face of the tile --- a vertex or an edge (\texttt{contact\_is\_edge}). A trace of positive length
excludes the vertex case, and the edge lengths are $a=21$, $b=40$, $c=49$.

(b) The closed region above $C$ is covered by the tiles, and a point of $C$ interior to the target is
a limit of points of that region, hence lies in a tile of $\mathcal U$, which is closed; so the traces
cover $C$. If two tiles of $\mathcal U$ had traces overlapping in a segment, both would occupy area
immediately above that segment, contradicting disjointness of interiors. Additivity of length gives
the sum.

(c) The flush lengths are integers drawn from $\{21,40,49\}$ and, by (b), sum to at most
$|C|=414/7=59.142\ldots$. The only multisets meeting that bound are the five listed: already
$21+40=61>|C|$, and $21+21+21=63>|C|$.

(d) By the computation in Proposition~\ref{prop:straddle} the area above $C$ is $138/49$ tiles, not an
integer, so a straddler exists. If no tile of $\mathcal U$ is flush, the straddling traces alone sum to
$|C|=59.142\ldots$; a segment contained in a tile has length at most the tile's diameter, which for a
triangle is its longest side $c=49$. Hence at least two straddlers.
\end{proof}

\begin{remark}[What the decomposition gives, and the rung that remains]\label{rem:chordscope}
Proposition~\ref{prop:chorddecomp} is the first statement that constrains the configuration
\emph{below} the chord at $J$, the region Remark~\ref{rem:dichotomy} identifies as untouched. It
reduces the possibilities there to a finite list: a flush total in $\{0,21,40,42,49\}$ together with a
straddling remainder of $59.142\ldots$, $38.142\ldots$, $19.142\ldots$, $17.142\ldots$ or
$10.142\ldots$, distributed over at least one and at most a bounded number of straddling tiles.

What it does not give is rigidity: the cross-section of a straddling tile varies continuously with its
position, so no individual value is forced. The natural hope is that the $\gamma$-grading of
\S\ref{sub:gammagrading} restores discreteness --- a straddler's edges have directions that are
integer multiples of $\gamma$ modulo $\pi$, so one might expect its cross-section to be drawn from a
countable set fixed by its label. \emph{That hope is unfounded}, and we record why, since it is the
first thing a reader will try.

The grading constrains directions, not positions. Tile vertices lie in the $\mathbb Z$-module
generated by the admissible edge vectors, and for a horizontal chord what matters is the projection
of that module to the vertical axis. Since $\gamma=(\pi+\alpha)/2$ with $\alpha/\pi$ irrational, that
projection is generated by the values $\ell\sin(k\gamma)$ for $\ell\in\{21,40,49\}$ and $k$ in the label
window, and these are rationally independent: the subgroup they generate is dense in $\mathbb R$.
So no cross-section is forced, and the chord decomposition cannot by itself close $p=1$.

What the decomposition does establish is that the obstruction is located strictly \emph{below} the
chord. Above it there is only $138/49=2.816\ldots$ of a tile's area, so that region is enumerable up
to boundedly many tiles --- which is precisely the enumeration Theorem~\ref{thm:ptwodead} performs.
The surviving configurations of Remark~\ref{rem:dichotomy} place the decisive tile below the chord,
where no area bound confines it. Closing $p=1$ therefore requires a constraint that reaches downward,
and we do not have one.
\end{remark}

\begin{remark}[The chord programme is exhausted]\label{rem:chordexhausted}
It is natural to hope that the family of chords does collectively what no single chord does: there
are $37$ admissible heights $\rho=n/49$ with $n=3i+7j$, each carrying its own non-integral area
$138\rho^{2}$, and a tile lying below the chord at $J$ lies above the chord at some lower junction. We
record that this hope is also unfounded, so that the reader need not pursue it.

For a tile $T$ and a chord $C$ write $\varphi_T(C)\in[0,1]$ for the fraction of $T$'s area lying above
$C$; thus $T$ straddles $C$ exactly when $\varphi_T(C)\in(0,1)$. The area statement at height $\rho$ is
\[
\sum_{T}\varphi_T(C_\rho)\;=\;138\rho^{2},
\]
and this is nothing but additivity of area over the tiling: the left-hand side \emph{is} the area
above $C_\rho$, decomposed tile by tile. It therefore holds identically, for every tiling of the
target, and the $37$ admissible heights supply $37$ instantiations of one identity rather than $37$
constraints.

The only information extractable from the family is the integrality bit --- where $138\rho^{2}\notin
\mathbb Z$, some $\varphi_T(C_\rho)$ lies strictly between $0$ and $1$ --- and that is precisely
Proposition~\ref{prop:straddle}. The nesting adds nothing further: the set of heights straddled by a
given tile is an interval, and $138$ tiles cover $37$ heights with room to spare.

This is the same degeneracy as Proposition~\ref{prop:nogocensus}, in a metric rather than a
combinatorial register: a single global identity, instantiated many times, cannot separate tilings.
Taken together with Remark~\ref{rem:chordscope} it closes the chord programme. What the chord method
yields at $(3,7)$ is exactly Propositions~\ref{prop:straddle} and \ref{prop:chorddecomp}, and closing
$p=1$ requires a mechanism that is neither counting nor chord-area. We do not have one, and we prefer
to say so.
\end{remark}





\begin{remark}[Close pairs: what the unconditional filters leave]\label{rem:closepairs}
At a close pair ($f^2\le 2ef+e^2$) Lemma~\ref{lem:basetri} no longer applies and the base admits
extra columns, all with $y=e+kf$ for some $k\ge1$. Three unconditional filters apply: the
$\gamma$-trap removes $z=0$; Theorem~\ref{thm:n1} removes $x=0$; and the corner parallelogram,
which keeps $b$-edges out of the first two and last two positions, removes any column with
$x+z\le3$ in a word of at least four edges (\texttt{cornerpara\_filter}). Over the range
$e\le9$, $f\le15$ these leave $54$ of the $62$ extra columns standing, the first being
$(3,7,2)$ at $(3,4)$; every survivor has $x<f$. Closing them requires an argument bounding the
number of base $b$-edges, and no such bound follows from size: writing $y=e+kf$, the size
constraint $yb\le e(3f^2-e^2)$ reads $k\le 2ef/(f^2-e^2)$, whose vanishing is precisely the
separation condition $f^2>2ef+e^2$. The extra columns at close pairs are permitted by size by
construction, so a bound on $y$ there must be geometric.

Close pairs are not a finite exception: the condition $f\le(1+\sqrt2)e$ admits infinitely many
coprime pairs, and among the smallest are $(5,6)$ and $(4,7)$, that is $N=83$ and $N=131$. Member
by member certification therefore cannot settle this family.
\end{remark}

\begin{remark}[The extra base columns, in closed form]\label{rem:columnreduction}
The extra columns at a close pair all have $y=e+kf$ with $k\ge1$
(Remark~\ref{rem:closepairs}), and the obstruction is that no bound on $k$ follows from size.
Substituting $y=e+kf$ into the base equation collapses it: the difference is $f$ times a relation in
$x$ and $z$ alone,
\[
x(ef)+(e+kf)(f^2-e^2)+zf^2-e(3f^2-e^2)\;=\;f\bigl(xe+zf-(2ef-kb)\bigr),\qquad b=f^2-e^2,
\]
an identity (\texttt{base\_column\_reduction}). So for each $k$ the surviving columns are the
solutions of the two-variable linear Diophantine equation
\[
xe+zf=2ef-kb ,
\]
not of a three-variable search. Verified against brute enumeration on every close pair with
$e\le25$, $f\le59$: $655$ surviving columns, no exception.

Two consequences. At $k=1$ the column is pinned, $(x,y,z)=(e,\,e+f,\,2e-f)$, so it occurs exactly
when $f\le 2e-1$ (\texttt{base\_column\_k\_one}). And the equation bounds $k$. Using only $x,z\ge1$ gives $kb\le 2ef-e-f$
(\texttt{base\_column\_k\_bound}). The corner parallelogram also forces $x+z\ge4$, and since $e<f$
the minimum of $xe+zf$ subject to $x,z\ge1$, $x+z\ge4$ sits at $(x,z)=(3,1)$, so
\[
kb\;\le\;2ef-3e-f
\qquad(\texttt{k\_bound\_sharp}),
\]
better by $2e$. Checked against brute enumeration on every close pair with $e\le25$, $f\le59$: no
violation, and the bound is \emph{attained} --- equal to the largest surviving $k$ --- at $186$ of
the $198$ members carrying survivors. There is exactly one surviving column for each $k$ in range. For fixed $k$ the solutions form a
chain $x=x_0+tf$, $z=z_0-te$, so there are $\lfloor(z_0-1)/e\rfloor+1$ with $z\ge1$, and exactly one
when $z_0\le e$; writing $ke=qf+r$ this is $e(1+q)\le kf$, and since $e<f$ we have $ke/f<k$, so
$q\le k-1$ and $e(1+q)\le ek<fk$ (\texttt{one\_column\_per\_k}; neither coprimality nor the
close-pair hypothesis is used). Hence the surviving set at a member is described completely by
\[
k=1,\dots,\Bigl\lfloor\frac{2ef-3e-f}{b}\Bigr\rfloor .
\]
The description can be made exact. Reducing $xe+zf=2ef-kb$ modulo $f$ gives
$e(x-ke)=f(2e-kf-z)$, so $f\mid e(x-ke)$ and, by coprimality, $f\mid x-ke$
(\texttt{column\_x\_congruence}). For each $k$ the admissible $x$ therefore form a single residue
class modulo $f$, and the surviving columns of that $k$ are the members of that class with
$z=(2ef-kb-xe)/f\ge1$ and $x+z\ge4$. Over every close pair with $e\le25$, $f\le59$ this criterion
reproduces the surviving $k$-set at all $198$ members carrying survivors, with no exception. The $12$
members at which the bound $kb\le2ef-3e-f$ is not attained are exactly those where
$xe+zf=2ef-kb$ has no solution with $x,z\ge1$ --- a Frobenius condition on $\langle e,f\rangle$, not
a further geometric filter; at $(e,f)=(15,23)$ and $k=2$, for instance, the residue $44$ is a gap of
$\langle15,23\rangle$.

All of this is a description of the surviving set, not the required $k=0$. What it buys is that the
residual question is one linear congruence per $k$, with $k$ in an explicit finite range, in place of
an unstructured three-variable family.
\end{remark}

\begin{proposition}[The extra base columns at a close pair]\label{prop:closepaircolumns}
Let $\gcd(e,f)=1$ and $1\le e<f$, let $b=f^2-e^2$, and let $(x,y,z)$ be a base column with
$y=e+kf$, $k\ge1$, satisfying $x,z\ge1$ and $x+z\ge4$. Then
\begin{enumerate}
\item[\textup{(i)}] the base equation is equivalent to $xe+zf=2ef-kb$;
\item[\textup{(ii)}] $f\mid x-ke$, so $x$ lies in a single residue class modulo $f$;
\item[\textup{(iii)}] $kb\le 2ef-3e-f$;
\item[\textup{(iv)}] for each $k$ the pair $(x,z)$ is unique.
\end{enumerate}
Consequently the extra columns at a close pair are exactly one for each
$k=1,\dots,\lfloor(2ef-3e-f)/b\rfloor$ that admits a solution, with $x$ determined modulo $f$ and $z$
determined by $x$.
\textup{(\texttt{close\_pair\_column}, \texttt{close\_pair\_column\_unique}, \texttt{one\_column\_per\_k}.)}
\end{proposition}

The clause ``that admits a solution'' in Proposition~\ref{prop:closepaircolumns} is not decided by
(iii): the bound there ignores the residue class of $x$, and over-permits. At $(e,f)=(8,17)$ it
allows $k=1$, while $8x+17z=47$ with $x\equiv8\pmod{17}$ has no solution with $x,z\ge1$, so no
column exists; the same happens at $551$ of the $1558$ close pairs with $e\le60$. Restoring the
residue class makes the test exact.

\begin{proposition}[Sharp criterion for an extra base column]\label{prop:sharpcolumn}
Let $\gcd(e,f)=1$, $1\le e<f$, $b=f^2-e^2$. For $k\ge1$ let $x_k$ be the least positive integer with
$x_k\equiv ke\pmod f$, and write $x_k=ke-m_kf$, so that $1\le x_k\le f$. Put
\[
g(k)=(m_k+2)e-1-kf, \qquad z_k=\frac{2ef-kb-x_ke}{f}.
\]
Then:
\begin{enumerate}
\item[\textup{(i)}] an extra base column with parameter $k$ exists if and only if $g(k)\ge0$ and
$x_k+z_k\ge4$, and it is then $(x_k,\,e+kf,\,z_k)$;
\item[\textup{(ii)}] $g(k+1)\le g(k)-(f-e)$, so $g$ is strictly decreasing;
\item[\textup{(iii)}] $g(1)=2e-1-f$. Hence no member with $f\ge 2e$ carries an extra column at any
$k$, and in general the admissible $k$ form an initial segment $1,\dots,K$ with
$K\le 1+\lfloor(2e-1-f)/(f-e)\rfloor$.
\end{enumerate}
\textup{Machine-checked in \texttt{lean/Frontier.lean} as
\texttt{column\_criterion\_identity}, \texttt{column\_criterion},
\texttt{k\_one\_forces\_f\_le\_two\_e\_sub\_one}, \texttt{residue\_step} and
\texttt{g\_strict\_drop}, axiom-free beyond \texttt{propext}, \texttt{Classical.choice} and
\texttt{Quot.sound}.}
\end{proposition}

\begin{proof}
(i) By Proposition~\ref{prop:closepaircolumns}(i),(ii) a column satisfies $xe+zf=2ef-kb$ and
$x\equiv ke\pmod f$; conversely, for any $x$ in that class $f$ divides $2ef-kb-xe$, because
$xe\equiv ke^2$ and $2ef-kb\equiv ke^2 \pmod f$. So the column exists precisely when some $x\ge1$ in
the class has $z=(2ef-kb-xe)/f\ge1$, that is $xe+f+kb\le 2ef$. The left side increases with $x$, so
it suffices to test $x=x_k$, and then $x=x_k$ by Proposition~\ref{prop:closepaircolumns}(iv).
Writing $x_k=ke-m_kf$, the identity
\[
(ke-mf)e+f+k(f^2-e^2)-2ef \;=\; f\bigl(kf-(m+2)e+1\bigr)
\]
turns that inequality, after dividing by $f>0$, into $kf\le(m_k+2)e-1$, i.e.\ $g(k)\ge0$.

(ii) From $x_k=ke-m_kf$ and $x_{k+1}=(k+1)e-m_{k+1}f$ one gets
$x_{k+1}-x_k=e+(m_k-m_{k+1})f$. Both $x_k,x_{k+1}$ lie in $[1,f]$, so the left side lies in
$[1-f,f-1]$; with $0<e<f$ this forces $m_k\le m_{k+1}\le m_k+1$. Hence
$g(k+1)-g(k)=(m_{k+1}-m_k)e-f\le e-f<0$.

(iii) Since $1\le e\le f-1$, the least positive member of the class of $e$ is $e$ itself, so
$m_1=0$ and $g(1)=2e-1-f$. If $f\ge2e$ then $g(1)<0$, and $g$ is decreasing, so $g(k)<0$ for every
$k\ge1$. In general $g(k)\le g(1)-(k-1)(f-e)$, and $g(K)\ge0$ gives the stated bound on $K$.
\end{proof}

\begin{remark}[What the criterion settles, and what it leaves]\label{rem:sharpcolumnscope}
The whole $k\ge1$ question is now confined to $e<f<2e$, a strict sub-range of the close-pair range
$f\le(1+\sqrt2)e$. Of the $17296$ close pairs with $e\le200$, $5065$ have $f\ge2e$ and so carry no
extra column at all: for those the base column is forced to $k=0$ outright. At $e=1$ the condition
$f\le2e-1=1$ is incompatible with $f>e$, so no $(1,f)$ member has an extra column, which is the
arithmetic content of Theorem~\ref{thm:e1reduce} recovered independently.

The criterion is exact, not merely necessary: it was checked against brute enumeration of the base
equation on all $17296$ close pairs with $e\le200$, covering $67218$ columns, with no disagreement,
and the initial-segment shape and the bound on $K$ hold in every case. The corner-parallelogram
condition $x_k+z_k\ge4$ is not redundant --- it rejects an otherwise admissible $k$ in $119$ of the
pairs with $e\le120$ --- so it is carried explicitly in (i).

What this does \emph{not} do is exclude the surviving columns: $21837$ of them exist as solutions of
the base equation, and no arithmetic condition can remove them, since they are genuine solutions.
The arithmetic layer of the close-pair problem is now exactly described; the exclusion must come
from the interior, as Remark~\ref{rem:closepairs} already recorded.
\end{remark}

\begin{remark}[The wedge-filler hypothesis, and where it fails]\label{rem:alphamin}
Several steps of the corner chain fill a wedge of angle exactly $\alpha$ with a single
$\alpha$-tile, justified in Lemma~\ref{lem:wallclimb} by ``$\alpha$ is the minimal angle''. That is
not the load-bearing fact. What the step needs is that no two smaller angles fit, that is
$2\beta>\alpha$, which with $3\alpha+2\beta=\pi$ reads
\[
\alpha<\tfrac\pi4 .
\]
The two conditions are different, and both directions of the difference occur. At $(e,f)=(3,4)$ one
has $b<a$, so $\beta$ is the minimal angle, yet $\alpha<\pi/4$ still holds and the step survives.

The criterion is exact and integral. Since $b^2+c^2-a^2=(f^2-e^2)(2f^2-e^2)$ and
$2bc=2f^2(f^2-e^2)$,
\[
\cos\alpha=\frac{2f^2-e^2}{2f^2}=1-\frac{e^2}{2f^2},
\]
which at $e=1$ is the familiar $(2f^2-1)/(2f^2)$. So $\alpha<\pi/4$ iff $2f^2-e^2>\sqrt2\,f^2$, and
squaring (both sides positive, as $f>e$) gives
\[
e^4-4e^2f^2+2f^4>0 .
\]
The separation condition $f^2>2ef+e^2$ implies it (\texttt{separated\_wedge\_ok}): from
$f^2-e^2>2ef$ one gets $f^4+e^4>6e^2f^2$ on squaring, and the claim follows with room to spare. So
the step is safe at every separated member, and every failure is a close pair. Among $e\le11$,
$f\le19$ the failures are
\[
(4,5),\ (5,6),\ (6,7),\ (7,8),\ (7,9),\ (8,9),\ (9,10),\ (9,11),\ (10,11),\ (10,13),\ (11,12),\ (11,13),
\]
including $(5,6)$, which is $N=83$. No claim in either paper is affected: Lemma~\ref{lem:wallclimb}
is invoked only at $e=1$, where $a=f<f^2-1=b$ and $\alpha<\pi/4$ both hold, and
Proposition~\ref{prop:unify} lists four ingredients not including it. What the criterion supplies is
a precise location for the obstruction: close pairs differ from separated members geometrically, not
merely in having weaker arithmetic filters, and an argument aimed at them must either avoid the
wedge-filler step or treat $e^4-4e^2f^2+2f^4\le0$ separately.

Two arithmetic strengthenings do \emph{not} help there, which is why the bound must be interior.
Reducing the base equation modulo $e$ gives $y+z\equiv0\pmod e$ by coprimality; and the
$\gamma$-injection is exactly a bipartite matching of the $a$- and $b$-edges to the junctions, whose
sharp criterion is Hall's condition rather than the counting bound. Neither removes any of the $54$
columns that survive the three unconditional filters over $e\le9$, $f\le15$: the first is implied by
the base equation, and the second collapses to the $\gamma$-trap $z\ge1$, since with $z\ge1$ a
matching always exists.
\end{remark}

\begin{lemma}[A $c$-corner carries a side $a$-edge]\label{lem:ccornerside}
The base corner angle is $\beta$, with flanks $\{a,c\}$. If a corner tile lays $c$ on the base it
lays $a$ on the side, so that side's $a$-count $P'$ is positive; by Lemma~\ref{lem:sidequant}
$P'=fp$, hence $P'\ge f$, and with $p e+R'=f$ and the $\gamma$-trap $R'\ge1$ one gets
$1\le p\le(f-1)/e$. So the side condition of Hypothesis~\ref{hyp:walls}, namely $p=0$, fails at
any $c$-corner: under that hypothesis both base ends are $a$-edges, exactly as
Theorem~\ref{thm:e1reduce} concludes at $e=1$. What remains for $e\ge2$ is to exclude
$1\le p\le(f-1)/e$, which is the analogue of the entire $e=1$ programme rather than a single
lemma; the boundary alone does not do it, since $a$ and $c$ are both available at the first base
position once the walls column is in force.
\end{lemma}

\begin{proposition}[The chain machinery is member-independent]\label{prop:unify}
Every ingredient of the corner chain holds at all members $(e,f)$, not only at $e=1$:
\begin{enumerate}
\item[\textup{(i)}] the branch relations $3\alpha+2\beta=\pi$ and $\gamma=2\alpha+\beta$ define
the family, so the unique fills of $\beta$, $\alpha$ and $\alpha+\beta$ (\texttt{fill\_beta},
\texttt{fill\_alpha}, \texttt{fill\_alpha\_beta}) and the straight and full vertex
classifications are unchanged;
\item[\textup{(ii)}] $\cos\gamma=-e/2f<0$, so $2\gamma>\pi$ and every reflected partner dies at a
straight junction, at every member;
\item[\textup{(iii)}] $b\sin\gamma=c\sin\beta$, so an $a$-up tile spans its strip, at every member;
\item[\textup{(iv)}] for $1\le j<f$ the only decomposition of $j\,b$ is $j$ copies of $b$. Writing
$j=y+z+w$, the equation gives $e^2(z+w)+x\,ef=w f^2$, so $f\mid e^2(z+w)$ and hence
$f\mid z+w$ by coprimality; since $z+w\le j<f$ this forces $z=w=0$, so $y=j$ and $x=0$. Both
reductions are derived rather than assumed (\texttt{partition\_jb\_gen}). Verified for all
coprime $(e,f)$ with $e<12$, $f<20$.
\end{enumerate}
The single parameter that changes is the overshoot $c-b=e^2$, which is $1$ only at $e=1$. The
$e=1$ arguments that used $c-b=1$ as a unit therefore carry over with $e^2$ in its place, while
those that used its being the smallest positive integer do not.
\end{proposition}

\begin{lemma}[The $\alpha$-vertex gap, general]\label{lem:avgen}
A cover piece of a horizontal $a$-edge has its foot at $(k+1)a$ from the row origin if it is
direct, and at $(k+3)a-e^3/f$ if it is mirrored, since $2\,cu_x-3a=-e^3/f$. The mirrored foot is a
lattice point only if $f^2\mid e^2$, and at row $1$, where feet land on the base and every
junction there is at an integer abscissa, it is a junction only if $f\mid e^3$; coprimality forces
$f=1$ in both cases (\texttt{alpha\_vertex\_gap\_gen}). So the mirrored placement is excluded at
row $1$ at every member with $f\ge2$, with $e^3$ playing the role that $1$ plays at $e=1$.
\end{lemma}

\begin{remark}[The two corners are not equally rigid]\label{rem:cornerasym}
At a $c$-corner the side junction carries $\gamma$ and the remaining $\alpha+\beta$ has the unique
fill $\{\alpha,\beta\}$, so the partner is forced immediately; this is what makes
Theorem~\ref{thm:n1} short. At an $a$-corner the side junction carries $\alpha$ and the remaining
$2\alpha+2\beta$ admits both $\{2\alpha,2\beta\}$ and $\{\beta,\gamma\}$, so the side branches at
once. The base end behaves the other way: there the $a$-corner leaves $\alpha+\beta$, whose unique
fill again forces the partner and confines the next base edge to $\{a,c\}$. The chain therefore
runs along the base from either corner, and what it cannot do is pass a base $b$-edge; this is the
mechanism behind Theorem~\ref{thm:e1cascade} at $e=1$ and it is available verbatim at every member.
\end{remark}

\begin{lemma}[The corner step, unconditional at every member]\label{lem:cornerstep}
Let a base corner tile lay $a$ on the base. It presents $\gamma$ at the far end of that $a$-edge,
where the junction is straight, so the remainder is $\alpha+\beta$ with the unique fill
$\{\alpha,\beta\}$. The corner tile's $b$-edge joins a base vertex to a side vertex, so both of
its ends lie on the boundary and the tile across it lays a whole $b$; that tile presents $\alpha$
at the junction, since a straight junction carries at most one $\gamma$. Hence the $\alpha$-tile
is adjacent to the corner tile and the $\beta$-tile is adjacent to the base, and the second base
edge, lying among the $\beta$-tile's flanks, is an $a$ or a $c$.
\end{lemma}

\begin{remark}[Where the chain stops being free]\label{rem:chainobstruction}
Lemma~\ref{lem:cornerstep} iterates only as far as the same two ingredients survive. The fill of
$\alpha+\beta$ is unique at every member, so the two tiles are always an $\alpha$ and a $\beta$;
what decides whether the next base edge is confined to $\{a,c\}$ is their angular order, and that
is settled by whether the $\alpha$-tile is the partner across the previous tile's whole $b$-edge.
At the corner this is unconditional, because the $b$-chord runs from a base vertex to a side
vertex. At the $k$-th step the chord's upper end is interior, and pinning it is precisely the
crossing question of Remark~\ref{rem:straddle}. Both remaining cases, the residual configurations
at $e=1$ and the exclusion of $1\le p\le(f-1)/e$ at $e\ge2$, reduce to that single question.

That question is not answered by the hypothesis $m=1$ alone. Instrumenting the certified
exhaustions to record every partial configuration they build shows the overshoot occurring
freely: at $(1,2)$ it appears in $31$ of the $52$ recorded frontiers and at $(1,3)$ in $1128$,
the first already at depth $3$ with three tiles placed, a $c$-edge from $(2,0)$ to
$\bigl(\tfrac{11}2,\tfrac{\sqrt{15}}2\bigr)$ covering a $b$-edge from $(2,0)$ and overshooting
it by $c-b=1$. Every such configuration is eventually refuted, but not locally: the partner
pinning cannot be established by excluding the overshoot near the chord.
\end{remark}

\begin{remark}[The thick program]\label{rem:thickprogram}
The three verified facts above change the character of the $e\ge2$ case. The chain machinery
ports with strictly stronger tools: the wall climb and the first-run chirality are
member-independent (their proofs use only the angle identities and the boundary overshoot
$c-b=e^2>0$), and Lemma~\ref{lem:monochotomy} removes the $c$-chord fork entirely, so the thick
chain branches strictly less than the thin one. What remains is the finite bookkeeping: the base
letters run through a trichotomy rather than a free word, the side violations are quantized to
$p\in f\mathbb{N}$, and the reach analysis of the thin case applies verbatim where the blocks are
short. The $(2,3)$, $(2,5)$ and $(3,4)$ members are certified by full exhaustions ($19{,}677$,
$553{,}417$ and $825{,}724$ nodes; the second is the prime $N=71\equiv11\pmod{12}$, the third a
close pair). The close pair refutes deeper ($8.3f$ against $5.7f$ and $5.2f$), consistent with
its additional base decompositions.
\end{remark}

\begin{remark}[First thick-member certificates]\label{rem:thickcerts}
The full $m=1$ search of the $(2,3)$ target ($N=23$, prime, $23\equiv11\pmod{12}$) is
\textsc{exhausted} at $19{,}677$ nodes: no tiling exists, and Hypothesis~\ref{hyp:walls} holds at
$(2,3)$. The refutation profile differs from the thin members: the semigroup-run prune carries
$57\%$ of the kills (against $8\%$ at $(1,5)$), consistent with Lemma~\ref{lem:monochotomy}
strengthening the run arithmetic, and the maximal depth is $5.7f$ against $4.8f$. \end{remark}

\begin{lemma}[The crossing kill, and solitude of branches]\label{lem:solitary}
A transversal edge crossing a ladder's line within parameter $b+c$ of its start is fatal: the
covers cannot cross it, every T-junction forces the next cover edge, and
Lemma~\ref{lem:termination} forbids termination below $b+c$. A branch's own ceilings never
trigger this --- the $k$-th strip crossing is endpoint-tangent, sitting exactly $(k-1)c$ east of
the corresponding ceiling's end (\texttt{crossing\_tangency}) --- but a \emph{different} same-row
branch's ceilings do whenever it sits $1$ to $2f-1$ columns west: the offset $D f$ lands the
crossing inside a ceiling for every $-2f<D<0$, at parameter $b$ or $2b<b+c$. Hence any strip
junction of an $A$-branch with another $A$-branch $1$ to $2f-1$ columns to its west admits no
straddler, and its fills are mates: at most one $A$-branch per $2f$-column window per row-band
carries any straddler at all.
\end{lemma}

\begin{lemma}[The mirrored piece is charged]\label{lem:charge}
A cover piece of a horizontal $c$-edge is direct, with foot at the row lattice, or mirrored, with
foot displaced by $1/f$; the mirrored piece presents $\beta$ at its left junction, as does its
left neighbour, so that junction carries $n_\beta\ge2$ and its figure is $\{3\alpha,2\beta\}$
(\texttt{mirrored\_left\_junction}). By the census (Lemma~\ref{lem:census}) each such junction
forces an additional $\{\beta,3\gamma\}$ vertex: a single mirrored piece already gives $v_1\ge2$
(\texttt{escape\_charge}). The off-lattice escape available at rows $j\ge3$ is therefore paid for
globally, in climber vertices, of which the kernel-checked tilings carry between one and four.
\end{lemma}

\begin{lemma}[Climbers detect deviation]\label{lem:climberdetect}
At an interior lattice point of the brick/mate structure exactly six tiles meet, contributing
$\beta,\gamma,\alpha$ from the row above and $\alpha,\gamma,\beta$ from the row below; the counts
are $(2,2,2)$ and the figure is $\{2\alpha,2\beta,2\gamma\}$. Hence a $\{\beta,3\gamma\}$ vertex
never occurs where the six incident tiles are the brick/mate six: by Lemma~\ref{lem:census} every
tiling contains such a vertex, and by Lemma~\ref{lem:charge} each mirrored cover piece forces a
further one, so each escape requires two points at which the tiling deviates from the forced
pattern. This is a detector, not an exclusion: the kernel-checked tilings do carry climbers at
lattice points, at places where their local structure is not brick/mate.
\end{lemma}

\begin{remark}[What reach buys, and what it does not]\label{rem:reachlaw}
The four kills at reach $R$ empty the window exactly for $f\le R+2$: reach $3$ gave $f\le5$,
reach $4$ gives $f\le6$, and each further row would give one further member. Theorem
\ref{thm:strippbound} raises the reach by one and no more, because the clearance bound it uses is
the fixed constant $T/b>2$. A proof for all $f$ therefore requires killing straddlers at
\emph{every} row, not at the first three. The natural route is that a ladder descends into strips
the chain has already forced, where every line of its direction is a shared $b$-edge chain broken
at the strip junctions on both sides at once --- unstaggered --- while a ladder is staggered by
$c-b=1$; a ladder in forced territory would have to run along such a line or cross tile interiors,
and both are impossible. What is missing is the bookkeeping: that the region a ladder descends
through, roughly one half-column east per strip, lies inside the forced staircase at the moment
the fork is being decided.
\end{remark}

\begin{remark}[The residue at $f\ge6$]\label{rem:f6residue}
For $f\ge6$ the surviving words form an explicit thin set, symmetric under reversal; at $f=6$ it
is a single class: $b$ at position $5$, $c$ at position $7$ ($a\,a\,a\,a\,b\,a\,c\,a$), read from
a corner whose block is at least $3$. The blocking structure is always the same: the chain's
column-$3$ line needs a row-$3$ brick, the row-$3$ fork's $A$-branch lays its south cover on the
row-$2$ floor, and the floor is edged only as far as the staircase has reached --- which the $b$
itself interrupts. Closing the row-$3$ fork is the entire remaining content of
Conjecture~\ref{conj:advance} at $e=1$.
\end{remark}

\begin{remark}\label{rem:e1residual}
The residual class for $f\ge5$ is precisely: some side's first $a$-run behind an initial $c$-block
of length at least $2$. At $(1,4)$ the single such side word dies by the line-total kill of
Remark~\ref{rem:walls14}, whose three gap facts $b-a$, $b-e^2$, $2b-a\notin\langle a,b,c\rangle$
hold for every $f\ge3$; what does not yet generalize is the case analysis behind them when the
initial block exceeds two or the run splits. Closing this class for all $f$ is the last step of
Conjecture~\ref{conj:advance} at $e=1$.
\end{remark}

\subsection{Assembly}
Fix a hypothetical tiling. Its base walk is one of the three; each funnels as above to a
configuration killed by the $T_{\mathrm{mid}}$ corner, the pentagon lemma, or the walk arithmetic
(A1-type shielding, the $a$-prefix forcing with P6-type end constraints, the $\beta$-wedge). No
case survives; no tiling exists. \qed

\subsection{Scope, certificates, verification}
Every arithmetic component is kernel-checked (\texttt{lean/}: BAdjacency, Rigidity, W2Core,
MidTriangle, SurplusLattice, GeneralPillars, MasterLemmas, Pentagon; core toolchain,
axiom-clean). For Lemma~\ref{lem:pentagon} the file \texttt{Pentagon.lean} proves, general in
$(e,f)$ and axiom-free, both facts the wedge analysis consumes: that $(0,\min(a,b))$ is a gap of
the numerical semigroup $\langle a,b,c\rangle$, and that the stub $e^2 \bmod b$ lies in that gap
\emph{immediately} --- splitting on $e^2<b$: below, the stub is $e^2$, smaller than $a=ef$ since
$e<f$ and smaller than $b$ by hypothesis; above, $b\le e^2\le ef$ bounds the remainder by both.
The second initial stub $|a-b|$ can exceed $\min(a,b)$ (e.g.\ $5>3$ at $(1,3)$), so it escapes that
gap; it is settled instead by Lemma~\ref{lem:stubgap} below, which removes the iteration from the
argument entirely. Both stubs are therefore machine-checked, general in $(e,f)$ and axiom-free.
Certified seeded refutations: $(1,3)$, $(1,4)$ at $j=1,2$, $(2,3)$ side-break, $(2,3)$ pentagon.
Engine exhaustions independently confirm $N=11,23,26,47,59,66,71,107$. The exposition above is
the assembled argument; an extended version with expanded strip/recursion diagrams is in
preparation and does not affect the claims.

\section{Realizability: the collar induction and the first complete member}\label{sec:comprealize}

Theorem~\ref{thm:basebeta-full} is the exclusion side of the branch: no member has an instance at
$m=1$. This section supplies the construction side. For the member $(e,f)=(1,2)$ --- tile
$(2,3,4)$, unit target $\Delta_1$ with base $Y_1=11$, sides $X_1=8$, $N_1=11$ --- it settles
\emph{every} remaining instance, making $(1,2)$ the first member of any incommensurable-angles
branch whose instance spectrum is completely determined.

Write $\Delta_m$ for the instance target at scale $m$ (base $mY_1$, sides $mX_1$, count
$N=m^2N_1$), and subdivide $\Delta_m$ by the standard $m$-fold edge subdivision into $m^2$ cells,
each a translate of $\Delta_1$ (``up-cells'') or of its half-turn (``down-cells''). A cell is
\emph{not} individually tileable --- that is exactly the $m=1$ exclusion --- so all constructions
below use two larger blocks whose interfaces are full straight segments.

\begin{lemma}[Unit parallelogram]\label{lem:pgram}
The parallelogram $P_1$ with vertices $(0,0)$, $(Y_1,0)$,
$(Y_1+X_1\cos\beta,\,X_1\sin\beta)$, $(X_1\cos\beta,\,X_1\sin\beta)$ --- one up-cell and one
down-cell, glued along a full cell edge --- admits a tiling by $2N_1=22$ congruent copies of the
tile. The tiling was produced by the exhaustive engine (125 nodes) and is kernel-verified in
\textup{\texttt{lean/PgramTiling22.lean}}: exact $\Z[\sqrt{15}]$ coordinates scaled by $4$, with
congruence of all $22$ tiles, containment in the closed parallelogram ($4$ half-planes), an
explicit separating edge-line for each of the $231$ tile pairs, and the exact area identity
$\sum 2A_i = 528\sqrt{15}$ --- no axioms, no imports.
\end{lemma}

\begin{lemma}[Collar decomposition]\label{lem:collar}
For $m\ge 3$, let $\Delta_{m-2}^{\mathrm{apex}}\subset\Delta_m$ be the scale-$(m-2)$ copy anchored
at the apex. The closure of $\Delta_m\setminus\Delta_{m-2}^{\mathrm{apex}}$ --- the \emph{collar},
the bottom two cell rows, $4(m-1)$ cells --- decomposes into
\[
(m-2)\ \text{parallelogram columns} \;+\; \text{one scale-2 corner triangle},
\]
where the $j$-th column ($0\le j\le m-3$) is the sheared strip with base
$[jY_1,(j+1)Y_1]$ and top $[(j+1)Y_1,(j+2)Y_1]$ at height $2X_1\sin\beta$ --- the union of two
translates of $P_1$ stacked flush --- and the corner triangle is the scale-$2$ up-triangle on the
base segment $[(m-2)Y_1,\,mY_1]$. Every interface (column-to-column, column-to-corner,
collar-to-$\Delta_{m-2}^{\mathrm{apex}}$) is a full straight segment, so the pieces tile
independently. In tiles: $44(m-2)+44=44(m-1)$, and
$N_1 m^2 = N_1(m-2)^2 + 44(m-1)$ exactly.
\end{lemma}

\begin{proof}
The subdivision arithmetic is immediate (and is kernel-checked, shape-free, as
\texttt{collar\_cells} in \textup{\texttt{lean/Collar.lean}}); the only geometric content is
flushness. The $j$-th column's lateral sides run parallel to the left leg (direction
$(X_1\cos\beta, X_1\sin\beta)$, slope $\tan\beta$); its right side, from $((j+1)Y_1,0)$ to
$((j+2)Y_1,\,2X_1\sin\beta)$, coincides with the left side of column $j+1$, and for $j=m-3$ with
the left edge of the corner triangle. The corner triangle's right edge lies on the right leg of
$\Delta_m$; the collar's top edge, from $(Y_1,\,2X_1\sin\beta)$ to $((m-1)Y_1,\,2X_1\sin\beta)$,
is exactly the base of $\Delta_{m-2}^{\mathrm{apex}}$ ($X_1\cos\beta=Y_1/2$ by the isosceles
altitude). Each column splits into two translates of $P_1$ along the full segment at height
$X_1\sin\beta$.
\end{proof}

\begin{theorem}[Realizability, member $(1,2)$]\label{thm:realize12}
$N=11m^2$ is realizable for every $m\ge 2$. With Theorem~\ref{thm:basebeta-full} the instance
spectrum of the member $(1,2)$ is completely determined:
\[
\Delta_m \text{ is tileable} \iff m\neq 1.
\]
\end{theorem}

\begin{proof}
Bases: $m=2$ and $m=3$ are the kernel-verified certificates \textup{\texttt{Tiling44.lean}} and
\textup{\texttt{Tiling99.lean}}. Step: if $\Delta_m$ is tileable ($m\ge2$) then so is
$\Delta_{m+2}$, by Lemma~\ref{lem:collar} applied at scale $m+2$: tile
$\Delta_{m}^{\mathrm{apex}}$ by hypothesis, each of the $m$ columns by two copies of
Lemma~\ref{lem:pgram}, the corner by the $m=2$ certificate. The two-step induction from bases
$2,3$ reaching all $m\ge2$ is \texttt{two\_step} in \textup{\texttt{lean/Collar.lean}}
(kernel-checked); the exclusion at $m=1$ is Theorem~\ref{thm:basebeta-full}.
\end{proof}

\begin{remark}[Relation to Zhang's constructions]\label{rem:zhangcompare}
Zhang's construction paper~\cite{Zhang} builds tiling families for the six sporadic targets whose
tile has a $2\pi/3$ angle (with a $\pi/3$ variant), using ideal trapezoids as macro-tiles; the
families there carry lower bounds and residue classes in $m$ (e.g.\ $m\ge9$, or
$m\in5\N_0+7\N_0$), leaving the small members undecided. The branch treated here,
$3\alpha+2\beta=\pi$, is disjoint from those six cases ($\gamma=2\pi/3$ forces $\beta=0$ in this
branch), and no construction family for it appears in the literature known to us; the $44$-, $99$-
and $22$-tilings above were found by our engine. The collar plays the role of Zhang's trapezoids
--- a macro-tile calculus--- but the induction it drives has no lower bound and no residue
class: every admissible $m$ is realized. Member $(1,2)$ is thus, to our knowledge, the first
incommensurable-angles (target, tile) pair with no undecided instance.
\end{remark}

\begin{theorem}[The family is not uniform]\label{thm:notuniform}
For the member $(e,f)=(1,3)$ --- tile $(3,8,9)$, $N_1=26$ --- the scale-$2$ target admits no tiling:
$N=104$ is not realized. \textup{(Exhaustive search, $43\,541\,916$ nodes.)} Hence the conclusion of
Theorem~\ref{thm:realize12} does not extend along the family: $\Delta_2$ exists for $(1,2)$ and does
not exist for $(1,3)$.
\end{theorem}

This is the reason every attempt to build the seeds at $f=3$ fails, and it is worth listing them
because each is now explained rather than merely unsuccessful: the cevian piece $C_2$ ($68$ tiles),
the symmetric middle $M_2$ ($32$ tiles), the maximal interior of the $a\times c$ lattice ($34$ tiles
left to place) and of the $b$-glued $Q_c$ lattice ($34$ likewise) all return \textsc{no tiling}, and
a seeding by unit parallelograms is impossible by counting, since $\Delta_M$ has $M(M{+}1)/2$
up-cells against $M(M{-}1)/2$ down-cells and each parallelogram consumes one of each, leaving $M$
untileable cells. There was nothing to construct.

Two consequences. First, Remark~\ref{rem:betapi3} acquires content: $(1,2)$ is the unique member of
its family with $\beta<\pi/3$, and it is also the only one known to realize its scale-$2$ target.
Second, the collar of Theorem~\ref{thm:realize12} has no base at $m=2$ for $(1,3)$; its bases must
come from $m\ge3$. Since $\Delta_4$ would require $\Delta_1+\Delta_3$ and $\Delta_5$ would require
$\Delta_2+\Delta_3$, both excluded, the set reachable from $\Delta_3$ alone is exactly the multiples
of $3$. So if $\Delta_3$ tiles, the spectrum of $(1,3)$ is plausibly $N=26m^2$ with $3\mid m$, a
shape quite unlike $\{m\ge2\}$; and if $\Delta_3$ does not tile, the member contributes nothing. We
state this as a question rather than a conjecture, and the search is running.

\begin{remark}[The program for the other members]\label{rem:collarprogram}
Lemma~\ref{lem:collar} is stated for $(1,2)$ but its geometry is member-free: the subdivision,
the column shear, and the flushness argument use only the isosceles altitude identity, valid for
every $(e,f)$. For any member, realizability of \emph{all} $m\ge2$ therefore reduces to three
finite seeds: a tiling of $\Delta_2$, of $\Delta_3$, and of the unit parallelogram $P_1$. The
$m$-direction is closed once and for all; what remains per member is the three seed searches (or
a uniform seed construction in $(e,f)$, which we have not yet found). The seeds grow quickly ---
for $(1,3)$ they are $104$, $234$ and $52$ tiles --- so this is engine territory; the collar
lemma converts each trio of successes into a complete member. The parallelogram seed of $(1,3)$
is already in hand: the tile $(3,8,9)$ tiles its unit parallelogram into $52$ copies, found by the
engine and kernel-verified in \textup{\texttt{lean/PgramTiling52.lean}} (exact $\Z[\sqrt{35}]$,
$1326$ separating lines, zero axioms).

For the family $e=1$ the seed is not needed at all: the parallelogram can be written down.

\begin{theorem}[The unit parallelogram, $e=1$]\label{thm:pgram-e1}
Let $e=1$ and $f\ge2$, so $(a,b,c)=(f,\,f^2-1,\,f^2)$ and $N_1=3f^2-1$. Then $P_1$ admits an
explicit tiling by $2N_1$ copies of the tile.
\end{theorem}

\begin{proof}
Work in the affine basis $u=(1,0)$, $w=(\cos\beta,\sin\beta)$, in which $P_1$ is the box
$[0,Y_1]\times[0,X_1]$ with $Y_1=3f^2-1$ and $X_1=f^3=f\,c$; recall $X_1\cos\beta=Y_1/2$. The
$a\times c$ cell of this basis is the parallelogram obtained by gluing two tiles along their
$b$-edge (sides $a,c$, angle $\beta$), so every whole cell carries two tiles and a cell cut along
its own $b$-diagonal carries one. Split $P_1$ by the two parallel lines through $(fa,0)$ and
$(fa+b,0)$ with direction $X_1w-fa\,u$:
\begin{enumerate}
\item[$\Lambda$] the triangle $(0,0),(fa,0),(0,X_1)$. Its hypotenuse drops exactly $a$ per row of
height $c$, hence passes through cell corners; cells with $i+j\le f-2$ are whole and those with
$i+j=f-1$ are cut along their $b$-diagonal, giving $2\binom f2+f=f^2$ tiles.
\item[$\Sigma$] the strip between the lines: the parallelogram spanned by $b\,u$ and
$X_1w-fa\,u$. Put $\sigma=(X_1w-fa\,u)/b$. Then
\[
|\sigma|^2=\frac{b^2+4f^6-N_1^2}{4b^2}=f^2,\qquad
\cos\angle(\sigma,u)=\frac1{2f}=\cos(\alpha+\beta),
\]
so $\sigma$ has length $a$ and makes angle $\alpha+\beta$ with $u$. Consequently the
parallelogram spanned by $b\,u$ and $\sigma$ has sides $b,a$ and angles $\alpha+\beta,\gamma$:
it is the parallelogram obtained by gluing two tiles along their $c$-edge, and its long diagonal
has length $c$ because $c^2=a^2+b^2+2ab\cos(\alpha+\beta)$. Stacking $b$ of them along $\sigma$
closes exactly on $(0,X_1)$ and gives $2b$ tiles.
\item[P] the remainder, which is the $\Lambda$-staircase in the lattice translated by $b$, run to
$u=Y_1$; this is legitimate since $Y_1-b=2f^2=2f\,a$. It gives $2\bigl(f^2+\binom f2\bigr)+f=3f^2$
tiles.
\end{enumerate}
The three pieces are disjoint with union $P_1$, and $f^2+2b+3f^2=4f^2+2(f^2-1)=2N_1$.
\end{proof}

The construction has been generated and verified exactly in $\Q(\sqrt D)$ --- congruence,
containment, pairwise separation, area --- for $f=2,\dots,12$, i.e.\ $22,52,94,148,214,292,382,
484,598,724$ and $862$ tiles; only $f=2,3$ had previously been obtained by search.

\begin{remark}[The obstruction at $e\ge2$]\label{rem:e2strip}
The mechanism is specific to $e=1$. For general $(e,f)$ the strip is always the parallelogram
spanned by $eb\,u$ and $X_1w-fa\,u$, whose slant has length exactly $fb$ and whose angle
satisfies $\cos\theta=e/(2f)=\sin(\alpha/2)$, i.e.\ $\theta=\alpha+\beta$; but in the sheared
basis it is combinatorially an $s\times fb$ rectangle to be filled by $a\times b$ bricks, and the
de Bruijn--Klarner criterion requires $a\mid s$ or $a\mid fb$. The latter says $e\mid f^2$, i.e.\
$e=1$; the former contradicts the compatibility condition $s\equiv-e^3\pmod{ef}$ unless
$f\mid e^2$. So no admissible width admits a brick interior when $e\ge2$, and exhaustive search
confirms that no interior of any kind exists in the cases tested: on $(2,3)$ the widths
$s=4,10,16,22$ all return \textsc{no tiling} ($1$, $30$, $181$, $855$ nodes), and on $(3,4)$ the
width $s=21$ returns \textsc{no tiling} ($72$ nodes), while the $e=1$ control is found in $10$
nodes.
\end{remark}

\begin{theorem}[The dichotomy]\label{thm:dichotomy}
The unit parallelogram of the member $(e,f)=(2,3)$ --- sides $Y_1=46$ and $X_1=27$ with angle
$\beta$, tile $(6,5,9)$, $46$ tiles --- admits no tiling. \textup{(Exhaustive search,
$4\,334\,789$ nodes.)} Together with Theorem~\ref{thm:pgram-e1} this separates the branch: at
$e=1$ the unit parallelogram exists for every $f$, and at $e=2$ it does not exist at all.
\end{theorem}

The obstruction of Remark~\ref{rem:e2strip} is therefore not a limitation of the two-lattice
architecture: at $e\ge2$ there is nothing for any architecture to find. At $(2,3)$ both the unit
cell (Theorem~\ref{thm:basebeta-full}) and the unit parallelogram are untileable.

\begin{remark}[What this does to the induction]\label{rem:e2induction}
The collar of Theorem~\ref{thm:realize12} runs on unit columns, so it is unavailable at $e\ge2$;
only the wide columns of Proposition~\ref{prop:widecol} remain, and they advance $m$ in steps of
$f$. Consequently the clean law ``tileable iff $m\ne1$'' proved for $(1,2)$ need not persist:
the spectrum may genuinely depend on $e$. For $(2,3)$, where $f=3$, the subdivision closure and
the wide-column law together would give every $m\ge4$ and every even $m$ from a single seed
$\Delta_2$, leaving exactly $m=3$ unreached from below; so the decisive experiment for the whole
$e\ge2$ story is the realizability of $\Delta_2$ at $(2,3)$, i.e.\ $N=92$. That search is
running; we record the reduction rather than guess its outcome.
\end{remark}

\begin{proposition}[Wide columns, all $(e,f)$]\label{prop:widecol}
A parallelogram with sides $p\,Y_1$ and $q\,X_1$ is a pure $a\times c$ lattice --- an
$(pN_1/f)\times(qf)$ grid of cells, hence tileable by $2pqN_1$ tiles --- if and only if $f\mid p$. Consequently, for
every member, $\Delta_M$ is tileable whenever $\Delta_t$ is tileable and $f\mid M-t$: run the
collar of $t$ rows with columns $f$ cells wide. Together with the subdivision closure
($\Delta_m$ tileable $\Rightarrow\Delta_{mk}$ tileable) this is a residue-class realizability law
valid in every member; at $e=1$ the unit column exists and the law upgrades to all $m\ge2$.
\end{proposition}
\end{remark}

\section{What is known to be dead}\label{sec:compdead}

The following approaches to this branch are closed, and are recorded so that they are not retried.

\begin{enumerate}
\item \textbf{The whole class of linear direction invariants.} Weighting a directed edge by a
function of its direction alone, antisymmetric under $\theta\mapsto\theta+\pi$, yields a
one-parameter family of identities; every member is a specialisation of the master identity of
\cite[Rem.~4.9]{Main}. The associated ideal-membership condition is satisfied identically for every
root of unity and every $(e,f)$, and the optimal $\ell^1$ bound the family can produce is
$m(3f-2e)(f+e)/f$, which is smaller than $N$ by a factor at least $11/6$ and growing like $mf$. So
no choice of weight, kernel, moment ladder or duality can exclude a candidate. Both statements are
machine-checked in \texttt{lean/LambdaFactor.lean}.
\item \textbf{Closing $e=1$ by proving $R\ge2$ on every side.} False at the base: the walk with
$R_{\mathrm{base}}=2$ is exactly the one the corner parallelogram kills, and the survivor has
$R_{\mathrm{base}}=1$. On the base, ``$R\ge2$'' restates the remaining problem rather than reducing
it. It does hold on the two equal sides, and is proved there.
\item \textbf{Boundary tile counts.} The number of tiles meeting $\partial ABC$ grows like the
number of boundary edges, and on the genuine certificates it is $28$ of $44$ and $44$ of $99$ --- a
falling fraction. Comparing it against $N$ therefore never yields a contradiction.
\item \textbf{The far side of the mismatch ray}, under the hypothesis that the ray is a complete
wall: that side is the tile scaled by $f$, and is tileable. No obstruction can come from there.
\item \textbf{Local combinatorics of an equal side, with or without the census}, for $p=1$ at
$(3,7)$. The junction automaton forbids a single adjacency and leaves $10\,560$ configurations
(Proposition~\ref{prop:nogoauto}); adjoining the corner census adds nothing, its two counting
equations being dependent (Proposition~\ref{prop:nogocensus}). Together with
Lemma~\ref{lem:census} this is the third confirmation that corner counting cannot see $p$.
\item \textbf{The chord/area programme}, beyond Propositions~\ref{prop:straddle}
and~\ref{prop:chorddecomp}. Two continuations both fail: the $\gamma$-grading does not discretize
straddle positions, because it constrains directions and the vertical projection of the vertex module
is dense (Remark~\ref{rem:chordscope}); and the family of $37$ admissible chord heights supplies $37$
instantiations of one identity --- additivity of area --- rather than $37$ constraints
(Remark~\ref{rem:chordexhausted}). Closing $p=1$ needs a mechanism that is neither counting nor
chord-area.
\item \textbf{Boundary-word invariants (Conway--Lagarias tiling groups)}~\cite{ConwayLagarias}. These
are the natural non-abelian strengthening of the direction invariants of item~1, and are known to be
strictly stronger than colouring and weighting arguments --- which is exactly the class item~1 kills
here. They nevertheless do not apply, for two independent reasons. First, the method is formulated
for regions assembled from cells of a regular lattice, and the vertices of a tiling of our target lie
in no lattice: the edge vectors have lengths in $\{21,40,49\}$ and directions $k\gamma$ with
$\gamma=(\pi+\alpha)/2$ and $\alpha/\pi$ irrational, so the vertical projection of the
$\mathbb Z$-module they generate is dense in $\mathbb R$. Second, tilings here need not be
edge-to-edge, so the boundary word of a region is not determined by its combinatorics.
\end{enumerate}

\begin{thebibliography}{9}
\bibitem{ConwayLagarias} J.~H.~Conway and J.~C.~Lagarias, \emph{Tiling with polyominoes and
combinatorial group theory}, J.\ Combin.\ Theory Ser.\ A 53 (1990), 183--208.
\bibitem{Laczkovich1995} M.~Laczkovich, \emph{Tilings of triangles}, Discrete Math.\ 140 (1995).
\bibitem{Main} V.~Bonfioli, \emph{A signed-direction invariant for triangle tilings, and the
exclusion of primes $\equiv3\pmod4$}, 2026.
\bibitem{BeesonIII} M.~Beeson, \emph{Triangle tiling III: the case $3\alpha+2\beta=\pi$},
arXiv:1206.2229.
\bibitem{BeesonIso} M.~Beeson, \emph{Tilings of an isosceles triangle}, arXiv:1206.1974.
\bibitem{Zhang} Y.~X.~Zhang, \emph{Tiling triangles with $2\pi/3$ angles}, arXiv:2512.22696.
\end{thebibliography}

\end{document}
